\documentclass{article}

\usepackage{iclr2027_conference,times}

\usepackage[utf8]{inputenc}
\usepackage[T1]{fontenc}
\usepackage{amsfonts}
\usepackage{amsmath}
\usepackage{amssymb}
\usepackage{colortbl}
\usepackage{amsthm}
\usepackage{booktabs}
\usepackage{fancyvrb}
\usepackage{etoolbox}
\usepackage{graphicx}
\usepackage{hyperref}
\usepackage{mathtools}
\usepackage{microtype}
\usepackage{tabularx}
\usepackage{url}
\usepackage{xcolor}
\usepackage{array}
\usepackage{float}
\usepackage{algorithm}
\usepackage{algorithmic}
% Front-matter artifact links. The document keeps ``hidelinks`` so cross-
% references and citations stay black; these two links are coloured locally
% instead, so adding them does not recolour every \ref in the paper.
\definecolor{artifactlink}{rgb}{0.10,0.25,0.55}
\newcommand{\micon}[1]{%
  \IfFileExists{icons/#1}{%
    \raisebox{-0.25\height}{\includegraphics[height=1em]{icons/#1}}%
  }{}%
}
\newcommand{\artifactlink}[2]{\href{#1}{\textcolor{artifactlink}{\textbf{#2}}}}
% An unnumbered footnote: no marker in the text, no number consumed, so the
% artifact links do not renumber the paper's real footnotes.
\newcommand{\artifactfootnote}[1]{%
  \begingroup
  \renewcommand{\thefootnote}{}%
  \footnotetext{#1}%
  \endgroup
}
\hypersetup{hypertexnames=false,hidelinks,pdftitle={There’s More Than One Way: Mode Collapse in RLVR \& ModeBench},pdfauthor={Anonymous Authors}}
\setlength{\textfloatsep}{4pt plus 1pt minus 2pt}
\setlength{\intextsep}{4pt plus 1pt minus 2pt}
\setlength{\floatsep}{4pt plus 1pt minus 2pt}
\setlength{\abovecaptionskip}{0pt}
\setlength{\belowcaptionskip}{0pt}

% Typography contract: natural-prose paragraphs may not end with a line that
% occupies less than half of the available measure.  Finite \parfillskip glue
% makes TeX prefer fuller endings; the post-build PDF audit in
% \path{ops/check_paper_line_fill.py} is the fail-closed guarantee.
\newcommand{\enforcehalffinalproseline}{%
  \setlength{\parfillskip}{0pt plus .20\linewidth}}
\AtBeginDocument{\enforcehalffinalproseline}
\AtBeginEnvironment{abstract}{\enforcehalffinalproseline}
% Bibliography entries are generated paragraphs. Ask each item to use one fewer
% line where feasible; the rendered audit below remains the final authority.
\makeatletter
\apptocmd{\@lbibitem}{\looseness=-2}{}{}
\apptocmd{\@bibitem}{\looseness=-2}{}{}
\makeatother

\newcolumntype{Y}{>{\raggedright\arraybackslash}X}
\newcolumntype{P}[1]{>{\raggedright\arraybackslash}p{#1}}
\newcommand{\1}{\mathbf 1}
\newcommand{\E}{\mathbb E}
\newcommand{\Bpos}{\mathcal B_x^+}
\usepackage{wrapfig}
\usepackage{paralist}
\newcommand{\mb}{\textsc{ModeBench}}
% Success-conditional breadth. \pmd is the headline breadth statistic; \neff is
% its count-valued reading, and \distk stays available for the registered
% endpoint that is retained in the appendix.
\newcommand{\pmd}{\texttt{PCMD}}
\newcommand{\neff}{\ensuremath{N^{(2)}_{\mathrm{eff}}}}
\newcommand{\distk}{\texttt{distinct@8}}
% Support counts regenerated with the payload; see mode_diversity_coverage.tex.
\input{results/mode_diversity_coverage.tex}
\input{results/pmd_retention_macros.tex}
\input{results/replay_bank_decomposition_macros.tex}
\input{results/domain_support_macros.tex}
\input{results/replay_cost_macros.tex}
\input{results/kl_replay_recovery_macros.tex}
\input{results/mathir_support_geometry_macros.tex}
\input{results/mode_diversity_hosted_cohort_macros.tex}
\input{results/qwen_level_trend_macros.tex}
\input{results/portfolio_withdrawal_macros.tex}
\definecolor{bettergreen}{RGB}{224,242,226}
\newcommand{\bettercell}[1]{\cellcolor{bettergreen}#1}
% Same five tints as each domain's card in Fig.~\ref{fig:modebench-examples}
% (\path{ops/plot_paper_modebench_examples.py}, \texttt{DOMAIN_PANEL}), so a
% domain named in prose can be matched back to its card by colour alone.
\definecolor{domGraph}{HTML}{E8F1FA}
\definecolor{domCountdown}{HTML}{E7FBF6}
\definecolor{domPython}{HTML}{EFFBE7}
\definecolor{domMathIR}{HTML}{E7FBEE}
\definecolor{domPantry}{HTML}{E7ECFB}
\DeclareMathOperator{\clip}{clip}
\newtheorem{lemma}{Lemma}[section]
\newtheorem{theorem}[lemma]{Theorem}
\newtheorem{corollary}[lemma]{Corollary}
\newtheorem{proposition}[lemma]{Proposition}
\theoremstyle{remark}
\newtheorem{remark}[lemma]{Remark}


\theoremstyle{definition}
\newtheorem{definition}[lemma]{Definition}
\newtheorem{assumption}[lemma]{Assumption}
\theoremstyle{plain}


% Provider marks: one definition each, so every table that names a model
% picks up the logo without its builder knowing about graphics.
\newcommand{\providermark}[1]{\raisebox{-.18ex}{\includegraphics[height=1.75ex]{icons/#1}}\,}
\newcommand{\qwenmark}{\providermark{qwen.png}Qwen}
% No Falcon/TII mark ships with the paper, and HuggingFace's logo would name
% the host rather than the maker. Text until a real one is added here.
\newcommand{\falconmark}{Falcon}
\newcommand{\openaimark}{\providermark{openai.png}GPT}

% Ten main figures across eight pages of body text exceed what the default
% float parameters will place, which strands the last floats past the end of
% the scientific main. These loosen placement only; no content moves.
% Float pages (pages holding only tables or figures) default to \@fpsep
% "8pt plus 2fil", so LaTeX spreads whatever floats land there across the full
% height. That is where the large gaps between stacked appendix tables came
% from: two tables of twenty lines on a fifty-line page, pushed apart to fill
% it. Packing them to the top with a fixed separation lets a third fit and
% lets the leftover space fall at the foot, where it reads as a page ending
% rather than as a gap between two objects.
\makeatletter
\setlength{\@fptop}{0pt}
\setlength{\@fpsep}{14pt}
\setlength{\@fpbot}{0pt plus 1fil}
\makeatother

\setcounter{topnumber}{3}
\setcounter{bottomnumber}{2}
\setcounter{totalnumber}{4}
\renewcommand{\topfraction}{0.92}
\renewcommand{\bottomfraction}{0.60}
\renewcommand{\textfraction}{0.06}
\renewcommand{\floatpagefraction}{0.80}


\title{There’s More Than One Way:\\
Mode Collapse in RLVR \& ModeBench}

\author{Anonymous Authors}

\begin{document}

\maketitle
\raggedbottom

\begin{abstract}

Most questions we ask a language model have more than one right answer, and
keeping several is worth something: the user gets a choice, and somewhere to go
when the first fails. Training against a verifier spends them: a correct
response earns one point whether it is the thousandth copy of a familiar
solution or one never produced before. Seeing this means deciding when two
right answers are the same, so \mb{} takes that decision from the verifier
itself: across five executable domains every accepted output is run and
identified by what it did, so breadth becomes how often two correct responses
differ --- which, unlike counting distinct answers, cannot rise merely because
accuracy did. Measured this way, training against a verifier narrows breadth
while accuracy holds or climbs, and the frontier models we measured are already
narrow. Re:Max keeps them: the same verifier that tells modes apart indexes a
memory holding one verified example per mode, each rehearsed equally, so a
solution found once is practiced as often as one found constantly. Across three
model scales, two RL objectives and a harder construction of every task, it
improves both how often a policy succeeds and how many ways it can, without
trading one for the other.
\end{abstract}

% Artifact links as an unnumbered footnote so they sit at the foot of page 1
% rather than between the abstract and the teaser figure, where the [H]-placed
% figure crowds them out.
\artifactfootnote{%
  \micon{huggingface.png}\;%
  \artifactlink{https://huggingface.co/datasets/XXXXXXXX}%
  {ModeBench dataset and models}\qquad
  \micon{gh_logo.png}\;%
  \artifactlink{https://anonymous.4open.science/r/XXXXXXXX}%
  {Code repository}\quad
  \textit{Links anonymized for review.}}

\begin{figure}[H]
  \centering
  \includegraphics[width=\linewidth]{figures/modecollapse_story.pdf}
  \caption{\textbf{Correctness can rise while sampled alternatives narrow.}
  One graph admits six valid colorings. Bars pool four eight-sample draws from
  matched Qwen2.5-3B runs. At pass 3.25 of 8, Dr.GRPO produces 22/32 correct outputs
  from one mode; Re:Max produces 32/32 across four. Circle fills mark
  vertex colors; dashed rings mark vertices left uncolored in the prompt. This illustrates the
  distinction; Fig.~\ref{fig:maxrl-factorial} separates the effects of the
  fresh objective and replay.}
  \label{fig:story}
\end{figure}

\section{Introduction}
\label{sec:introduction}

We usually evaluate reinforcement learning with verifiable rewards (RLVR)
post-training by how often a model succeeds. But test-time search and reranking
depend on having alternatives to choose from
\citep{brown2024monkeys,snell2024testtime,cobbe2021training,wang2023selfconsistency}.
This raises a question: as a model gets better at finding a correct answer
during RLVR, \textit{what happens to other correct alternatives}?

RLVR trains against an executable check \citep{guo2025deepseekr1}, but its
binary reward does not distinguish a solution the model has produced a thousand
times from one it has never produced before. Nothing in the objective asks it
to maintain diversity. Fig.~\ref{fig:story} shows what this can look like on a
graph-coloring problem with six valid answers: under Dr.GRPO
\citep{liu2025understanding}, the policy becomes more accurate while its
correct outputs concentrate onto a single mode.

Mode collapse, or narrowing of the output space, has been observed in alignment
\citep{kirk2024understanding,cui2025entropy}, open-ended response
\citep{jiang2025hivemind}, creative generation
\citep{doshi2024generative}, brainstorming and ideation
\citep{anderson2024homogenization,meincke2025chatgpt}, public argument
\citep{kim2026argument,shahid2025homogenize}, and LM-assisted writing
\citep{abdulhai2026distort}. Quantifying this narrowing requires deciding when
two outputs are actually the same answer. Token-level diversity
\citep{li2016diversity,zhu2018texygen} counts two spellings of one solution as
variety, while a judge that groups answers by meaning
\citep{farquhar2024detecting,zheng2023judging} puts a second model between the
data and the claim.

Here, we investigate the prevalence and impact of mode collapse in
problem-solving tasks under RLVR post-training. Our contributions are
threefold:

\begingroup
\setlength{\leftmargini}{1.1em}
\begin{compactenum}
\item \mb{}: We contribute a benchmark for empirically studying mode collapse
in reasoning tasks. \mb{} provides both output correctness and maintains a
unique identity for each available modal solution to a given problem. The
benchmark spans five executable domains, each with five difficulty levels
(Sec.~\ref{sec:modebench}).

\item We run \mb{} across 4 open-weight model families at 3 scales and
\MDfrontierdeployments{} frontier models. Solution modes collapsed under both
Dr.GRPO and GRPO in every domain with measurable initial breadth, and every
frontier model we measured reached fewer than $1.5$ effective modes
(Sections~\ref{sec:results-collapse} and~\ref{sec:hosted-concentration}).

\item We propose \emph{mode replay}, which stores one verified exemplar per
mode and rehearses modes equally. Added to Dr.GRPO
\citep{liu2025understanding} it raises accuracy in all \PMDblocksReDr{}
complete model--domain blocks and pairwise correct mode diversity at every
scale we trained. Added to MaxRL \citep{tajwar2026maxrl} (which we dub as Re:Max), it
adds diversity and accuracy across \mb{} difficulty levels
(Sections~\ref{sec:method} and~\ref{sec:results-retention}--\ref{sec:results-levels}).
We also prove why the intervention has to be a memory: no objective built from
binary verifier outcomes can prefer one correct mode to another, at any group
size and no matter how it weights them (Thm.~\ref{thm:objective-inertness}).
\end{compactenum}
\endgroup

Overall, our findings shed light on reasoning mode collapse and suggest that
incorporating mode replay as part of RLVR post-training can lead to both
increased diversity of solutions and improved model performance.

\section{\mb: Measuring Successful Alternatives}
\label{sec:modebench}
\label{sec:collapse}

\begin{figure}[!b]
  \centering
  \includegraphics[width=\linewidth]{figures/modebench_examples.pdf}
  \caption{\textbf{One prompt can have different verified outputs.}
  % colour the (1) and (2) like we have inline for solutions for each domain
  (1) and (2) identify different reasoning modes across the five domains, which
  include graph coloring, a normalized computation tree, executable Python
  code, algebraic problems and constraint optimization.}
  \label{fig:modebench-examples}
\end{figure}

%% In graph, take out the 1234 2235 bit, we're not going to do this here. No
%% numbers or colors. Just a 1 and 2 AND USE A SUPER DIFFERENT PASTEL COLORS.

Producing a benchmark to study reasoning mode collapse is challenging for two
reasons. First, we must ensure that a given question has several valid discrete
solutions. Second, we must be able to identify, from a reasoning trace, which
solution was reached --- for the algebra problems of
Fig.~\ref{fig:modebench-examples}D, both generating questions that admit several
solutions and recovering which mode a correct answer used.

Ensuring these two properties hold is domain specific.
Fig.~\ref{fig:modebench-examples} illustrates a subset of reasoning modes across
the considered domains; e.g., in Countdown and MathIR, modes mean the model
reached the same answer by a different route, whereas
%in Python it means that the program it wrote behaves differently on the test inputs. Finally,
in Graph coloring and PantryPlan, modes mean a different valid answer
altogether, so we report per-domain results beside every average.
App.~\ref{app:domain-overview} and App.~\ref{app:prompts} give the per-domain
rules and the prompts.

For each domain we provide a verifier that evaluates whether a solution is
correct and which reasoning mode it belongs to, built by enumerating in advance
a comprehensive set of potential solutions to each problem so that at training
time a reasoning trace is sorted into a pre-computed mode. Precisely, for a
prompt $x$ with specification $s_x$, the validator maps an output $y$ either to
failure $\bot$ or to a key $V(y,s_x)\in\mathcal C_x$, and two accepted outputs
are the same mode if and only if their keys agree. The full set of valid keys
$\mathcal C_x$ is never revealed, neither its members nor its size. A policy,
and the memory that stores its successes, only ever sees keys for answers the
policy actually produced.

\subsection{Success and verified-mode diversity}
\label{sec:collapse-metrics}

Two properties of a policy matter here and are easy to confuse: how often it is
right, and --- given that it is right --- which of the available right answers
it gives. A benchmark reporting only the first cannot see the failure we are
after, and the obvious way to report the second mostly reports the first again.

That obvious way is to count distinct correct answers in $K$ samples. Write $P$
for the chance a sample is correct and $q_c$ for how that correct probability is
divided among the valid answers. Then

\begin{equation}
 \underbrace{\Pr[D_K>0]}_{\texttt{pass@}K}=1-(1-P)^K,
 \qquad
 \underbrace{\E[D_K]}_{\texttt{distinct@}K}
 =\sum_c[1-(1-Pq_c)^K],
 \label{eq:main-occupancy}
\end{equation}

where $D_K$ counts the distinct correct answers. Both sides climb with $P$: a
policy that is simply right more often returns more distinct answers without
having broadened at all. Across the \MDcells{} untrained measurements of
App.~\ref{app:base-levels-all-scales} that count correlates
$\MDdistinctcorr$ with \texttt{pass@8} --- close enough to be a second copy of
the accuracy axis.

We want a question whose answer cannot move just because the model got more
questions right. Of two independent correct responses to the same prompt, how
often are they different answers? That is pairwise correct-mode diversity,

\begin{equation}
 \pmd = 1-\sum_c q_c^2
 = \Pr[Z_1\ne Z_2\mid \text{both responses verified}].
 \label{eq:main-pmd}
\end{equation}

Correctness has cancelled: \eqref{eq:main-pmd} depends on $q$ alone, so a policy
that doubles its success rate while allocating those successes the same way
scores what it scored before. It is the fraction of correct pairs in a draw that
disagree, unbiased given two verified draws
(App.~\ref{app:metric-details}), and easy to read --- $0$ means every correct answer is the same one, an even spread over
$m$ modes gives $1-1/m$, and $1/(1-\pmd{})$ converts back into an effective
number of distinct answers.

A fall in \pmd{} is what we call \emph{mode collapse}: a statement about where
correct probability sits, not that any answer reached probability zero, which
finite samples cannot establish. It also has a floor --- two correct draws are
needed for it to exist, so a policy that almost never succeeds supports no
statement about how its successes are spread. Every \pmd{} we quote therefore
carries the number of prompts defining it, and a measurement with too few is
left blank rather than filled with noise.

\subsection{Harder tasks with the same notion of a mode}
\label{sec:levels-design}

A benchmark is useful only for as long as models cannot saturate it, so every
domain is built at five levels of difficulty.
Fig.~\ref{fig:base-levels-all-scales} measures the Qwen2.5-Instruct family
across those levels at every scale we could run, which holds the training recipe
fixed and lets size be the only thing that varies. Two things show up at once.
Harder levels push success down within a scale, and breadth runs mildly against
accuracy even before any training: across the \MDreportable{} untrained
measurements with enough support, \pmd{} correlates $\MDpmdcorr$ with
\texttt{pass@8}. And as the models get bigger they
get more accurate and answer in fewer ways --- \pmd{} falls with scale in
\MDscalefall{} of the \MDscalesupport{} domain--level blocks that support the
comparison, while \texttt{pass@8} rises.

That second pattern is what the rest of the paper has to explain. If breadth
simply disappeared with capability there would be little to be done about it,
but size is not what settles it: at matched parameter count, different recipes
land at visibly different \pmd{}, and a recipe swap at fixed scale moves the
measure about as far as a family's whole scale range
(App.~\ref{sec:families-appendix}). The open question is therefore whether
narrowing is something a model grows into or something a particular kind of
training does to it. Section~\ref{sec:collapse-mechanism} gives a reason to
expect the latter, and the matched comparisons in Section~\ref{sec:results} test
it directly.

\begin{figure}[!t]
  \centering
  \includegraphics[width=0.95\linewidth]{figures/mode_diversity_levels_appendix.pdf}
  % This panel carries its own whitespace under the axes, so the default gap on
  % top of that reads as a break; pulled back, but not all the way.
  \vspace{-2pt}
  \caption{\textbf{Breadth narrows as models get larger, and success does not
  buy it back.} Frozen Qwen2.5-Instruct scales: colour is parameter count; black
  is the frontier model of Section~\ref{sec:hosted-concentration}, off the ramp.
  Each mark averages one scale over the five levels;
  Fig.~\ref{fig:level-construction} reads them apart. Hollow marks cleared the
  support bar at no level, and ticks are too rarely solved to place.}
  \label{fig:base-levels-all-scales}
\end{figure}

Each level is a fresh construction, not a re-scored copy of the one before it.
The validator, the mode equivalence and the distribution of how many modes a
prompt admits stay fixed; what grows is the structure to be searched --- more
hidden graph vertices, more Countdown operands, longer factor case vectors. The
fit targets only pass@1 and pass@8 of a reference model, never breadth
(App.~\ref{app:data-levels}), and the levels behave accordingly: averaged over
all \MDladderdomainstotal{} domains and every Qwen scale, \texttt{pass@8} falls at
each of the \MDladdersteps{} level steps, from \MDladderpassfirst{} at
Level 1 to \MDladderpasslast{} at Level \MDladderrungs{}.

The frontier models make the separation of the two axes starkest. They
answer $\MDfrontierpass$ of prompts correctly, far above every open checkpoint
measured here, yet their mean \pmd{} is $\MDfrontierpmdmean{}$ --- below the
$\MDsmallbestpmd{}$ of a 1.7B model that answers $\MDsmallbestpassall$ of its
own prompts.

\subsection{Concentration in frontier models}
\label{sec:hosted-concentration}

\begin{figure}[!htbp]
  \centering
  \includegraphics[width=\linewidth]{figures/gpt56_all_levels32_sampling_budget.pdf}
  % The legend row sits inside the image, so the default caption gap reads as a
  % second break; close it as for Fig.~\ref{fig:base-levels-all-scales}.
  \vspace{-9pt}
  \caption{\textbf{More samples need not recover the available modes.}
  GPT-5.6 Sol, medium reasoning; 32 problems per domain and level, 512 responses each.
  Curves are distinct verified modes in random $k$-subsets as a share of the
  modes those problems admit, so the shaded band is what no level ever
  produced; shading is the pointwise 95\% problem bootstrap. Printed counts are
  the enumerated support those problems admit, in Level 1/2/3 order, and every
  panel shares one scale (App.~\ref{app:gpt56-all-levels-discovery}).}
  \label{fig:gpt56-sampling-budget}
\end{figure}

The same measurement applies to models we did not train, and they are already
concentrated. Asked a Level-3 Graph problem with six valid colorings, GPT-5.6
Sol is right in all 1,024 draws and almost always right the same way, offering
barely one of the six. Correctness is saturated there, so \pmd{} isolates the
collapse exactly; on MathIR it reaches zero
(App.~\ref{app:hosted-concentration}).

\textbf{Sampling harder does not recover what was lost.}
The obvious remedy is to draw more. Fig.~\ref{fig:gpt56-sampling-budget}
follows five domains and three levels out to 512 draws per problem, and the
curves flatten far below what the problems admit: a model that has settled on
two colorings does not find the other four by being asked five hundred times.
Python is the extreme case. Its problems admit several hundred distinct return
vectors on average, and five hundred draws surface between one and four of
them --- under two percent of what is there, and under one percent at
Levels 2 and 3. The curves have not saturated: fourteen of the fifteen
domain--level panels still gain between 256 and 512 draws. But each doubling of
the budget buys a fraction of one mode, against gaps of several modes
everywhere and several hundred in Python, so sampling does not come remotely
close to closing them. These are finite-budget measurements and do not
establish a limit or rule out rare unseen modes
(App.~\ref{app:gpt56-all-levels-discovery}).


The pattern is not one model's quirk, and it deepens with difficulty: no
deployment reaches one and a half effective verified modes --- one and three
fifths with MathIR set aside --- and all seven narrow
from Level 1 to Level 3 while staying near-perfect on correctness, including on
the two domains whose prompts add no strategy guidance
(Fig.~\ref{fig:frontier-level-grid}, Apps.~\ref{app:hosted-mode-diversity}
and~\ref{app:prompt-hint-ablation}). The concentration is not an artifact of
how we look, and the interventions that widen it without replacing it ---
mixing deployments, quadrupling $K$, disabling reasoning --- are reported in
Apps.~\ref{app:cross-model-overlap},~\ref{app:decoding-objection}
and~\ref{app:sampling-budget-ablation}. Raising temperature widens breadth too,
but stops paying early (Fig.~\ref{fig:gpt56-temperature-curve}), and disabling
reasoning costs correctness and observed breadth together while moving \pmd{}
either way by deployment (Table~\ref{tab:hosted-level-averages},
App.~\ref{app:hosted-reasoning-control}).

\section{Re:Max: Learn from New Samples and Remember Successes}
\label{sec:method}

\subsection{Why memory addresses a different pressure}
\label{sec:collapse-mechanism}

Collapse follows from two properties of the objective, neither about model
size: a binary verifier pays the same for every correct answer, so nothing asks
a second way of being right to survive, and a fresh-sample update acts only on
what the policy just sampled --- which is what it already prefers. They close a
loop, the mode that is ahead being drawn most often, reinforced most and drawn
more often still, whose end Fig.~\ref{fig:replay-bank-balance} draws: under the
independent-logit geometry of Theorem~\ref{thm:grpo-collapse} a single winner,
a property of that geometry rather than of binary reward alone
(App.~\ref{app:theory-entropy}).

Breaking the loop does not depend on the geometry. A rare mode is sampled
rarely, so the expected force any bounded fresh-sample score puts on it
vanishes: no reweighting of the advantage reaches a mode the policy has stopped
producing (Theorem~\ref{thm:objective-inertness}). A stored target does,
whether or not the policy rediscovers it, so a fixed bank with a constant
replay dose holds every banked mode above a probability floor
(Theorem~\ref{thm:replay-retention}). The difference from a diversity-shaped
objective is not its strength but where it can act: an objective shapes the
group it just drew, a memory acts on a mode that group no longer contains. What
the retained modes are worth shows up when the task changes --- withdraw an
option from a solved prompt and the broader policies still answer, and recover
in fewer attempts
(Apps.~\ref{app:portfolio-withdrawals},~\ref{app:pantry-adaptation}).

\subsection{Replaying distinct verified successes}
\label{sec:replay}

The analysis asks for one specific thing: a target that survives its own rarity.
Each training prompt therefore owns an initially empty, bounded bank
$\mathcal B_x$. The first time a mode appears, the bank keeps one verified
exemplar $e_c$ that realizes it; later rediscoveries of the same mode add
nothing. A deterministic schedule revisits nonempty banks throughout training.
For an exemplar of token length $|e_c|$, the replay loss is
\begin{equation}
 L_{\mathrm{replay}}(x)=-\frac{1}{|\mathcal B_x|}
 \sum_{c\in\mathcal B_x}\frac{\log\pi_\theta(e_c\mid x)}{|e_c|}.
 \label{eq:lm-verified-replay}
\end{equation}
The uniform weight is the whole point: a mode found once is rehearsed exactly as
much as a mode found a hundred times, so popularity cannot compound. It is also
what makes the bank a counterweight to the fresh update rather than
another copy of it --- App.~\ref{app:replay-metric-alignment} decomposes this
loss into a term that rewards more banked probability and a term minimized by an
even spread across the banked modes, which are the two things the fresh flow
works against.


Adding this loss to the fresh update from MaxRL gives Re:Max; adding it to
Dr.GRPO gives Re:Dr. Exemplars stay outside the fresh reward group, so the two can be
varied independently, and controls run the same bank and scoring code with the
replay gradient set to zero. Banks hold only training-prompt responses, so
held-out evaluation measures what training changed in the policy rather than
what the bank can recite (App.~\ref{app:algorithm}).

\subsection{Learning from fresh rollouts}
\label{sec:maxrl-method}

% wrapfig reserves whole lines and rounds up, so the figure leaves a gap
% it did not need; reclaim the rounding, not the surrounding text.
\vspace{-0.9\baselineskip}
\begin{wrapfigure}{r}{0.56\linewidth}
  \centering
  \includegraphics[width=\linewidth]{figures/replay_bank_balance.pdf}
  \caption{\textbf{What the bank preserves.} Verified modes of one prompt,
  each area the probability the policy produces that mode. RLVR-only training
  concentrates it on the mode found first; uniform replay keeps every banked
  mode and leaves room to find more.}
  \label{fig:replay-bank-balance}
\end{wrapfigure}
Replay is added to a fresh objective, not substituted for one, so we vary that
objective to check that memory is not merely covering for a weak choice. MaxRL
combines success objectives across sampling budgets \citep{tajwar2026maxrl}; its
binary-reward group estimator weights a correct response more strongly when
successes are scarce. Running replay under both MaxRL and Dr.GRPO
\citep{liu2025understanding} therefore asks whether memory still contributes
once fresh learning is already doing better.

Neither objective avoids the starvation above. A rare mode has to be sampled
before it can contribute a score term, and a group whose responses are all
correct or all incorrect carries zero centered task advantage, so even a policy
that is succeeding often can stop receiving any comparative signal.
App.~\ref{app:binary-maxrl} gives the exact advantages and the finite-group
interpretation, and App.~\ref{app:degenerate-groups} measures how often this
actually happens during our runs.

Fig.~\ref{fig:verified-support-story} shows the two learning paths we combine:
binary rewards learn from fresh samples, while replay returns to successes the
policy has already found. Both run on the policy's own outputs, and neither
needs gold solutions.

\begin{figure}[!htbp]
  \centering
  \includegraphics[width=\linewidth]{figures/verified_support_story.pdf}
  \caption{\textbf{Fresh learning and memory use the same verified samples.}
  MaxRL learns from their binary rewards. Replay stores one exemplar per mode
  and revisits retained modes with equal weight. In the schematic, fresh
  samples add C while replay revisits A and B. A stored exemplar can continue
  to receive a training signal after it becomes rare in fresh samples.}
  \label{fig:verified-support-story}
\end{figure}

\section{Experimental Design}
\label{sec:experiments}
\label{sec:shared-protocol}

Every comparison below is paired: the two arms share initialization, prompt
stream, seed, budget, and evaluation requests, and differ only in the
intervention named. Cost is matched where it dominates: the same 3,072
updates and $G=16$ fresh rollouts, with the control at a zero replay derivative,
so replay buys no extra sampling; what it adds is one forward and one backward
pass over at most $M=16$ short banked exemplars per visited prompt
(Table~\ref{tab:run-contract}).
\phantomsection\label{sec:retention-design}%
\phantomsection\label{sec:factorial-design}%
We train Qwen2.5-0.5B, Falcon3-1B, and Qwen2.5-3B on all five Level-1 domains
for eight passes with five seeds each, adding replay to Dr.GRPO and crossing
Dr.GRPO/MaxRL with replay off and on, then repeat that factorial on Levels 2
and 3 at 0.5B. Frequency-weighted replay, UCPO, RLEP, and fixed Semantic-MaxEnt are
controls for alternative explanations. Evaluation averages four eight-sample
draws on 128 held-out prompts that never enter training or the bank
(Table~\ref{tab:run-contract} gives every shared setting).

We separately evaluate five hosted models --- GPT-5.6 Sol, GPT-5.4, Grok 4.3,
Kimi K3, and Claude Opus 4.8 --- on 480 fixed prompts spanning all five domains
and Levels 1--3, eight stateless responses each, graded through a fixed
formatting normalizer so that presentation differences are not counted as
verified alternatives (App.~\ref{app:hosted-comparison}). These are
inference-only: no replay is applied and no training effect is identified.

\section{Results}
\label{sec:results}

RLVR narrows a policy onto a single verified
answer, and it does so across scales, domains and both objectives
(Section~\ref{sec:results-collapse}). Adding a memory of distinct successes
reverses that narrowing without costing correctness, and the active ingredient
is the equal weighting rather than the reuse
(Section~\ref{sec:results-retention}). The gain is not an artifact of a weak
baseline: it survives a stronger fresh objective
(Section~\ref{sec:results-maxrl}) and a harder construction of every task
(Section~\ref{sec:results-levels}). Both arms train without a reference-KL
penalty (Table~\ref{tab:run-contract}), which leaves the control freer to
concentrate than a regularized one would be; App.~\ref{app:theory-kl} works
out what switching it on does and does not change.

\subsection{RLVR concentrates verified outputs}
\label{sec:results-collapse}

% wrapfig reserves whole lines and rounds up, so the figure leaves a gap
% it did not need; reclaim the rounding, not the surrounding text.
\vspace{-0.9\baselineskip}
\begin{wrapfigure}{l}{0.500\textwidth}
  \centering
  \includegraphics[width=0.485\textwidth]{figures/concentration_story_resampled.pdf}
  \caption{\textbf{Training concentrates verified outputs.} Final minus
  initial \pmd{} at Qwen2.5-3B, Level 1, under three objectives, five domains,
  from the independently seeded resampling: both sides of the difference are measured
  on the same surface. Marks are means over training seeds with nominal 95\%
  intervals. Fourteen of the fifteen comparisons move left; the exception is within
  half a point of zero. Support pools a comparison's seeds, from 134
  eligible prompts at the thinnest to 546; below thirty pooled, hollow. Other scales and every level:
  Fig.~\ref{fig:concentration-all-scales};
  counts: App.~\ref{app:conditional-concentration}.}
  \label{fig:concentration-story}
  % wrapfig rounds its reserved lines up; reclaim the rounding so the
  % following text does not sit in a gap.
  \vspace{-0.5\baselineskip}
\end{wrapfigure}


\textbf{Training can replace varied attempts with one repeated answer.}
The starkest case is Python. Untrained, Qwen2.5-3B returns two verified
answers on 73 of the 128 held-out prompts and a second mode on half of those.
Training takes most of that away: Dr.GRPO seeds end with as few as 22 such
prompts, and repeat themselves more on what is left, \pmd{} falling from $.168$
to $.095$. These are the resampling's counts, over its 1,024 responses per
prompt rather than the eight of
Table~\ref{tab:core-terminal-endpoints}.
The loss is coverage first, breadth second, and what it settles on is not always
correct.

This is also where the fresh objective runs out of signal. A binary-reward group
teaches nothing unless it contains both a success and a failure, and by then
almost none of them do: under the control, one update in seventy carries any
task gradient at all. Replay keeps roughly one in five informative, because a
banked exemplar does not have to be rediscovered to be trained on
(App.~\ref{app:degenerate-groups}).

\textbf{Concentration among correct outputs is separately measurable.}
Collapsing onto one answer is the extreme; the subtler and more general effect
is a policy that still succeeds in several ways but increasingly prefers one of
them. Across all five domains no comparison moves the other way: fourteen of the
fifteen move toward repetition, and the exception sits within half a point of
zero. The size of the move is a property of the domain rather than of the
objective: PantryPlan loses about fifty points, Graph and Countdown about
fourteen, Python seven to sixteen, MathIR about five. The one flat case is
MaxRL on Graph, which holds its breadth where both other objectives lose
theirs.
Replay stays above its matched control in 111 of 120 pairs, turning loss into
gain in 25
(Fig.~\ref{fig:concentration-story}; Apps.~\ref{app:conditional-concentration}
and~\ref{app:pre-rl-regression}).


\subsection{Replay improves both outcomes}
\label{sec:results-retention}

\textbf{Verified memory improves both outcomes across scales.}
Re:Dr raises correctness in every tested model--domain comparison
(Fig.~\ref{fig:maxrl-factorial}; per domain and scale in
Table~\ref{tab:cross-scale-terminal-effects}), and the breadth claim survives
on the success-conditional axis: replay ends above its control in \PMDabove{}
of the \PMDblocks{} comparisons where both sides define \pmd{}. The gaps are
not marginal --- controls finishing at \pmd{} of zero, every success the same
success, against replay arms with a third of their verified pairs differing.
Python is where the two axes separate: correctness roughly triples at
Qwen2.5-0.5B while \pmd{} moves $+.082$, and at Falcon3-1B the control never
returns two verified solutions to one prompt.
Succeeding more often is not succeeding in more ways
(App.~\ref{app:retention-numerical-detail}).

\textbf{It is the balancing, not merely the reuse.}
Replay could help simply by training on known-good outputs again. To separate
that from the equal weighting, we change one thing: instead of giving every
retained mode the same share, we weight modes by how often fresh rollouts
rediscovered them. Reuse is identical; only the balance differs. Equal
weighting wins on breadth --- clearly in Graph, Countdown and PantryPlan, with
no aggregate difference in correctness
(App.~\ref{app:frequency-replay-progress},
Fig.~\ref{fig:replay-key-weighting}). Refusing to let popularity decide how
much a mode is rehearsed is doing real work.

\subsection{Replay complements a stronger fresh objective}
\label{sec:results-maxrl}

\begin{table}[!tbp]
  \centering
  \caption{\textbf{The alternatives at Qwen2.5-0.5B, Level 1.} Each method minus
  matched terminal Dr.GRPO; Avg weights the five domains equally within seed.
  Bold marks a column's best, over the trained arms only. A superscript marks a
  paired seed count below five, an em dash too few verified pairs on one side
  to difference. The alternatives are ordered by their \pmd{} average. Colour runs red to green down each column, scaled to that
  column's largest magnitude. Per-seed values and paired intervals stay in
  the record.}
  \label{tab:direct-comparator-matrix}
  \scriptsize
  % Thirteen columns with the domains spelled out: the two widest headers,
  % Countdown and PantryPlan, set the measure rather than the digits do, so the
  % padding comes down to keep the names at full size inside \linewidth.
  \setlength{\tabcolsep}{1.4pt}
  \renewcommand{\arraystretch}{0.92}
  \begin{tabular}{@{}l*{12}{r}@{}}
    \toprule
    & \multicolumn{6}{c}{\textbf{$\Delta$ \texttt{pass@8}}}
      & \multicolumn{6}{c}{\textbf{$\Delta$ \pmd{}}} \\
    \cmidrule(lr){2-7}\cmidrule(l){8-13}
    & Graph & Countdown & Python & MathIR & PantryPlan & Avg
      & Graph & Countdown & Python & MathIR & PantryPlan & Avg \\
    \midrule
    \input{results/retention_matrix_panel_b_table_body.tex}
  \end{tabular}
\end{table}

\textbf{Memory adds value when fresh learning already uses MaxRL.}
The paired comparisons in Fig.~\ref{fig:maxrl-factorial} separate replay
from the choice of fresh objective. At all three scales, replay improves
domain-averaged correctness and \pmd{} under both Dr.GRPO and MaxRL. On
breadth every unmodified arm ends below its untrained model and every replay
arm above its control. Replay's benefit therefore persists beyond MaxRL's
sampling-aware fresh objective; the two address complementary parts of
training.

\begin{wrapfigure}{r}{0.72\textwidth}
  \centering
  % wrapfig reserves whole lines and rounds up; reclaim that rounding so the
  % plate sits against the text it belongs to.
  \vspace{-0.9\baselineskip}
  \includegraphics[width=0.70\textwidth]{figures/e118_all_scale_factorial_progress.pdf}
  \caption{\textbf{Replay complements both fresh objectives at multiple
  scales.} Diamonds are initial models; open/filled markers are training
  without/with replay. (A) averages all five domains within seed; faint paths
  show seeds. (B) draws both paths from the untrained model --- one to the
  fresh objective alone, one to it with memory --- over the domains where every
  arm clears the support bar; a path too short to carry a head is drawn without
  one. Solid paths are complete five-seed blocks, dashed paths partial with
  their counts in the record. App.
  Fig.~\ref{fig:maxrl-factorial-scale-extensions} gives each domain at every
  scale.}
  \label{fig:maxrl-factorial}
\end{wrapfigure}

Replay helps most where the policy has least to lose: tens of points of
correctness at the smallest scale, and double digits at every scale. Domain
differences survive that average and we report them --- Falcon Graph moves the
wrong way on raw support, and two of the 3B correctness intervals include zero.
Adding the replay update is what carries the gain
(Fig.~\ref{fig:replay-factorial-effects},
Apps.~\ref{app:factorial-numerical-detail}
and~\ref{app:factorial-training-curves}).

\textbf{Existing objectives that act on the same pressure.} Six controls ---
UCPO, RLEP, GAPO, SetPO, a fixed semantic-entropy bonus, and replay weighted by
rediscovery frequency --- give local rather than general gains
(Table~\ref{tab:direct-comparator-matrix}). GAPO comes closest. It penalizes a
rollout that repeats a mode the group has already produced, and how hard it
penalizes depends on how many valid answers the prompt admits. When that count
is near the group size the pressure is sharp, and GAPO doubles our breadth gain
on PantryPlan. PythonFactors prompts can admit thousands, and there it does
nothing (App.~\ref{app:diversity-comparators}). Frequency-weighted replay isolates the
same point from the other side: same bank, same reuse, popularity deciding
rehearsal. Equal weighting beats it by $+.318$ modes per prompt
($[+.272,+.358]$) at no cost in correctness
(Apps.~\ref{app:ucpo-progress} and~\ref{app:frequency-replay-progress}).

\subsection{Replay transfers to a harder construction}
\label{sec:results-levels}

\begin{wrapfigure}{r}{0.36\textwidth}
  \centering
  \includegraphics[width=0.34\textwidth]{figures/modebench_level_admission.pdf}
  \caption{\textbf{What replay adds, level by level, at Qwen2.5-0.5B.}
  Each bar is a replay arm minus its own fresh objective at the same level.
  \texttt{pass@8} averages the four domains complete at every level, \pmd{} the
  three that also clear the support bar
  (App.~\ref{app:level2-interim}).}
  \label{fig:level2-admission}
  % Reclaim wrapfig's line rounding, as for Fig.~\ref{fig:concentration-story}.
  \vspace{-0.5\baselineskip}
\end{wrapfigure}

\textbf{The benefit extends to harder task constructions.}
At Qwen2.5-0.5B, replay improves average correctness and verified-mode
diversity under both fresh objectives within each of the three matched levels,
every comparison holding the level's prompts and fresh objective fixed
(Fig.~\ref{fig:level2-admission}).


\textbf{The retained alternatives are worth having.} Withdraw one option from a
task already answered eight ways --- an ingredient, a colour, an operation, a
divisor, a menu action --- and replay portfolios more often still hold a valid
answer, in all \PWDomainsWord{} domains, and in \PWBroadDomainsWord{} of them
with the verified-draw count held fixed, and they recover in fewer attempts than
ordinary or tuned-temperature sampling in every domain
(Apps.~\ref{app:portfolio-withdrawals},~\ref{app:pantry-adaptation}).

\section{Related Work}
\label{sec:related}

\textbf{Alignment narrows what a model produces.} RLHF and reasoning RL reduce
output diversity, entropy and high-budget coverage
\citep{kirk2024understanding,dang2025diversity,yue2025rlvrlimit,jiang2025hivemind},
and \citet{cui2025entropy} trace the entropy fall to a covariance mechanism.
The collapse is theirs; ours is what survives it and what an intervention
reaches.

\textbf{What counts as a different answer.} Semantic clustering infers sameness
from language \citep{farquhar2024detecting} but needs a judge; \mb{} takes
identity from the executable check that assigns reward: two outputs differ
exactly when they execute differently. \pmd{} reads breadth given success;
entropy over all outputs and coverage counts climb with correctness
\eqref{eq:main-occupancy}. Alternatives matter because sampling-based selection
consumes them
\citep{brown2024monkeys,snell2024testtime,cobbe2021training,lightman2023lets},
and because breadth predicts survival when requirements change
(Sec.~\ref{sec:collapse-mechanism},
Apps.~\ref{app:portfolio-withdrawals},~\ref{app:pantry-adaptation}).

\textbf{Shaping the fresh group, or retaining what was found.} Outcome,
set-level and group-aware objectives reshape the fresh update
\citep{song2025outcome,chen2025passktraining,li2025divergence,li2026setpo,anschel2025groupaware,chen2026diverse};
UCPO penalizes non-uniform correct mass \citep{lochab2026ucpo}, DMPO matches a
reward-proportional target \citep{li2026distribution}, MaxRL combines sampling
budgets \citep{tajwar2026maxrl}. Two different limits separate these from
memory. An objective reading only binary outcomes cannot prefer one correct
mode to another at all (Thm.~\ref{thm:objective-inertness}), which covers MaxRL
and every pass@$k$ reweighting. One reading the group's modes escapes that, and
we share UCPO's target --- but its force decays with the mode's own
probability, so it reaches nothing the group stopped sampling
(Lem.~\ref{lem:sampled-score-bound}, App.~\ref{app:related-extended}). A bank
acts on a mode the group no longer contains
(Thm.~\ref{thm:replay-retention}); rehearsing by discovery frequency, as replay
\citep{lin1992selfimproving,rolnick2019replay} and RLEP \citep{zhang2025rlep}
do, costs verified modes against rehearsing equally
(App.~\ref{app:frequency-replay-progress}), which connects our memory to
quality-diversity archives \citep{mouret2015illuminating,ecoffet2021return}.

\section{Conclusion and Limitations}
\label{sec:conclusion}
\label{sec:discussion}
\label{sec:limitations}

\mb{} makes reliable answering and verified breadth separately observable, and
Re:Max keeps discovered successes in play after fresh sampling stops
producing them. Four limits remain: replay retains only what sampling finds;
a changed task is tested as one or two withdrawn options, not an arbitrary
revision;
\pmd{} says least where breadth most needs
establishing; and the equivalence is executable by construction, which in
MathIR certifies derivations of one answer rather than alternative answers,
so replay finds nothing there to rebalance,
leaving untested whether these
results carry where sameness cannot be executed. Measure and intervention both
point toward more spread; whether a \emph{chosen} distribution over the
admitted modes can be trained for is the question after that.

\label{sec:main-end}

% ICLR wants the AI Use, Ethics and Reproducibility statements after the
% Conclusion and before the references, where they do not count against the
% page limit. They used to close the appendix, which put them after the
% bibliography and failed the layout contract.
\clearpage
\subsubsection*{AI Use Statement}
\label{sec:ai-use}

Generative AI assisted manuscript preparation, mathematical analysis, and
supporting code. The authors reviewed and verified the resulting content and
take responsibility for the final manuscript.

\subsubsection*{Ethics Statement}

The benchmark and experiment-specific data are synthetic; this study collects
no human-subject or personal data. More verified modes are not inherently
better: a distinct execution may be risky or costly even when it passes the
validator, and \texttt{distinct@}$K$ does not measure downstream value. Systems
that evaluate or optimize breadth should report correctness, invalidity, and
application-specific costs alongside mode counts.

\subsubsection*{Reproducibility Statement}

Every reported number is computed from a versioned result file.
App.~\ref{app:reproducibility} records protocols, hashes, validations,
and regeneration commands, and links the public dataset and model archive;
App.~\ref{app:data} records split identities.
The AI Use Statement above discloses generative-AI assistance.

\clearpage
\phantomsection\label{sec:references}
\bibliographystyle{iclr2027_conference}
\bibliography{example_paper}

\long\def\crossscaletablematerial{%

% Floats: forced [H] here could not fit the page its section starts on, so
% it began the next one and left the first part empty -- and it blocked the
% cross-scale table behind it.
\begin{table}[!tbp]
  \centering
  \caption{\textbf{Where each arm ends, after eight training passes.} The levels
  behind the differences in Table~\ref{tab:cross-scale-terminal-effects}: both
  arms on correctness and on breadth over the successes alone, as paired-seed
  means at pass 8. Blocks are five seeds unless a superscript on the domain
  says otherwise (App.~\ref{app:source-integrity}); a superscript on a single
  \pmd{} value is that cell's own seed count, which can fall below its block's
  when the support bar thins one arm. An em dash is a cell where no seed defines
  \pmd{} on 30 prompts, so breadth is unmeasured rather than zero; every one is
  a Dr.GRPO control. Table~\ref{tab:cross-scale-terminal-effects} still
  differences three of the four, because a paired difference carries the lower
  bar of 20 defined prompts rather than this table's 30
  (App.~\ref{app:registered-endpoints}); the one cell where even that bar fails
  is the uncoloured entry there. Shading runs red through yellow to green from
  zero up to the largest level of that metric, on one scale shared by both arms,
  so the two columns of a pair can be read against each other.}
  \label{tab:core-terminal-endpoints}
  \scriptsize
  % No \resizebox here. It scales to \linewidth and keeps the aspect, so
  % dropping a column made the table wider per column and therefore taller.
  % Six columns set at their natural width, which is narrower and shorter.
  \setlength{\tabcolsep}{5pt}
  \begin{tabular}{@{}ll*{4}{r}@{}}
    \toprule
    \textbf{Model} & \textbf{Domain} &
      \multicolumn{2}{c}{\textbf{Dr.GRPO}} &
      \multicolumn{2}{c}{\textbf{Re:Dr}} \\
    \cmidrule(lr){3-4}\cmidrule(l){5-6}
      & & \textbf{\texttt{pass@8}} & \textbf{\pmd{}}
      & \textbf{\texttt{pass@8}} & \textbf{\pmd{}} \\
    \midrule
    \input{results/core_terminal_endpoints_table_body.tex}
  \end{tabular}
\end{table}

}

\appendix
\onecolumn
\raggedbottom

\section*{Appendix Organization}
\label{app:organization}
The supplement is organized in seven parts, mapped here to their pages.

\vspace{0.6em}

\begingroup
\providecommand{\apxpart}[1]{\par\vspace{2pt}\noindent\begingroup\parfillskip=0pt plus 1fil\relax\textbf{#1}\par\endgroup}
\providecommand{\apxdots}{\nobreak\leaders\hbox{\hspace{0.26em}.\hspace{0.26em}}\hfill}
\providecommand{\apxline}[3]{\par\noindent\begingroup
  \rightskip=2.1em plus 1fil\relax \parfillskip=-\rightskip
  \hangindent=#1\hangafter=1 \leavevmode
  #2\apxdots\nobreak\makebox[2.1em][r]{#3}\par\endgroup}
\providecommand{\apxsec}[3]{\apxline{2.4em}{\makebox[2.2em][l]{#1}#2}{#3}}
\providecommand{\apxsub}[3]{\apxline{4.0em}{\hspace*{1.4em}\makebox[2.6em][l]{#1}#2}{#3}}
\scriptsize\linespread{0.97}\selectfont\setlength{\parindent}{0pt}\setlength{\parskip}{0pt}
\apxpart{Part I. Scope and measurement}
\apxsec{\ref{app:metric-details}}{Metric and Objective Details}{\pageref{app:metric-details}}
\apxsub{\ref{app:metric-sampling}}{Measuring success and diversity}{\pageref{app:metric-sampling}}
\apxsub{\ref{app:registered-endpoints}}{Endpoints and the reported breadth axis}{\pageref{app:registered-endpoints}}
\apxsub{\ref{app:replay-metric-alignment}}{Why balance within the replay bank}{\pageref{app:replay-metric-alignment}}
\apxpart{Part II. Benchmark and method}
\apxsec{\ref{app:data}}{Benchmark and Evaluation Protocol}{\pageref{app:data}}
\apxsub{\ref{app:data-materialization}}{Split identities and executable validation}{\pageref{app:data-materialization}}
\apxsub{\ref{app:python}}{Python execution contract}{\pageref{app:python}}
\apxsub{\ref{app:data-levels}}{Matched level constructions}{\pageref{app:data-levels}}
\apxsub{\ref{app:base-levels-all-scales}}{Frozen base models across all five levels}{\pageref{app:base-levels-all-scales}}
\apxsub{\ref{sec:families-appendix}}{Other model families at matched scale}{\pageref{sec:families-appendix}}
\apxsub{\ref{app:repro-run-contract}}{Training and evaluation settings}{\pageref{app:repro-run-contract}}
\apxsec{\ref{app:prompts}}{Exact Domain Prompts}{\pageref{app:prompts}}
\apxsub{\ref{app:prompt-graph}}{Graph coloring}{\pageref{app:prompt-graph}}
\apxsub{\ref{app:prompt-countdown}}{Countdown}{\pageref{app:prompt-countdown}}
\apxsub{\ref{app:prompt-python}}{Python factors}{\pageref{app:prompt-python}}
\apxsub{\ref{app:prompt-mathir}}{MathIR action menu}{\pageref{app:prompt-mathir}}
\apxsub{\ref{app:prompt-pantry}}{PantryPlan}{\pageref{app:prompt-pantry}}
\apxsec{\ref{app:algorithm}}{Replay Algorithm}{\pageref{app:algorithm}}
\apxsub{\ref{app:binary-maxrl}}{Fresh task advantages}{\pageref{app:binary-maxrl}}
\apxsub{\ref{app:algorithm-update}}{State and optimizer update}{\pageref{app:algorithm-update}}
\apxsub{\ref{app:algorithm-bank}}{Admission, capacity, and deterministic scheduling}{\pageref{app:algorithm-bank}}
\apxsub{\ref{app:algorithm-control}}{Matched controls and recoverable state}{\pageref{app:algorithm-control}}
\apxsec{\ref{app:semantic-maxent}}{Fixed Semantic MaxEnt Comparator}{\pageref{app:semantic-maxent}}
\apxpart{Part III. Training results}
\apxsec{\ref{app:supporting}}{Empirical Results and Comparisons}{\pageref{app:supporting}}
\apxsub{\ref{app:scale-matrix}}{Terminal comparisons across scales}{\pageref{app:scale-matrix}}
\apxsub{\ref{app:level2-interim}}{Terminal comparison on the harder level}{\pageref{app:level2-interim}}
\apxsub{\ref{app:factorial-training-curves}}{Accuracy and verified modes throughout training}{\pageref{app:factorial-training-curves}}
\apxsub{\ref{app:ucpo-progress}}{Supporting comparators}{\pageref{app:ucpo-progress}}
\apxsub{\ref{app:diversity-comparators}}{Diversity-preserving comparators}{\pageref{app:diversity-comparators}}
\apxsub{\ref{app:frequency-replay-progress}}{Uniform versus fresh-frequency replay}{\pageref{app:frequency-replay-progress}}
\apxsec{\ref{app:semantic-current-factorial}}{Fixed Semantic-Entropy Factorials}{\pageref{app:semantic-current-factorial}}
\apxsec{\ref{app:ablations}}{Collapse and Replay Mechanism Checks}{\pageref{app:ablations}}
\apxsub{\ref{app:conditional-concentration}}{Conditional concentration from saved training outputs}{\pageref{app:conditional-concentration}}
\apxsub{\ref{app:pre-rl-regression}}{Breadth before and after RLVR-only training}{\pageref{app:pre-rl-regression}}
\apxsub{\ref{app:degenerate-groups}}{Fresh-gradient degeneracy and observed output collapse}{\pageref{app:degenerate-groups}}
\apxsub{\ref{app:fixed-bank-survival}}{Fixed-bank exemplar score trajectories}{\pageref{app:fixed-bank-survival}}
\apxpart{Part IV. What breadth represents}
\apxsec{\ref{app:prompt-hint-ablation}}{Matched Prompt-Hint Ablation}{\pageref{app:prompt-hint-ablation}}
\apxsub{\ref{sec:prompt-hint-ablation}}{Removing strategy hints from unchanged problems}{\pageref{sec:prompt-hint-ablation}}
\apxsec{\ref{app:decoding-objection}}{The Decoding Grid Behind the Objection}{\pageref{app:decoding-objection}}
\apxsec{\ref{app:sampling-budget-ablation}}{Sampling-Budget Ablation}{\pageref{app:sampling-budget-ablation}}
\apxsub{\ref{sec:discovery-curves}}{Protocol and estimands}{\pageref{sec:discovery-curves}}
\apxsub{\ref{app:hosted-discovery-complete}}{Hosted discovery: guidance, formatting, and late modes}{\pageref{app:hosted-discovery-complete}}
\apxsub{\ref{app:gpt56-all-levels-discovery}}{All five domains and all three levels at 512 draws}{\pageref{app:gpt56-all-levels-discovery}}
\apxsub{\ref{app:sampling-budget-effects}}{Observed sampling-budget effects}{\pageref{app:sampling-budget-effects}}
\apxsec{\ref{app:inference-followups}}{Inference Follow-ups}{\pageref{app:inference-followups}}
\apxsub{\ref{app:pantry-adaptation}}{Pantry portfolios under changed requirements}{\pageref{app:pantry-adaptation}}
\apxsub{\ref{app:cross-model-overlap}}{Complementarity across hosted models}{\pageref{app:cross-model-overlap}}
\apxsub{\ref{app:coarser-keys}}{Sensitivity to task-relevant equivalences}{\pageref{app:coarser-keys}}
\apxsec{\ref{app:portfolio-withdrawals}}{Portfolio Survival under Withdrawn Options}{\pageref{app:portfolio-withdrawals}}
\apxsub{\ref{app:withdrawal-severity}}{Two withdrawals at once}{\pageref{app:withdrawal-severity}}
\apxsub{\ref{app:withdrawal-decoding}}{Decoding settings as a substitute}{\pageref{app:withdrawal-decoding}}
\apxsub{\ref{app:withdrawal-recovery}}{Recovery after the change}{\pageref{app:withdrawal-recovery}}
\apxpart{Part V. Fixed deployments}
\apxsec{\ref{app:hosted-concentration}}{Inference-Only Concentration in a Hosted Model}{\pageref{app:hosted-concentration}}
\apxsub{\ref{app:hosted-mode-diversity}}{Hosted concentration on the success-conditional axis}{\pageref{app:hosted-mode-diversity}}
\apxsec{\ref{app:hosted-comparison}}{Comparison across Hosted Deployments}{\pageref{app:hosted-comparison}}
\apxsub{\ref{app:hosted-python-prompt}}{A separate Python prompt sensitivity test}{\pageref{app:hosted-python-prompt}}
\apxsub{\ref{app:hosted-temperature}}{A matched requested-temperature sensitivity test}{\pageref{app:hosted-temperature}}
\apxsub{\ref{app:gpt56-temperature}}{GPT-5.6 Sol: temperature, pass@8 and verified modes}{\pageref{app:gpt56-temperature}}
\apxsub{\ref{app:hosted-provider-outcomes}}{Native provider outcomes and undefined concentration}{\pageref{app:hosted-provider-outcomes}}
\apxsub{\ref{app:hosted-reasoning-control}}{Matched reasoning-control comparison}{\pageref{app:hosted-reasoning-control}}
\apxpart{Part VI. Idealized analysis}
\apxsec{\ref{app:theory}}{Mode Collapse and Verified Replay}{\pageref{app:theory}}
\apxsub{\ref{app:theory-grpo-mean}}{Categorical abstraction and mean updates}{\pageref{app:theory-grpo-mean}}
\apxsub{\ref{app:theory-collapse}}{Winner-take-all collapse in the categorical mean flow}{\pageref{app:theory-collapse}}
\apxsub{\ref{app:theory-gradient-availability}}{A replay signal when fresh groups supply none}{\pageref{app:theory-gradient-availability}}
\apxsub{\ref{app:theory-replay}}{Fixed-bank replay and qualitative retention}{\pageref{app:theory-replay}}
\apxsub{\ref{app:theory-exemplar-bridge}}{Complete responses and finite evaluation visibility}{\pageref{app:theory-exemplar-bridge}}
\apxsub{\ref{app:theory-entropy}}{Scope relative to entropy, admission, and neural optimization}{\pageref{app:theory-entropy}}
\apxpart{Part VII. Extended comparison}
\apxsec{\ref{app:related-extended}}{Diversity-Preserving RLVR: Extended Comparison}{\pageref{app:related-extended}}
\apxpart{Part VIII. Records}
\apxsec{\ref{app:reproducibility}}{Source Integrity and Reproducibility}{\pageref{app:reproducibility}}
\apxsub{\ref{app:source-integrity}}{Admissible sources and missing checkpoints}{\pageref{app:source-integrity}}
\apxsub{\ref{app:initial-provenance}}{Conflicted initial checkpoints, admitted by provenance}{\pageref{app:initial-provenance}}
\apxsub{\ref{app:repro-artifacts}}{Data, models, and scientific records}{\pageref{app:repro-artifacts}}
\apxsub{\ref{app:repro-recovery}}{Exact continuation and validation}{\pageref{app:repro-recovery}}
\par\endgroup

\clearpage
\phantomsection
\section*{Notation}
\label{app:notation}
Symbols are used consistently across the benchmark, the algorithm and the
idealized analysis. Several letters are conventionally overloaded --- $B$, $M$,
$m$ and $\tau$ each carry more than one meaning in this area --- so this table
is the paper's usage, and nothing below departs from it.

\begin{table}[H]
\centering
\caption{\textbf{Notation.} Symbols reused across the benchmark, the
algorithm, and the idealized analysis.}
\label{tab:notation}
\small
\setlength{\tabcolsep}{5pt}
\begin{tabularx}{\linewidth}{@{}lX@{}}
\toprule
\textbf{Symbol} & \textbf{Meaning} \\
\midrule
\multicolumn{2}{@{}l}{\emph{Benchmark and metrics}}\\
$V(y,s_x)$ & Validator: a canonical execution key, or $\bot$ for failure \\
$\mathcal C_x$ & Set of valid keys for prompt $x$ (never revealed to the learner) \\
$m$ & Number of correct modes a prompt admits, $m=|\mathcal C_x|$ \\
$P$ & Probability that a sampled response is verified correct \\
$q_c$ & Conditional probability of key $c$ given correctness \\
$Z_i$ & Canonical key of the $i$th draw, or $\bot$ \\
$C(q)$ & Conditional collision $\sum_c q_c^2$; the concentration measure \\
\pmd{} & Success-conditional breadth $1-C(q)$; the reported breadth axis \\
\neff{} & Effective number of verified modes, $1/(1-\pmd{})$ \\
$K$ & Draws per evaluation group ($K=8$ for every headline number) \\
$n_c$, $n_+$ & Draws verified on key $c$, and verified in all, $n_+=\sum_c n_c$ \\
$D_K$, $F_K$ & Distinct verified keys, and correct fraction, in $K$ draws \\
$\bar D_K$ & Expected distinct verified keys, $\E[D_K]=\texttt{distinct@}K$ \\
$B$ & Extra verified modes, $\texttt{distinct@}K-\texttt{pass@}K$ \\
\midrule
\multicolumn{2}{@{}l}{\emph{Algorithm}}\\
$\mathcal B_x$ & Replay bank for prompt $x$; $M=16$ its capacity \\
$k$ & Number of keys a bank holds, $k=|\mathcal B_x|\le M$; separately, the
  budget index in pass@$k$ \\
$P_{\mathcal B}$ & Total probability mass of banked keys \\
$e_c$ & Stored exemplar realizing key $c$; $|e_c|$ its token length \\
$\ell_c$ & Mean response-token log probability of the exemplar $e_c$ \\
$G$ & Fresh rollout group size ($G=16$ in training) \\
$R_i$, $R$ & Binary reward of response $i$, and group success count $\sum_iR_i$ \\
$A_i$ & Fresh task advantage of response $i$ \\
$\sigma$ & Serialized round-robin replay cursor \\
$\eta$, $\kappa$ & Scale and clip constants of the Semantic-MaxEnt
  comparator \\
\midrule
\multicolumn{2}{@{}l}{\emph{Idealized analysis}}\\
$\mathcal A$ & Category set, correct modes $\mathcal C$ and incorrect $\mathcal N$ \\
$c_G(P)$, $\Psi$ & Mean-flow correctness coefficient and its potential \\
$\rho$ & Effective replay dose in the mean flow; $\lambda_{\mathrm{replay}}$ its implemented weight \\
$\tau$ & Effective optimization time in the mean flow \\
$C_T$ & Energy constant in the retention bound \\
$\mu$, $\beta$ & Reference policy and its KL coefficient in (A4$'$) \\
$D$ & Reverse divergence $\mathrm{KL}(p\,\|\,\mu)$ to that reference \\
\bottomrule
\end{tabularx}
\end{table}


\clearpage
\phantomsection
\begin{center}\large\textbf{Part I. Scope and measurement}\end{center}
\addcontentsline{toc}{part}{Part I. Scope and measurement}

\section{Metric and Objective Details}
\label{app:metric-details}


Correctness and verified-mode diversity are distinct properties of a policy. For a
prompt $x$ with reference specification $s_x$, the validator maps a response
to either failure or a canonical key:
\[
  V(\cdot,s_x):\mathcal Y\rightarrow \mathcal C_x\cup\{\bot\}.
\]
Here $\mathcal C_x$ is the set of valid solution keys and $\bot$ denotes
failure. The learner observes a key only for a response it generates; neither
the full set $\mathcal C_x$ nor its size is revealed.

For $Y\sim\pi_\theta(\cdot\mid x)$, the policy induces both a total probability
of verified success and a distribution over how that probability is allocated:
\[
 P_\theta^+(x)=\Pr[V(Y,s_x)\neq\bot],
 \qquad
 p_\theta^+(c\mid x)=\Pr[V(Y,s_x)=c\mid V(Y,s_x)\neq\bot].
\]
The main text writes these two quantities as $P$ and $q_c$, dropping the
policy and the prompt where they are clear.
The conditional distribution is defined only when $P_\theta^+(x)>0$;
when $P_\theta^+(x)=0$, all three sampling metrics below are zero, while
\pmd{} is undefined rather than zero.
The first quantity measures correctness. The second describes the distribution
of verified behaviors conditional on being correct. We call concentration of
$p_\theta^+$ onto fewer keys \emph{mode collapse}, particularly when
$P_\theta^+$ remains stable or increases. That undefinedness sets a per-prompt
threshold, and the noise in an average taken over few such prompts sets two
more, which
App.~\ref{app:registered-endpoints} states: a prompt defines \pmd{} only from
two verified responses, a reported value rests on at least 30 such prompts,
and a paired difference, whose two sides are read on the same prompts, on at
least 20.
Moving correct probability onto keys that already hold more of it can only
lower $\pmd{}=1-\sum_c q_c^2$ (Lemma~\ref{lem:success-breadth}), so the
concentration defined here always registers as the fall in \pmd{}
\eqref{eq:main-pmd} that the main text reports. The converse does not follow:
\pmd{} summarizes an allocation in one number, so a fall in it need not arise
from any such ordered movement of mass. Neither statement claims that a key
reached probability zero, which finite samples cannot establish.

\subsection{Measuring success and diversity}
\label{app:metric-sampling}

Write $n_c$ for the number of a prompt's draws verified on key $c$, and
$n_+=\sum_c n_c$ for its verified draws in all. Where $n_+\ge2$, the estimator
of \eqref{eq:main-pmd} is the fraction of verified pairs that differ:
\begin{equation}
 \widehat{\pmd}
 =1-\frac{\sum_c \binom{n_c}{2}}{\binom{n_+}{2}}
 =1-\frac{\sum_c n_c(n_c-1)}{n_+(n_+-1)},
 \label{eq:main-pmd-estimator}
\end{equation}
which is unbiased for \eqref{eq:main-pmd} at every $n_+\ge2$: conditioning on
the number of verified draws leaves that many independent labels drawn from
$q$, so every ordered pair agrees with probability $\sum_c q_c^2$. What the
estimator runs on is the verified count $n_+$, not the draw budget: under the
pooling rule below a local prompt's four $K=8$ groups contribute 32 draws, a
hosted prompt contributes its eight, and $n_+$ counts the verified ones. The
identity is promptwise, and does not survive joint selection on other
conditions or dependent prompts. Two equivalent readings are sometimes useful:
\texttt{rarefied-distinct@}$2=1+\pmd$ is the expected number of distinct modes
in exactly two verified draws, which fixes sampling depth rather than success
count; and $\neff=1/(1-\pmd)$ is an effective number of verified modes. Both
are monotone transforms of \pmd{} and order policies identically.

Every printed \pmd{} is an average of that estimator. Where a protocol draws
more than one group per prompt, those draws are pooled before scoring --- $q$
is a property of the policy and prompt rather than of the draw budget, so
pooling estimates the same quantity from more pairs and defines it far more
often. The local evaluations draw four groups of eight and pool to 32; the
hosted deployments answer each prompt eight times, with nothing to pool. Either
way a reported value is the unweighted mean over the prompts where \pmd{} is
defined, carried with its across-prompt standard error. Prompts are weighted
equally on purpose. Weighting a prompt by how many verified draws it
contributed --- or, equivalently, pooling verified pairs across prompts before
taking the ratio, as a collision counter naturally does --- lets the prompts a
policy solves most often dominate the estimate, which reintroduces the accuracy
dependence \pmd{} exists to remove.

For $K$ independent samples $Y_{1:K}$, let
\[
  D_K=\bigl|\{V(Y_i,s_x):i\le K\}\setminus\{\bot\}\bigr|,
  \qquad
  F_K=K^{-1}\sum_i \1[V(Y_i,s_x)\neq\bot].
\]
\mb{} reports three sampling averages with different units. \texttt{distinct@}$K
=\E[D_K]$ is an expected count in $[0,K]$ of unique verified keys;
\texttt{pass@}$K=\Pr[D_K>0]$ and \texttt{mean@}$K=\E[F_K]$ are probabilities
in $[0,1]$.

\begin{lemma}[Success, sampled breadth, and \pmd{}]
\label{lem:success-breadth}
Fix a prompt, a finite or countable valid-key set, and $K\ge1$ independent
draws from the same fixed sampling law. Let
$\mu_c=\Pr[V(Y,s_x)=c]$ and $P=\sum_c\mu_c$. Then
\begin{equation}
 \texttt{pass@}K=1-(1-P)^K,\qquad
 \bar D_K:=\texttt{distinct@}K=\sum_c[1-(1-\mu_c)^K].
 \label{eq:success-breadth-identities}
\end{equation}
For a finite valid-key set of size $m\ge1$ and $P>0$, write $q_c=\mu_c/P$.
At fixed $P$, concentration lowers $\bar D_K$: if $q$ majorizes $q'$
(its decreasingly sorted partial sums are at least those of $q'$), then
$\bar D_K(P,q)\le\bar D_K(P,q')$. In particular,
\begin{equation}
 1-(1-P)^K\le\bar D_K(P,q)\le m[1-(1-P/m)^K].
 \label{eq:fixed-success-breadth-bounds}
\end{equation}
For $K\ge2$, the lower bound is attained only by concentration on one key,
and the upper bound only by uniform $q$. For $K=1$, $\bar D_1=P$ for every $q$.
At fixed $q$, increasing $P$ cannot decrease $\bar D_K$.
By contrast $\pmd(q)=1-\sum_c q_c^2$ does not involve $P$: it is unchanged when
correctness rises or falls at a fixed allocation. Concentration lowers it: if
$q$ majorizes $q'$ then $\pmd(q)\le\pmd(q')$, strictly unless $q'$ is a
permutation of $q$. It runs from $0$ at a point mass to
$1-1/m$ at the uniform allocation, so its ceiling is set by the number of
valid keys rather than by correctness.
\end{lemma}
\begin{proof}
All $K$ draws fail with probability $(1-P)^K$. For each key $c$, its
appearance indicator has expectation $1-(1-\mu_c)^K$; summing indicators
gives the second identity, also for countably many keys by nonnegativity.
For fixed $P>0$ and $K\ge2$, the function
$f(v)=1-(1-Pv)^K$ is strictly concave on $[0,1]$.
Transferring mass from a larger coordinate to a smaller one without reversing
their order increases $\sum_c f(q_c)$; successive such transfers give the
majorization assertion. The point mass and uniform vector are its extreme
cases, giving the bounds and equality conditions. Finally,
$\partial\bar D_K/\partial P=K\sum_c q_c(1-Pq_c)^{K-1}\ge0$,
which proves monotonicity in correctness.
For the final claim, $\sum_c q_c^2$ is symmetric and convex, so the
same transfer strictly decreases it whenever the two coordinates differ;
$\pmd$ is one minus that sum and no term in it involves $P$. The point mass and
the uniform vector again give its extremes, by the same succession of transfers
that gave the bounds on $\bar D_K$.
\end{proof}

These are promptwise identities, before averaging over prompts. In particular,
applying the nonlinear formula to a mean correctness generally changes the
result. They explain the shared measurement in the training and hosted
studies: success depends on total correct mass; sampled breadth also depends
on its allocation.

Raw \texttt{distinct@}$K$ measures sampled coverage, not the total number of
positive-probability modes. For an uninformed reference, suppose a finite
space has $N$ actions and $N_c$ actions map to valid key $c$. Uniform
independent action sampling then gives $\sum_c[1-(1-N_c/N)^K]$ expected
verified keys. The simpler expression $m[1-(1-1/N)^K]$ applies when the $m$
valid actions have distinct keys, as in our reported finite-action references.
Open-ended program and derivation spaces have no canonical nominal
``uniform'' generator. Table~\ref{tab:absolute-support-references} applies these
references to the trained policies.

\subsection{Endpoints and the reported breadth axis}
\label{app:registered-endpoints}

Two terminal endpoints are registered as co-primary, meaning neither is
subordinate to the other and both are carried for every
model--domain--level block: \texttt{pass@}$K$, the probability of finding any
correct output, and raw \texttt{distinct@}$K$, the number of verified keys
found. \texttt{pass@}$K$ is printed in the tables here and in
App.~\ref{app:supporting}; \distk{} is retained beside each of them in the
machine-readable records, which is where it is read, because it rises with
correctness and so restates accuracy wherever it is printed as breadth.

A separate question is which statistic should carry the breadth axis in the
figures and the prose. Raw \texttt{distinct@}$K$ is coupled to the number of
successful draws: by \eqref{eq:success-breadth-identities} it grows with $P$
at fixed $q$ (Lemma~\ref{lem:success-breadth}), and empirically it correlates
$\MDdistinctcorr$ with \texttt{pass@8} across the \MDcells{} untrained
measurements, while \pmd{} correlates $\MDpmdcorr$ on the \MDreportable{} of
them with sufficient support. The two coefficients are not read on the same
population: definedness is itself accuracy-dependent, as the support paragraph
below sets out. Reported as breadth, raw \texttt{distinct@}$K$ therefore
largely restates correctness, and the comparison it invites between a
high-\texttt{distinct@8} domain and a
low-\texttt{distinct@8} one is mostly a comparison of accuracy. We
consequently report \pmd{} \eqref{eq:main-pmd} as the breadth axis, which is a
functional of $q$ alone (Lemma~\ref{lem:success-breadth}).
Fig.~\ref{fig:level-construction} shows what that separation buys.
From 7B up, Python is answered correctly on at least \MDpyPasslo{} of held-out
prompts at Levels 1--3 while its \pmd{} never exceeds \MDpyPmdhi{}:
near-perfect accuracy reached by returning the same answer almost every time.
Read on the \texttt{distinct@8} axis that same block would look like the
broadest domain in the table.

\paragraph{Support and explicit gaps.}
\pmd{} is undefined for a prompt returning fewer than two verified responses,
and definedness is itself accuracy-dependent: the fraction of prompts defining
\pmd{} correlates $\MDdefinedcorr$ with \texttt{pass@8}. This is an information
limit rather than an artifact, since a policy that almost never succeeds
exhibits no distribution of successes to measure. We mark it everywhere it
applies.
Every reported \pmd{} carries the number of prompts that define it, and
measurements resting on fewer than 30 such prompts sit below the support bar.
Figures draw those hollow rather than leaving them empty, from five defining
prompts upward, so a weak measurement shows up as a thin estimate carrying its
own error bar rather than as an absence. Below five nothing is drawn: a
standard error of zero on two or three prompts records those prompts agreeing,
not a precise estimate, and a mark there would read as the firmer measurement
it is not. Of the
\MDcells{} measurements of frozen base models across all four families,
\MDreportable{} clear the 30-prompt bar,
and the gaps concentrate where success is rarest: \MDtopgapcount{} of
the \MDgaps{} are \MDtopgapdomain{}, whose frozen success rate is too low to
support any breadth statement. A hollow or absent entry means breadth was
not reliably measurable there; it does not mean breadth was low.

A paired difference carries a separate bar of 20, because both of its sides are
measured on the same prompts and only their difference is read. The distinction
is not a relaxation for convenience: the side that falls short is almost always
the RLVR-only control, whose definedness is depressed by the very concentration
under study, so holding a difference to the stricter bar would censor exactly
the comparisons where replay separates most. A cell with either side under 20
remains a gap.

\pmd{} is not a new measurement of the saved responses. It is exactly one minus
the conditional-collision U-statistic of
App.~\ref{app:conditional-concentration}; we report the diversity orientation
because a metric that increases with breadth reads more directly alongside
\texttt{pass@8}. Every \pmd{} value
is recomputed from canonical keys already stored with the saved responses,
without resampling or regrading.

The derived contrast $\texttt{distinct@}K-\texttt{pass@}K=\E[D_K-\1[D_K>0]]$
estimates the number of verified modes beyond the first. It is computable
without knowing the full valid set $\mathcal C_x$ but remains coupled to the
number of successful draws, so we keep it as an interpretive decomposition
rather than as a diversity measure. Because every
headline evaluation fixes $K=8$, reporting $\texttt{distinct@8}/8$ would add no
information beyond a linear rescaling; we retain the raw count and label its
units explicitly.


\subsection{Why balance within the replay bank}
\label{app:replay-metric-alignment}

Uniform sampling over keys prevents discovery frequency from determining replay
frequency. To explain its relationship to the evaluation metric we analyze an
equal-weight categorical surrogate, without length normalization
(Fig.~\ref{fig:replay-bank-balance} draws what the two terms preserve);
App.~\ref{app:algorithm} states the objective as implemented.
A stored response's probability need not equal its key's total probability.
For $k=|\mathcal B_x|$ banked keys, write $p(c)$ for their probabilities and
$P_{\mathcal B}$ for their total mass. Then
\begin{equation}
 -\frac1k\sum_{c\in\mathcal B_x}\log p(c)
 =-\log P_{\mathcal B}+\log k+
   \mathrm{KL}\!\left(U_k\,\|\,
   p(c\mid c\in\mathcal B_x)\right).
 \label{eq:replay-bank-decomposition}
\end{equation}
The objective rewards greater banked verified mass and a more balanced
conditional distribution, aligning with \texttt{pass@}$K$ and
\texttt{distinct@}$K$. The combined update need not improve both at every step:
banked-key coverage contributes $\sum_{c\in\mathcal B_x}[1-(1-p(c))^K]$ to
expected \texttt{distinct@}$K$ and, at fixed $P_{\mathcal B}$ and $k$, concavity
makes it largest under uniform allocation, while $1-(1-P_{\mathcal B})^K$, the
probability of any banked key, depends only on total mass
(Fig.~\ref{fig:replay-bank-decomposition}). Uniformity also
maximizes the smallest conditional banked-key mass, making equal retention
explicit.

\begin{figure}[!htbp]
  \centering
  \includegraphics[width=.86\linewidth]{figures/replay_bank_decomposition.pdf}
  \caption{\textbf{The decomposition against measured prompts, by domain.}
  Bands are the coverage sum above with the banked restriction dropped, over
  all of a prompt's correct keys: $\sum_c[1-(1-Pq_c)^{8}]$, where $p(c)=Pq_c$.
  They are evaluated over allocations from a point mass to uniform; at each
  breadth a band spans only the key counts that can attain it, over the range
  those prompts realize.
  Points are \MDbankprompts{} held-out prompts returning at least two verified
  responses, at the terminal checkpoint of the Qwen2.5-0.5B replay arms, binned
  by measured verified mass $P$. Panels are the five domains and carry the
  tints of their cards in Fig.~\ref{fig:modebench-examples}, because a mode
  means something different in each and the number of keys a prompt can occupy
  is a property of the domain. Graph and Pantry occupy the full range; Python
  and MathIR sit at the left edge, MathIR almost entirely so --- nearly every
  one of its prompts returns a single key, where two correct answers are the
  same answer, because its certified modes are derivations of one answer
  rather than alternatives. These prompts are held out and so carry no bank; the keys a
  policy occupies here stand in for what a training prompt's bank would have to
  hold. That count stays far below the capacity $M=16$: what bounds these
  policies is how many modes they find, not how many the bank can hold, so
  raising $M$ alone would not relieve what binds them.}
  \label{fig:replay-bank-decomposition}
\end{figure}

That alignment is deliberate, which is a reason for caution rather than a
result: the bank objective raises \texttt{distinct@}$K$ by construction, so
\texttt{distinct@}$K$ cannot also serve as the evidence that replay worked. The
objective is aligned with \texttt{pass@}$K$ as well, and what makes that
endpoint evidence rather than recitation is where it is measured: banks hold
only training-prompt responses, and evaluation responses never enter training.
\pmd{} carries the breadth axis because the objective does not target it at
all.


% Shared long-paper and workshop benchmark/protocol appendix.

\clearpage
\phantomsection
\begin{center}\large\textbf{Part II. Benchmark and method}\end{center}
\addcontentsline{toc}{part}{Part II. Benchmark and method}

\section{Benchmark and Evaluation Protocol}
\label{app:data}
\label{app:experimental-design}
\label{app:domain-overview}
\label{app:data-inventory}

\mb{} takes correctness and identity from one executable validator, the map
$V(\cdot,s_x)$ of App.~\ref{app:metric-details}: invalid outputs receive reward
zero and no key, valid outputs reward one and a canonical execution key. What
matters for the protocol is who may see the catalogue, and the answer is nobody:
neither its members nor its size reach training, replay, scheduling, or
stopping.
Table~\ref{tab:tasks} specifies the five domains and which surface forms merge.
The \href{https://huggingface.co/datasets/od2961/ModeBench}{public dataset}
preserves each domain's split roles and construction identity. The training
comparisons reported here use Levels 1 to 3.

\begin{table}[H]
  \centering
  \caption{\textbf{\mb{} at Level 1.} Correctness and identity come from
  one executable path; training starts without the mode catalogue.
  \emph{Valid modes} is the mean number of correct execution modes per
  evaluation prompt with the per-prompt range beside it, exact over the
  released Level-1 multi-answer split. \emph{Reached} is what one frontier
  deployment actually produces on the same construction --- distinct verified
  modes per prompt at \MDreachedDraws{} draws, over the \MDreachedPrompts{}
  problems per domain of Fig.~\ref{fig:gpt56-sampling-budget}, an
  outcome-independent subset of the same evaluation split --- so the two
  columns separate the modes a domain admits from the modes anything has been
  observed to give on the same prompts. MathIR is the domain where that gap is structural rather
  than a sampling limit, because its modes are derivations of one answer and
  most of the certified ones are longer routes. Each higher level is a separate
  construction with its own support distribution
  (App.~\ref{app:data-levels}). The exact prompt each policy receives is
  reproduced verbatim in App.~\ref{app:prompts} for a direct, byte-level audit
  of the evaluation interface, chat-template wrapper included.}
  \label{tab:tasks}
  \scriptsize
  \setlength{\tabcolsep}{5pt}
  \renewcommand{\arraystretch}{1.08}
  \begin{tabularx}{\linewidth}{@{}l
      >{\hsize=1.3\hsize}Y
      >{\hsize=.8\hsize}Y
      >{\hsize=.9\hsize}Y
      rr@{}}
    \toprule
    \textbf{Domain} & \textbf{Generated answer / check} &
    \textbf{Executed key} & \textbf{Aliases that merge} &
    \textbf{Valid modes} & \textbf{Reached} \\
    \midrule
    Graph coloring & Complete assignment; check fixed colors and edges &
    Color vector & Spacing only & \MDvmGraphMean{} & \MDreachedGraph{} \\
      & & & & {\scriptsize\MDvmGraphRange} & \\
    \addlinespace[2pt]
    Countdown & Parse operands; execute exactly &
    Normalized AST & Commutative order & \MDvmCountdownMean{} & \MDreachedCountdown{} \\
      & & & & {\scriptsize\MDvmCountdownRange} & \\
    \addlinespace[2pt]
    Python factors & Execute restricted lambda externally &
    Return vector & Source programs with same behavior & \MDvmPythonMean{} & \MDreachedPython{} \\
      & & & & {\scriptsize\MDvmPythonRange} & \\
    \addlinespace[2pt]
    MathIR & Execute bounded action program &
    State trajectory & Programs with the same states & \MDvmMathIRMean{} & \MDreachedMathIR{} \\
      & & & & {\scriptsize\MDvmMathIRRange} & \\
    \addlinespace[2pt]
    PantryPlan & Exact ingredient allocation; check all bounds &
    Ingredient support & Quantity and ordering aliases & \MDvmPantryMean{} & \MDreachedPantry{} \\
      & & & & {\scriptsize\MDvmPantryRange} & \\
    \bottomrule
\end{tabularx}
\end{table}


\subsection{Split identities and executable validation}
\label{app:data-materialization}

Every domain's training and evaluation splits are disjoint, at the sizes
Table~\ref{tab:run-contract} records. PantryPlan and the Level-2 constructions
carry development splits besides, which no reported result reads
(App.~\ref{app:data-levels}).
Python materialization independently executes at least two functions with
different valid keys for every row; MathIR executes every certified route
through the same interpreter used for model outputs.

\begin{table}[H]
\centering
\caption{\textbf{Deterministic materialization identities.} The two domains
whose materialization is itself executable, as
App.~\ref{app:data-materialization} describes: Python runs
functions, MathIR runs certified routes. Hashes bind the released rows to those
checks, and the full hashes and construction records accompany the data.
PantryPlan's row hashes follow below. Exact modes per prompt are the training
split here; Table~\ref{tab:tasks} reports the evaluation split, which is
disjoint from this one.}
\label{tab:materialization-identities}
\small
\begin{tabular}{@{}lrr@{}}
\toprule
& \textbf{Python factors} & \textbf{MathIR action menu} \\
\midrule
Training rows & 384 & 384 \\
Evaluation rows & 128 & 128 \\
Exact modes per prompt (train) & \MDvmPythonTrainRange & \MDvmMathIRTrainRange \\
Training-row SHA-256 & \texttt{bcfa9acc\ldots37cf7} & \texttt{bfc98362\ldots5d3835} \\
Evaluation-row SHA-256 & \texttt{be0f621c\ldots94d183} & \texttt{df5783dd\ldots328b96} \\
\bottomrule
\end{tabular}
\end{table}

\paragraph{PantryPlan.}
\label{app:data-pantry}
The deterministic v2 release has 384/64/128 training/development/evaluation
prompts across four formulation families. Exact support ranges from
\MDvmPantryTrainRange{} over the training split and \MDvmPantryRange{} over
the 128 evaluation prompts, which
average \MDvmPantryMean{} feasible supports. Training and evaluation row hashes are
\texttt{a4883c5e\ldots a757f} and \texttt{5d81ba8d\ldots5923}.
The validator checks exact inventory, serving steps, mass, nutrient bounds,
and dietary tags; identity is ingredient support. A verified allocation is
not a claim of palatability, cookability, or safety.

\subsection{Python execution contract}
\label{app:python}
\label{app:python-language}

Candidates are stripped, limited to 240 characters and 64 AST nodes, and
parsed in expression mode. The only accepted form is \texttt{lambda n: EXPR}
with exactly one ordinary argument named \texttt{n}: no defaults,
positional-only, keyword, or variadic arguments. Integer literals have
absolute value at most 10,000. The expression grammar permits \texttt{n},
\texttt{+}, \texttt{-}, \texttt{*}, floor division, remainder, integer
comparisons, Boolean \texttt{and}/\texttt{or}, unary plus/minus/\texttt{not},
and conditionals. Boolean and non-integer literals, calls, imports,
attributes, comprehensions, containers, subscripts, assignment, loops, and
other identifiers are rejected. The specification format admits $J$ distinct
integers in $[4,1000]$, with $J$ between two and eight. Every released Level-1
row uses \MDpyCases{} of them, drawn from $[\MDpyCaseLo,\MDpyCaseHi]$, and
carries at least two valid behavior vectors.

\paragraph{External execution.}
\label{app:python-process}
The trainer never executes candidate code. A locked parent client launches
an isolated Python \texttt{-I} process in a fresh session, with a persistent
JSON-lines worker accepting only restricted source and trusted task JSON.
The worker repeats the AST checks, compiles the expression with empty
\texttt{\_\_builtins\_\_}, and executes the frozen inputs under a
0.25-second real-time alarm; the parent enforces a separate one-second
response deadline. A valid response contains the integer vector and key.
Timeouts, empty or malformed responses, broken pipes, non-integer outputs,
and failed divisor checks yield $\bot$ and zero reward. Broken pipes,
operating-system errors, and malformed responses permit one clean worker
restart; a timed-out request remains failed. Worker stdout is never parsed
as a candidate answer, and exceptions never produce a key.

\paragraph{Behavioral identity.}
\label{app:python-identity}
For each of the $J$ frozen inputs $n_j$, the returned $d_j$ must be an integer
but not a Boolean, satisfy $1<d_j<n_j$, and divide $n_j$ exactly. All checks
must pass. The key is the ordered vector
$\texttt{python\_factor:}d_1\texttt{,}\cdots\texttt{,}d_J$;
programs merge exactly when their vectors agree. The exact support count
$\prod_j|\{d:1<d<n_j,\ d\mid n_j\}|$ is used only for dataset validation.

\subsection{Matched level constructions}
\label{app:data-levels}

Levels 2 to 5 are separate constructions, each calibrated against a fixed
Level-1 reference at a larger model, and each admitted on its own terms. What
they share is the verifier, the canonicalizer and the notion of a mode; what
grows is the structure to be searched.

\paragraph{Level 2.}
Level 2 changes problem structure while preserving the verifier and
canonicalizer. Its 384/128/128 training/development/evaluation rows have
disjoint identities across the Level-1 and Level-2 splits. Training histograms
match the native Level-1 training sets. Development and evaluation histograms
match designated Level-1 development and confirmation reserves in Graph,
Countdown, Python, and MathIR; Pantry uses native references, with the counts
of its 64-row Level-1 development histogram scaled by two to give the 128-row
target that its Level-2 development split is built to match.

These construction reserves are not the native Level-1 terminal test sets.
The displayed terminal evaluations have different valid-key histograms
between levels in Graph, Countdown, and Python. Cross-level changes therefore
combine task structure and test population; they do not isolate difficulty
at fixed test-set support. Replay comparisons remain paired within each
level on the same prompts, seeds, fresh objective, and training budget.

\begin{table}[H]
\centering
\caption{\textbf{What Level 2 changes.} The verifier, the canonicalizer and the
notion of a mode are the same at both levels; only the structure to be searched
grows.}
\label{tab:level2-structure}
\small
\begin{tabularx}{.96\linewidth}{@{}lYY@{}}
\toprule
\textbf{Domain} & \textbf{Level 1} & \textbf{Level 2 structural change} \\
\midrule
Graph coloring & Three hidden vertices & Four hidden vertices \\
Countdown & Three operands & Four operands, with values up to 14 \\
Python factors & Four executed cases, all at most 96 & Four cases, including a value above 96 \\
MathIR & One-sided linear families & Variable-on-both-sides families \\
PantryPlan & Base dietary/composition constraints & More active dietary exclusions and tighter composition \\
\bottomrule
\end{tabularx}
\end{table}

Before treatment training, frozen Qwen2.5-0.5B development evaluation admitted
a Level-2 domain only when \texttt{pass@8} lay in $[.10,.90]$ and was strictly below
its construction-reference Level-1 value under identical decoding. The matched-level
identity manifest records support histograms, row hashes, construction
seeds, disjointness checks, and these admission receipts.


\paragraph{Scale-calibrated Level 3.}
Level 3 uses frozen Qwen2.5-3B against the same fixed Qwen2.5-0.5B Level-1
reference. Recipes are selected on development problems before fresh held-out
confirmation, requiring absolute pass@1 and pass@8 differences of at most
$.04$ and $.08$ in every domain, with four independent groups of eight samples
per prompt. \distk{} is reported but is not a fitting objective. The
interface uses native chat, temperature 1, top-$p$ 1, and 192 generated
tokens: boxed-direct answers for Graph and the hybrid solver with
legal-syntax guidance for the other domains. This matches fixed measured
targets empirically; it is not a two-sided equivalence test.

The Python Level-3 interface used in the hint ablation removes the strategy
hints (App.~\ref{app:prompt-hint-ablation}), so its difficulty is not
covered by that match. All reported results use the original datasets and
prompts.

\paragraph{Levels 4 and 5.}
Both continue the level sequence against the same fixed Qwen2.5-0.5B Level-1 reference
and the same $.04$ and $.08$ tolerances on pass@1 and pass@8, calibrated at 7B
and 14B. Level 5 is difficulty-matched in all five domains and carries a
composite record naming each domain's source and protocol digests. Level 4 is
admitted with a recorded exception. Four of its domains --- Graph, Countdown,
Python and PantryPlan --- are difficulty-matched; its MathIR construction is
not, for a reason that belongs to the target rather than to the search for it. Matching there demands a pass@8/pass@1 ratio of $5.50$; the best
ratio any construction in this family attains at 7B is $4.16$, against a ceiling
of $8$ for a homogeneous population, and a ratio of averages cannot exceed its
best component. No mixture of these constructions reaches the target, so the
target sits outside the achievable set: MathIR at 7B is far more bimodal than
its Level-1 reference, with measured heterogeneity gaps of $.169$--$.398$
against the reference's $.060$, and three further difficulty mechanisms measured
inert. All five Level-4 datasets are frozen at 384/128/128, verified, and
disjoint from every other level, and the level is admitted on that basis: a
target proved unreachable is a property of the construction, not a gap in it.
What admission does not extend to is a matched-difficulty claim for MathIR,
which its entries carry as an exception rather than a match.
Figs.~\ref{fig:base-levels-all-scales} and~\ref{fig:level-construction} plot
both levels without marking that exception, and should be read against this
paragraph.

\subsection{Frozen base models across all five levels}
\label{app:base-levels-all-scales}

Fig.~\ref{fig:base-levels-all-scales} reports the frozen Qwen2.5-Instruct
scales across all five levels and five domains;
Fig.~\ref{fig:families-appendix} adds the other three families, for
\MDcells{} completed measurements in all. Every model uses the same 128 held-out
prompts within a domain and level, with four independent groups of eight draws
per prompt, so each plotted point rests on 4,096 responses. Each group contributes an indicator
of any verified output to pass@8 and the number of distinct verified canonical
keys to the registered \distk{}; we average groups within prompts and then
average prompts. \pmd{} instead pools a prompt's four groups before scoring it,
under the rule and for the reason given in App.~\ref{app:metric-sampling}.
For \texttt{pass@8} and \distk{} the four groups stay separate $K=8$
observations throughout; only \pmd{} pools them.

These measurements use native chat, $T=1$, top-$p=1$, and a 192-token output
limit, with boxed-direct Graph answers and the registered hybrid interface
elsewhere. Every constructed level is measured on the prompts it was admitted
with, unchanged. That interface is not the one the training runs evaluate
through, so this grid and the untrained endpoint of a training comparison are
two measurements of the same weights rather than one number reported twice.
They can differ by a lot: on Python at Level 1 this grid answers most prompts,
where the same weights under the training template and the resampling's wider
draw budget clear two verified answers on 73 of 128 and sit at \pmd{} $.168$,
the figure Section~\ref{sec:results-collapse} quotes. Read a grid row against
other grid rows, and a training contrast against its own untrained arm; the
two are not interchangeable. These held-out
measurements are separate from development fitting and do not imply that
success decreases monotonically with level for every model and domain:
\texttt{pass@8} ends below its Level-1 value in \MDladderfall{} of the
\MDladdersupport{} model--domain series the levels support, not in all of them.
Moving along the two axes also costs breadth by different amounts. Taking the
median range of \pmd{} along a row of the grid, scale moves it \MDscalespan{}
at fixed level against \MDlevelspan{} for level at fixed scale, so the level axis and the parameter
count are not interchangeable ways of narrowing how a model answers.

\begin{figure}[!htbp]
  \centering
  \includegraphics[width=.62\linewidth]{figures/qwen_level_trends.pdf}
  \caption{\textbf{The ladder read along one axis, for one family.}
  Both axes averaged equally over the five domains, one line per frozen
  Qwen2.5 scale in the colours of Fig.~\ref{fig:base-levels-all-scales}, with
  GPT-5.6 Sol dashed in black. Left: \texttt{pass@8} falls with level at every
  scale while the hosted deployment stays near one. Right: \pmd{} over the same
  cells, pooling the held-out prompts with the training split where that split
  has been measured. The two panels must be read together rather than
  separately. \pmd{} exists only where a cell clears the 30-prompt support bar,
  and of the \QLTpoints{} scale--level points on the right only
  \QLTsupported{} rest on five domains that all clear it; the rest average
  whichever domains remain, and are drawn hollow. Those thinned points put the
  smallest scales highest at the hard levels, which follows the support loss
  that comes with falling accuracy on the left rather than any gain in breadth.
  Breadth per domain, where it is measurable without that averaging, is in
  Figs.~\ref{fig:base-levels-all-scales} and~\ref{fig:level-construction}.}
  \label{fig:qwen-level-trends}
\end{figure}

\begin{figure}[!htbp]
  \centering
  \includegraphics[width=\linewidth]{figures/mode_diversity_level_construction.pdf}
  \caption{\textbf{Every frozen scale on its own.}
  Each row is one scale drawn without the others over it, across every family
  and every hosted deployment of Fig.~\ref{fig:families-appendix}; columns
  are the five domains and marker shape identifies the level. Each point pools four independent eight-draw groups on
  128 held-out prompts at $T=1$. Reading across a row gives one model over all
  five domains; reading down a column gives one domain over every scale.
  The two axes separate cleanly at 7B: Python answers almost every prompt
  correctly and almost always with the same mode (\texttt{pass@8}
  $\MDsevenbPypasslo$--$\MDsevenbPypasshi$, \pmd{}
  $\MDsevenbPypmdlo$--$\MDsevenbPypmdhi$), while Graph reaches \texttt{pass@8}
  $\MDsevenbGraphpasslo$--$\MDsevenbGraphpasshi$ with \pmd{}
  $\MDsevenbGraphpmdlo$--$\MDsevenbGraphpmdhi$. A filled mark clears the
  30-prompt bar; a hollow one carries its estimate and standard error from five
  to twenty-nine prompts, which is thin rather than absent and is drawn so it
  can be read as thin. Ticks below an axis mark measurements resting on fewer
  than five prompts, where a standard error of zero would record two or three
  prompts agreeing rather than a precise estimate; those are given no height. The smallest scales are mostly such: the measurable entries
  number \MDtinyAentries{} at \MDtinyAlabel{} and \MDtinyBentries{} at
  \MDtinyBlabel{}, so their rows are largely empty by measurement rather
  than by omission. The hosted frontier deployments
  carry no parameter count and so belong to no row; they appear in
  Figs.~\ref{fig:base-levels-all-scales} and~\ref{fig:families-appendix},
  where breadth is read against level rather than against scale.}
  \label{fig:level-construction}
\end{figure}

\begin{figure}[!htbp]
  \centering
  \includegraphics[width=\linewidth]{figures/frontier_level_grid.pdf}
  \caption{\textbf{Harder levels narrow the frontier.} Success-conditional
  breadth against the level axis, one line per hosted deployment, over 128 held-out
  prompts and eight stateless requests per prompt per cell. Colour runs with
  collapse, darkest most concentrated. Only the rungs every deployment covers
  are drawn: Levels 4 and 5 exist for one deployment alone, and a line running
  past its neighbours would read as that deployment holding up where the others
  fall rather than as the only one measured there
  (App.~\ref{app:hosted-mode-diversity} reports those two probes separately).
  Correctness is not drawn, because there is nothing to see --- every hosted
  cell sits at \texttt{pass@8} $\MDfrontierminpass$ or above, so the levels cost these deployments
  accuracy they do not have to give and breadth they do. The Average panel
  weights domains equally over a set held fixed for each deployment, so the
  domain mix cannot move with the level; on it every one of the
  \MDfrontierdeployments{} falls from Level 1 to Level 3, by
  $\MDfrontierdroplo$ to $\MDfrontierdrophi$. The domains do not agree on
  how: MathIR runs to zero, Python and Graph fall, and Countdown does not
  narrow at all. Open marks sit below the support bar and stay off the line,
  as in Fig.~\ref{fig:concentration-levels}. The higher levels also add
  strategy guidance to the MathIR, Python and PantryPlan prompts that Level 1
  does not carry, so part of the fall on those three is instruction-following
  rather than difficulty; Graph and Countdown add none, and every deployment
  falls on those two alone (Apps.~\ref{app:prompt-hint-ablation}
  and~\ref{app:hosted-mode-diversity}).}
  \label{fig:frontier-level-grid}
\end{figure}

\begin{figure}[!htbp]
  \centering
  \includegraphics[width=.92\linewidth]{figures/mode_diversity_families_appendix.pdf}
  \caption{\textbf{Four model families on the axes of
  Fig.~\ref{fig:base-levels-all-scales}.} Frozen instruction-tuned models
  from SmolLM2, Qwen2.5, Falcon3 and OLMo-2, plotted as \texttt{pass@8} against
  \pmd{}; colour runs green to red with parameter count on a log scale, and
  marker shape identifies the level. This figure carries every model family and
  every level collected, where Fig.~\ref{fig:base-levels-all-scales} shows
  Qwen2.5 alone. Grey markers are all
  \MDfrontierdeployments{} hosted deployments of
  Section~\ref{sec:hosted-concentration}, of which the main-body figure draws
  only GPT-5.6 Sol; they have no parameter count and so sit off that ramp.
  Fig.~\ref{fig:level-construction} separates the same data one scale per
  row. Ticks below each axis mark measurements whose success is too rare
  for \pmd{} to be estimated on enough prompts; those are given no
  height. Because colour encodes size rather than family, models of different
  families at the same scale share a colour and can be compared directly.}
  \label{fig:families-appendix}
\end{figure}


\subsection{Other model families at matched scale}
\label{sec:families-appendix}

The main text holds the training recipe fixed and reads the scale axis within
Qwen2.5. Fig.~\ref{fig:families-appendix} instead places SmolLM2, Falcon3 and
OLMo-2 beside those scales on the same axes, under the identical protocol: the
same 128 held-out prompts per domain and level, the same four groups of eight
draws, the same verifier and the same support requirement. The comparison is
taken on the \MDcommonlevels{} levels every family covers, since only Qwen2.5
reaches Level 5 and a family average over more levels would put the level
axis inside a contrast about recipes. Parameter count sets
colour, so two families at the same size take the same colour and are read
against each other directly.

At matched parameter count the families do not agree on breadth. Near 7B,
where Qwen2.5, Falcon3 and OLMo-2 each place a model and SmolLM2 does not, the
three separate on \pmd{} while their \texttt{pass@8} stays close; the ordering
holds into the 10--15B band, though Falcon3 and OLMo-2 sit close enough there to
read as a pair rather than as two ranks. Size therefore does not fix how
many ways a model answers: at matched scale a change of recipe moves \pmd{} by a
median \MDrecipespan{}, against \MDwithinrecipespan{} for a family's entire
scale range --- a same-size change of recipe is worth about as much breadth as
the whole parameter range within one. Both medians are taken within a family
over the \MDcommonlevels{} shared levels, so they are a matched pair; the
\MDscalespan{} quoted above is a different statistic, read across the Qwen
grid's seven scales at every level. We report this as a measured property
of the released checkpoints, not as a claim about which pretraining or
post-training choice is responsible, which these frozen measurements cannot
separate.


\subsection{Training and evaluation settings}
\label{app:repro-run-contract}
\label{app:comparison-design}

% [H] pinned this table to the head of the subsection, and being too tall for
% what the previous page had left it took a page of its own and stranded the
% remainder. Floating lets the prose that follows it close the gap.
\begin{table}[!htbp]
  \centering
  \caption{\textbf{Matched run contract.} Settings and systems changes are
  shared by methods within each model--domain--seed comparison, including
  the zero-weight replay control. Optimizer updates and fresh-rollout budget
  are therefore identical across arms, and replay buys no additional sampling.
  What it does add, and what is not matched away, is one forward and one
  backward pass over at most 16 short banked exemplars on each visited prompt.}
  \label{tab:run-contract}
  \small
  \begin{tabularx}{\linewidth}{@{}P{.23\linewidth}Y@{}}
    \toprule
    \textbf{Component} & \textbf{Setting} \\
    \midrule
    Data & 384 training and 128 disjoint evaluation prompts per domain and
      level; no evaluation row enters training or replay. \\
    Fresh rollouts & $G=16$, temperature 1, top-$p=1$; one prompt group per update. \\
    Optimization & Eight passes (3,072 updates), one PPO epoch,
      $\beta_{\mathrm{KL}}=0$; paired initialization, prompt stream,
      optimizer, hardware class, and checkpoint cadence. \\
    Evaluation & Passes $0,.5,\ldots,8$; four fixed temperature-1,
      top-$p=1$, $K=8$ draws per prompt, plus a separate temperature-0
      greedy trace. No best-checkpoint selection. \\
    Replay & Capacity $M=16$ per prompt; one round-robin bank per update;
      $\lambda_{\mathrm{replay}}=.10$. Controls execute the same code with
      an exact-zero replay derivative; realized banks and workloads can differ. \\
    Training seeds & Qwen2.5-0.5B: 43--47; Falcon3-1B: 55--59;
      Qwen2.5-3B: 70--74. \\
    Learning rate & Qwen-0.5B and Falcon-1B: constant $2\times10^{-7}$,
      with optimizer-state offload for Falcon. Qwen-3B: cosine decay with
      $10^{-7}$ peak, $10^{-8}$ floor, and 10\% warmup. \\
    Optimizer & Fused Adam in \texttt{bf16} with decoupled weight decay,
      $\beta_1=.9$, $\beta_2=.95$, weight decay $0$, gradient-norm clip $1.0$,
      gradient checkpointing on. Offload swaps the fused kernel for its CPU
      twin and changes no hyperparameter. \\
    Policy update & One PPO epoch per batch with parameters synchronized every
      update, so each rollout is used once and on-policy. The surrogate keeps
      the standard symmetric clip at $\varepsilon=.2$; at one epoch the
      importance ratio is $1$ when the clip is evaluated, so it binds on no
      reported update and is listed to pin the loss, not because it acts. \\
    % All twelve cells are the OAT_ZERO_GENERATE_MAX_LENGTH actually submitted,
    % read per domain from the job ledgers rather than from a launcher default:
    %   0.5B  var/artifacts/e72_frontier_source_runs.json
    %           (inherited_eval_config.generate_max_length; the 3B cohort
    %            inherits this record and overrides only the base model)
    %   1B    var/artifacts/e79_falcon1b_aligned_verified_replay_jobs.json
    %   3B    var/artifacts/e80r1_qwen3b_aligned_verified_replay_jobs.json
    % ops/train.sh defaults this to 192, so a domain that is 192 looks the same
    % whether it was chosen or inherited; the ledgers are what distinguish them.
    Response tokens & Graph and Countdown: 192; Python: 192/512/192
      at 0.5B/1B/3B; MathIR: 64/128/64; PantryPlan: 8. \\
    Evaluation seeds & Base 610100/610200/610300/610400/76299 for
      Graph/Countdown/Python/MathIR/PantryPlan; draw $d$ uses base plus $d$. \\
    \bottomrule
  \end{tabularx}
\end{table}

\begin{table}[!htbp]
\centering
\caption{\textbf{Pinned base-model revisions.} Every run in this paper starts
from these exact commits; a different revision of the same model name is a
different initialization and is not admissible as a matched arm, whatever
name the checkpoint carries.}
\label{tab:pinned-revisions}
\small
\begin{tabular}{@{}ll@{}}
\toprule
\textbf{Base model} & \textbf{Pinned revision} \\
\midrule
Qwen2.5-0.5B-Instruct & {\scriptsize\texttt{7ae557604adf67be50417f59c2c2f167def9a775}} \\
Falcon3-1B-Instruct & {\scriptsize\texttt{28ba2251970a01dd1edc7ba7dad2eb71216ccfdf}} \\
Qwen2.5-3B-Instruct & {\scriptsize\texttt{aa8e72537993ba99e69dfaafa59ed015b17504d1}} \\
\bottomrule
\end{tabular}
\end{table}

Responses are independent within each evaluation draw. For each training
seed, metrics average over all 128 prompts and four draws; raw responses
and draw-level records are retained. The co-primary terminal outcomes
registered in App.~\ref{app:registered-endpoints} are realized here as
\texttt{pass@8} and raw \texttt{distinct@8}; \texttt{mean@8} is diagnostic.
Evaluation never supplies training data, bank entries, checkpoint selection,
or a stopping rule. Complete paired five-seed blocks receive nominal
pointwise 95\% Student-$t$ intervals; partial primary intersections report
exact counts and descriptive means without intervals. Intervals condition
on the frozen prompts and evaluation requests. Domain averages are
explicitly labeled, and no missing terminal result is replaced by a
shallower checkpoint. Source admission and curve gaps follow
App.~\ref{app:source-integrity}; comparison-specific intersections and
the replay-weighting bootstrap are specified with their results.

% Appendix B ends with four heavy floats (the run contract, the pinned
% revisions, the frontier grid and the families scatter). Left to drift they
% placed inside Appendix C, which split C's opening paragraph across four
% pages with B's tables sitting in the middle of it. Flush them before C.
\clearpage

\section{Exact Domain Prompts}
\label{app:prompts}

Each domain uses one system message and one user message, with no few-shot
examples. The same template renders training and evaluation for every
method. What it does not do is impose one model's chat surface on another.
The listings below are the Qwen rendering, and a Falcon3 run trains and is
evaluated on a paired template that carries the identical system and user
text onto Falcon's own \texttt{<|system|>}/\texttt{<|user|>} markers; the two
are validated against one shared role contract, so the surface differs and
the task does not. The frozen base-model grid goes further and renders every
scale through its own tokenizer's chat template, refusing to run a model
whose native template is absent, which is how SmolLM2 and OLMo-2 are measured
(App.~\ref{app:base-levels-all-scales}). No model is asked to read a foreign
chat format. These are the first rows of the frozen Level-1 evaluation splits,
including the exported chat-template wrapper; Levels 2 to 5 are separate
constructions on the same interface (App.~\ref{app:data-levels}). The checked
generator \path{ops/check_paper_domain_prompts.py} reproduces them from the
datasets and fails closed when the typeset appendix has drifted from them.
A two-space continuation indent marks display wrapping: rejoining such a line
to the one above it with a single space recovers the original string, and in
the PantryPlan listing a break without that indent is a real newline in the
prompt.

% BEGIN generated by ops/check_paper_domain_prompts.py

\filbreak
\subsection{Graph coloring}
\label{app:prompt-graph}

\noindent Template \texttt{qwen\_boxed}; 384 training and 128 evaluation\newline rows in \path{var/data/graph_coloring_modebench_v2}.

\begin{Verbatim}[fontsize=\scriptsize,frame=single,framesep=4pt,rulecolor=\color{black!25},samepage=true]
<|im_start|>system
Return only the final answer inside \boxed{}. Do not explain.<|im_end|>
<|im_start|>user
Color this graph with vertices 1 through 6. The undirected edges are: (1,2), (1,3),
  (3,5), (3,6), (4,6), (5,6). Assign each vertex one color from {1, 2, 3} so that no
  edge has the same color on both endpoints. A partial coloring is 2??1?2, where ?
  means uncolored. Keep every shown digit fixed. There are exactly 3 question marks.
  Return exactly one final answer and no explanation: exactly 3 digits, each one 1,
  2, or 3, for the missing positions from left to right, inside \boxed{}.<|im_end|>
<|im_start|>assistant
\end{Verbatim}

\filbreak
\subsection{Countdown}
\label{app:prompt-countdown}

\noindent Template \texttt{qwen\_boxed}; 384 training and 128 evaluation\newline rows in \path{var/data/exact_countdown_easy3_probe}.

\begin{Verbatim}[fontsize=\scriptsize,frame=single,framesep=4pt,rulecolor=\color{black!25},samepage=true]
<|im_start|>system
Return only the final answer inside \boxed{}. Do not explain.<|im_end|>
<|im_start|>user
Using the numbers [3, 6, 9], create an arithmetic expression that equals 18. Use
  each given number exactly once. You may use +, -, *, /, and parentheses. Return
  exactly one final answer and no explanation: the expression inside
  \boxed{}.<|im_end|>
<|im_start|>assistant
\end{Verbatim}

\filbreak
\subsection{Python factors}
\label{app:prompt-python}

\noindent Template \texttt{qwen\_boxed}; 384 training and 128 evaluation\newline rows in \path{var/data/python_factor_modebench_v1}.

\begin{Verbatim}[fontsize=\scriptsize,frame=single,framesep=4pt,rulecolor=\color{black!25},samepage=true]
<|im_start|>system
Return only the final answer inside \boxed{}. Do not explain.<|im_end|>
<|im_start|>user
Write one pure Python function with the exact form lambda n: EXPR. An external
  Python tool will call it once for each n in [18, 82, 91, 93]. For every call,
  return an integer d satisfying 1 < d < n and n % d == 0. Different valid return
  vectors are different solution modes. You may use integer literals, n, +, -, *,
  //, %, comparisons, Boolean operators, and conditional expressions; calls,
  imports, attributes, containers, and other names are forbidden. Return exactly the
  one-line lambda inside \boxed{} and no explanation.<|im_end|>
<|im_start|>assistant
\end{Verbatim}

\filbreak
\subsection{MathIR action menu}
\label{app:prompt-mathir}

\noindent Template \texttt{qwen\_boxed}; 384 training and 128 evaluation\newline rows in \path{var/data/mathir_action_menu_v1}.

\begin{Verbatim}[fontsize=\scriptsize,frame=single,framesep=4pt,rulecolor=\color{black!25},samepage=true]
<|im_start|>system
Return only the final answer inside \boxed{}. Do not explain.<|im_end|>
<|im_start|>user
MathIR linear-menu-v1. Solve for x by choosing an executable action sequence. Every
  selected action is applied exactly to BOTH sides of the current equation, followed
  by exact simplification. Bindings: a=2, b=-9, c=-1. Initial equation: x/a + b = c.
  Action menu: A means div(a); B means add(b); C means sub(b); D means sub(c); E
  means sub(mul(a,b)); F means mul(a). Choose 1-4 action IDs. Illegal operations,
  repeated states, and paths that do not finish with x isolated are rejected. Return
  only the IDs separated by semicolons inside the answer box; for example formatting
  only: \boxed{A;C}. Do not return x's numeric value.<|im_end|>
<|im_start|>assistant
\end{Verbatim}

\subsection{PantryPlan}
\label{app:prompt-pantry}

\noindent Template \texttt{qwen\_pantry\_support\_mask}; 384 training and 128 evaluation\newline rows in \path{var/data/pantry_plan_modebench_v2}.

\begin{Verbatim}[fontsize=\scriptsize,frame=single,framesep=4pt,rulecolor=\color{black!25},samepage=false]
<|im_start|>system
Return only six binary digits and no other text. Each digit maps to the
  corresponding Pantry row in printed order. A trusted environment will choose exact
  quantities on precisely the selected rows.<|im_end|>
<|im_start|>user
Create one feasible high fiber snack from this frozen pantry.
Nutrition values are per 100 g. Quantities are exact integer grams.

Pantry:
- almonds: available=100g, step=25g, minimum-if-used=50g; energy=584kcal,
  protein=21.5g, fiber=10.8g, sodium=0.0mg; tags=tree_nut
- banana: available=125g, step=25g, minimum-if-used=50g; energy=97.0kcal,
  protein=0.74g, fiber=4.62g, sodium=0.0mg; tags=none
- carrots: available=150g, step=25g, minimum-if-used=50g; energy=45.0kcal,
  protein=0.941g, fiber=3.1g, sodium=86.6mg; tags=none
- grape_tomatoes: available=100g, step=25g, minimum-if-used=50g; energy=27.0kcal,
  protein=0.83g, fiber=2.1g, sodium=6.0mg; tags=none
- navel_orange: available=150g, step=25g, minimum-if-used=50g; energy=47.0kcal,
  protein=0.91g, fiber=2.0g, sodium=9.0mg; tags=none
- sunflower_seeds: available=125g, step=25g, minimum-if-used=50g; energy=571kcal,
  protein=18.9g, fiber=7.22g, sodium=0.0mg; tags=none

Use 2 to 4 ingredients.
Forbidden tags: none
Targets:
- mass_g: minimum 125, maximum 200
- energy_kcal: minimum 239.4, maximum 432.25
- protein_g: minimum 6.423
- fiber_g: minimum 3.478
- sodium_mg: maximum 37.6

Choose the ingredient support as a six-bit mask in the exact Pantry row order. Bit 1
  includes that row and bit 0 excludes it.<|im_end|>
<|im_start|>assistant
\end{Verbatim}

% END generated by ops/check_paper_domain_prompts.py

% Shared exact replay update; empirical comparisons are in results.tex.

\section{Replay Algorithm}
\label{app:algorithm}

Table~\ref{tab:notation} lists the symbols used here. This section states the
objective as implemented; App.~\ref{app:replay-metric-alignment} reads the same
objective against the evaluation metric, which is why that analysis sits with
the measurement rather than here.
For each prompt $x$, $\mathcal B_x$ stores at most $M$ verified policy
exemplars, one per canonical key; the serialized cursor $\sigma$ selects
banks for replay. The same update supports either fresh objective.


\subsection{Fresh task advantages}
\label{app:binary-maxrl}

For a group of $G$ fresh responses with verified binary rewards $R_i$ and
$R=\sum_i R_i$, Dr.GRPO uses $A_i=R_i-R/G$, while MaxRL uses
\[
 A_i^{\mathrm{MaxRL}}=
 \begin{cases}
 G R_i/R-1, & R>0,\\
 0, & R=0.
 \end{cases}
\]
The centered MaxRL estimator \citep{tajwar2026maxrl} gives failures advantage
$-1$ in mixed groups and weights each success more strongly when successes
are scarce. All-correct and all-failure groups have zero task advantage.
Lemma~\ref{lem:maxrl-mean} gives its harmonic pass@\(k\) mean objective.
These advantages apply only to fresh responses; bank exemplars enter the
separate verified-likelihood loss.

\subsection{State and optimizer update}
\label{app:algorithm-update}

\begin{algorithm}[H]
  \caption{One verified-replay optimizer update: $o=\mathrm{Dr.GRPO}$
    gives Re:Dr, $o=\mathrm{MaxRL}$ gives Re:Max}
  \label{alg:verified-replay-update}
  \small
  \begin{algorithmic}[1]
    \REQUIRE Policy $\pi_\theta$; prompt $x$ and verifier specification
      $s_x$; group size $G$; fresh objective
      $o\in\{\mathrm{MaxRL},\mathrm{Dr.GRPO}\}$; replay weight
      $\lambda_{\mathrm{replay}}$; states $(\mathcal B,\sigma)$
    \ENSURE Atomically updated $(\theta',\mathcal B',\sigma')$
    \STATE Draw the on-policy group
      $\mathcal G=\{y_i\sim\pi_\theta(\cdot\mid x)\}_{i=1}^G$
    \FOR{$i=1,\ldots,G$}
      \STATE Execute $\widetilde c_i\gets V(y_i,s_x)$ and observe task
        reward $R_i$
      \STATE Set $c_i\gets\widetilde c_i$ iff
        $R_i>0\wedge\widetilde c_i\ne\bot$; otherwise set $c_i\gets\bot$
    \ENDFOR
    \STATE Compute fresh task advantages $A_i^{o}$ from $\{R_j\}_{j=1}^G$;
      for $o=\mathrm{MaxRL}$ use the formula in App.~\ref{app:binary-maxrl}
    \STATE After the full group is validated, deterministically retain one
      policy exemplar for each new $c_i\ne\bot$ in $\mathcal B_x$, up to
      capacity $M$
    \STATE $(z,\sigma')\gets\textsc{NextBank}
      (\{\mathcal B_u\}_u,\sigma)$ by a fixed round-robin cursor over
      prompts; the replayed prompt $z$ need not be the prompt $x$ that
      supplied the fresh group
    \IF{a replay bank $\mathcal B_z$ exists}
      \STATE For each exemplar $e_c\in\mathcal B_z$, compute
        $\ell_c=|e_c|^{-1}\log\pi_\theta(e_c\mid z)$
      \STATE $L_{\mathrm{replay}}\gets-|\mathcal B_z|^{-1}
        \sum_{c\in\mathcal B_z}\ell_c$
    \ELSE
      \STATE $L_{\mathrm{replay}}\gets0$
    \ENDIF
    \STATE $w_{\mathrm{rep}}\gets(G-1)/G^2$ for one replay
      pseudo-group per $G$ on-policy rows
    \STATE $\theta'\gets\textsc{Update}
      (\theta,L_{\mathrm{PPO}}(A^{o})
      +w_{\mathrm{rep}}\lambda_{\mathrm{replay}}L_{\mathrm{replay}})$
    \STATE Checkpoint $(\theta',\mathcal B',\sigma')$ as one atomic state
\end{algorithmic}
\end{algorithm}

\subsection{Admission, capacity, and deterministic scheduling}
\label{app:algorithm-bank}

A bank is keyed by the SHA-256 digest of the exact prompt-token sequence.
Only loss-active rows with positive reward and a nonempty validator key
qualify. A row is loss-active when it carries weight in the policy-gradient
loss, which the framework withholds from a response that reaches the
generation cap without emitting an end-of-sequence token. That exclusion is
disabled in every run reported here --- MaxEnt length control refuses to run
with it on, since it needs the truncated rows --- so the qualifier admits
every rollout row and excludes none; it is stated because a deployment that
enables it would silently change what the bank can hold. The full group is validated against the pre-group bank before
any new key is committed. For each new key, the lexicographically smallest
normalized response-token tuple is retained; new keys are admitted in
sorted-key order. Admission is therefore invariant to rollout-row order.
The bank keeps one exemplar per key, up to $M=16$, with no eviction or
replacement. Later keys may enter diagnostics but not replay support when
capacity is full. Prompt-hash collisions or changed tokens under an
existing hash are hard errors.

At each update, the scheduler sorts nonempty prompt-bank hashes and selects
one by a fixed round-robin cursor; singleton banks are eligible. The cursor
is independent of reward, evaluation, desired support, or entropy, and is
serialized with the bank. Stored prompt and response tokens are
materialized unchanged into a right-padded teacher-forcing batch. Prompt
and padding positions are masked. Each exemplar contributes its mean
response-token log probability, averaged uniformly across keys in the
selected bank before applying $\lambda_{\mathrm{replay}}$ and the exact
$(G-1)/G^2$ rollout scaling in
Algorithm~\ref{alg:verified-replay-update}. That constant is the weight one
lucky success already carries in a centered fresh group, so a replay
pseudo-group enters at the same size as the event it stands in for: the proof
of Lemma~\ref{lem:grpo-mean} gives a lone correct row the centered reward
$(G-1)/G$, and the outer $1/G$ that averages the group supplies the second
factor. Neither exhaustive support nor
evaluation outputs enter this loss.

\subsection{Matched controls and recoverable state}
\label{app:algorithm-control}

Controls share validation, bank updates, scheduling, materialization, and
the score pass, with the replay score-gradient vector replaced by exact
zeros. Treatment activates only that derivative. Policies can subsequently
differ in samples, bank contents, selected banks, exemplar lengths, and
realized computation. Logs retain bank/key counts, replayed token counts,
raw and applied losses, score-gradient summaries, cursor phase, and the
exact-zero switch. Model and optimizer state, bank counts and exemplar
tokens, replay and data cursors, and request-stream state are checkpointed
atomically; the loader validates schema, capacity, scheduler and estimator
constants, token types, prompt hashes, key/exemplar consistency, and cursor
bounds before accepting a resume.

\subsection{What the bank needs from a verifier}
\label{app:algorithm-identity}

Replay reads one thing from the task that ordinary RLVR does not: an
equivalence on outputs. Everything else in Algorithm~\ref{alg:verified-replay-update}
is objective-agnostic. That equivalence has to be decidable while the rollout
is in hand, cheap enough to run on every row of every group, a function --- the
same output always yields the same key --- and consistent with the reward, so
that ``still correct'' and ``a different answer'' cannot disagree. \mb{}
supplies it from the executable check that already grants credit, which is why
no judge or embedding threshold stands between the data and the claim.

The widely used mathematical reasoning sets do not supply it. A grader for
MATH, GSM8K or AIME returns correctness on a final answer, and the final
answer is unique, so an identity keyed on what the grader returns gives every
correct response to a prompt the same key. A bank under that identity holds
one exemplar per prompt, the replay term reduces to rehearsing the single
success the policy already found, and the mechanism this paper measures has
nothing to rebalance. Getting more than one mode out of such a task means
keying identity on the derivation instead of the answer, and MathIR is the
cautionary case: it keys on the executed state trajectory, its certified
alternatives are mostly longer routes, and across \MIRdraws{} verified draws
no model in this paper produced one. Our own evidence therefore places open-ended derivations at the hard end
of this method's range, not the easy end.

A model judging whether a candidate derivation is new relative to the bank
would restore an equivalence where execution cannot supply one, and it is the
natural adaptation. It is also a trade rather than a free extension, and the
terms are worth stating. It gives up the property that identity and reward
come from one check, so the breadth being optimized and the breadth being
reported would no longer be the same object. It gives up determinism unless
the judge is pinned to a fixed model, prompt and temperature and frozen for
the run; admission here is invariant to rollout-row order
(App.~\ref{app:algorithm-bank}), and a judge that answers the same pair two
ways forfeits that. It also puts the identity inside the optimization: replay
reinforces whatever the judge calls new, so the judge's blind spots become
a target, which is the reason the comparators in this paper that require an
in-loop judge are cited rather than run.

Those costs are measurable rather than hypothetical, and \mb{} is the
instrument for measuring them. A judge can be scored against ground-truth
identity on the domains where an executable key exists, before it is trusted
on domains where none does, and the disagreement rate on that calibration set
bounds what it would contribute as training signal. Reporting \pmd{} under the
executable key while training under the judge key would keep measurement
independent of the signal on exactly the cells that permit both. This paper
takes neither step; the design is set out because the conditions it would have
to meet are specific, and the MathIR result is a reason to expect a smaller
gain where a mode is a derivation rather than an answer.

\subsection{What the bank does not rank}
\label{app:algorithm-quality}

Uniform rehearsal over discovered keys says nothing about whether a discovered
key deserved to be kept. A verified output can be long-winded, or can reach a
correct answer by reasoning that would not survive inspection, and the
verifier will still grant it credit: a canonical key describes an executed
output, not the thinking behind it. Re:Max preserves modes; it does not
appraise them.

Two parts of the update already blunt the specific worry that verbosity is
what gets reinforced. An exemplar is scored by its \emph{mean} response-token
log probability, $\ell_c=|e_c|^{-1}\log\pi_\theta(e_c\mid z)$, so its
contribution does not grow with its length, and the bank holds one exemplar
per key, so a verbose mode cannot occupy more of a bank than a terse one. The
idealized account is sharper still: under the mean-token exemplar loss the
effective per-key weights go as $L_b^{-1}$, so the limit tilts toward shorter
exemplars rather than longer ones, and is uniform only when lengths are equal.
Length inflation is therefore not a direction this objective rewards.

What remains true is that the bank applies no quality order at all, and that
is a choice rather than an oversight. Equal weighting over verified keys is
the ingredient these experiments isolate; mixing a brevity or quality prior
into the same term would make a measured gain ambiguous between keeping more
modes and preferring better ones, which is the one distinction the design
exists to make.

The seam for anyone who wants both is admission. A row already enters a bank
only when it is loss-active, carries positive reward and parses to a canonical
key (App.~\ref{app:algorithm-bank}); a token budget, a process-reward score or
a trace-quality classifier is a fourth conjunct on that rule, and the
scheduler, the pseudo-group scaling and the matched control are untouched by
it. Exemplar selection within a key and eviction at capacity are the other two
places such a signal naturally attaches. Nothing in this paper's claim
competes with that literature: the claim is that when a verified output is a
\emph{new} mode, keeping it is worth doing, and worth tracking beside whatever
quality signal is already being optimized.

The composition is not free, and this benchmark shows where it binds. Where a
domain's alternative modes are intrinsically longer than its preferred one, a
brevity filter removes exactly the modes that make its support non-trivial:
MathIR's alternatives are detours through extra states, and Python's rarer
return vectors
need longer lambdas than the modal one. Breadth and terseness therefore pull
against each other hardest in the domains where breadth is already hardest to
hold, which is a reason to read \pmd{} after such a filter rather than before
it.

\section{Fixed Semantic MaxEnt Comparator}
\label{app:semantic-maxent}

Fixed Semantic MaxEnt is a fresh-rollout comparator, not a replay method: it
adds no bank, reuses no verified exemplar, and touches none of the state
App.~\ref{app:algorithm} specifies. It is set out here because it is the one
alternative objective this paper implements rather than cites, and because
App.~\ref{app:ucpo-progress} reads its measured effects against that
specification. It is entropy-inspired: for validator-positive canonical
outcomes, an empirical predictor combines successful policy history and
successful leave-one-out peers. Scoring row $i$ of a group of $G$ on prompt
$x$, write $n_x(a)$ for the historical count of key $a$ at that prompt,
$m_{-i}(a)$ for its count among that row's $G-1$ peers,
$N_x=\sum_a n_x(a)$, and $\mathcal K_i$ for the union of the historical and
peer-observed keys. Unit pseudocounts and one reserved unseen bucket give
\begin{equation}
  \hat p_i(a)=\frac{n_x(a)+m_{-i}(a)+1}{N_x+G-1+|\mathcal K_i|+1},\qquad
  \hat p_i(\textsc{unseen})=\frac{1}{N_x+G-1+|\mathcal K_i|+1}.
  \label{eq:open-set-predictor}
\end{equation}

\paragraph{What the counts accumulate.}
The estimator keeps one count table per prompt, addressed by the SHA-256
digest of the unpadded prompt-token sequence, and keeps nothing global: two
prompts never share a count, and a key's count at one prompt says nothing
about how often that key was produced at another. A row reaches the table
only when it is loss-active, carries positive task reward, and parses to a
canonical key --- the three conditions the replay bank admits on
(App.~\ref{app:algorithm-bank}) --- so invalid, unsuccessful and unparseable
rows leave it unchanged. The table is strictly cumulative over every update
the run has taken on that prompt: counts only rise, and there is no window,
decay, eviction or reset. It is serialized into the checkpoint's client state
beside the optimizer, so a resumed run continues the history rather than
restarting it. Within a group the order is score first, then commit: every
candidate is scored against the history that earlier groups left plus its own
$G-1$ peers, and the group is added only once all of its candidates have been
scored, so no row ever contributes to the count that scores it.
The added detached advantage on an eligible successful row is
\begin{equation}
  A^{\rm sem}_{i,t}=\eta\,
  \frac{\min\{-\log\hat p_i(a_i),\kappa\}
    -\E_{\hat p_i}[\min\{-\log\hat p_i(a),\kappa\}]}{\kappa},
  \qquad \eta=.10,\ \kappa=5.
  \label{eq:legacy-semantic-advantage}
\end{equation}
The score does not depend on the token index, so row $i$ carries the same
$A^{\rm sem}_{i,t}$ at every $t$. Invalid and unsuccessful rows receive zero,
and the centered score is bounded by $\eta$ in magnitude. It can favor an underrepresented verified outcome
present in a fresh rollout, but supplies no direct reinforcement to an
absent remembered exemplar. Its heterogeneous empirical effects therefore
remain a comparison with exploration rather than a retention guarantee.

For the true conditional distribution $q_\theta=p_\theta(\cdot\mid\mathcal C)$,
\begin{equation}
 \nabla_\theta H(q_\theta)
 =\E_{a\sim q_\theta}[(-\log q_\theta(a))\nabla_\theta\log q_\theta(a)].
 \label{eq:entropy-score-function}
\end{equation}
This motivates the direction of redistribution, not an unbiased-gradient
claim for the empirical estimator. Hold a detached predictor score vector
$s$ and baseline $b$ fixed independently of the scored action, for example
conditional on its leave-one-out peers. Averaging successful rows under the
unconditional policy then gives
\[
 \E_{Z\sim p}[\1\{Z\in\mathcal C\}(s_Z-b)\nabla_\theta\log p_Z]
 =P\E_{c\sim q}[s_c\nabla_\theta\log q_c]
   +(\E_q[s]-b)\nabla_\theta P.
\]
A historical baseline need not equal the current conditional mean, so this
bonus can change correctness as well as diversity. Even with the exact
unclipped score $s_c=-\log q_c$ and baseline $b=H(q)$, averaging over
unconditional rollout rows gives $P\nabla_\theta H(q)$.
For a single correct key, conditional entropy is zero, yet the implemented
unseen bucket can give that key negative semantic advantage. The comparator
thus carries neither an exact entropy-gradient interpretation nor an
individual-mode preservation guarantee.

% Every part divider opens a fresh page, this one included. It used to run on
% from the comparator section to save the white space below that section's end;
% uniformity across the seven dividers is worth the part page.
\clearpage
\phantomsection
\begin{center}\large\textbf{Part III. Training results}\end{center}
\addcontentsline{toc}{part}{Part III. Training results}

\section{Empirical Results and Comparisons}
\label{app:supporting}

\begin{figure}[!htbp]
\centering
\includegraphics[width=\linewidth]{figures/replay_factorial_effects.pdf}
\caption{\textbf{Canonical replay complements both fresh objectives.}
Each point compares replay with its matched no-replay control at pass eight,
within domain and scale, on the paper's two axes: correctness in probability
points, and \pmd{} on the successes alone. The rows do not share a seed
population. For \texttt{pass@8} all 75 MaxRL pairs are complete and Falcon
Countdown uses four Dr.GRPO pairs. For \pmd{} a seed counts only where both arms
define it: on the Dr.GRPO side that leaves Falcon Countdown at four, Falcon
MathIR at one, Qwen2.5-3B Python at two, and Falcon Python blank rather than
zero. Nominal intervals accompany complete blocks only.}
\label{fig:replay-factorial-effects}
\end{figure}

\subsection{Terminal comparisons across scales}
\label{app:scale-matrix}
\label{app:scale-snapshot}
\label{app:maxrl-scale-extensions}
\label{app:current-campaign-results}

Within each model--domain--seed comparison, methods share initialization,
data identities, the eight-pass horizon, evaluation requests, verifier, and
replay scheduling code. The replay control runs the same bank and scoring
path with zero replay derivative. Resulting policies can build different
banks. The admissible intersections are defined in
App.~\ref{app:source-integrity}.

\crossscaletablematerial

\paragraph{Dr.GRPO retention contrasts.}
\label{app:retention-numerical-detail}

% Floats rather than [H] so it can rise to the page its discussion starts
% on instead of forcing a break and leaving the preceding page part empty.
\begin{table}[!tbp]
  \centering
  \caption{\textbf{Replay improves correctness and verified breadth at every
  scale.} Re:Dr minus matched Dr.GRPO at the terminal checkpoint, per domain
  and scale behind the cross-domain averages of
  Fig.~\ref{fig:maxrl-factorial}. Columns, the equally weighted Avg, the
  superscripted paired seed counts and the em dashes for a block with too few
  verified pairs on one side all read as they do in
  Table~\ref{tab:direct-comparator-matrix}. The two Avg columns are equally
  weighted within seed over a basis fixed per row: for \texttt{pass@8} that is
  all five domains, on the seeds all five share; for \pmd{} it is the domains
  every seed of that row defines, which drops Python at both larger scales and
  Countdown and MathIR at Falcon3-1B, so the \pmd{} average is never taken
  over a domain set that changes from seed to seed.}
  \label{tab:cross-scale-terminal-effects}
  \scriptsize
  % Thirteen columns with the domains spelled out: the two widest headers,
  % Countdown and PantryPlan, set the measure rather than the digits do, so the
  % padding comes down to keep the names at full size inside \linewidth.
  \setlength{\tabcolsep}{1.4pt}
  \renewcommand{\arraystretch}{0.92}
  \begin{tabular}{@{}l*{12}{r}@{}}
    \toprule
    & \multicolumn{6}{c}{\textbf{$\Delta$ \texttt{pass@8}}}
      & \multicolumn{6}{c}{\textbf{$\Delta$ \pmd{}}} \\
    \cmidrule(lr){2-7}\cmidrule(l){8-13}
    & Graph & Countdown & Python & MathIR & PantryPlan & Avg
      & Graph & Countdown & Python & MathIR & PantryPlan & Avg \\
    \midrule
    \input{results/retention_matrix_panel_a_table_body.tex}
  \end{tabular}
\end{table}

Table~\ref{tab:cross-scale-terminal-effects} gives the same Re:Dr contrast per
domain and scale that Table~\ref{tab:core-terminal-endpoints} gives as levels,
and reading the two together is what separates a gain from a ceiling: a large
difference on a domain whose control already sits high is a different result
from the same difference on one that starts at zero. Both co-primary block
means are positive in all 14 complete five-seed comparisons, and the registered
\distk{} endpoint increases in all 74 admissible pairs; its per-cell values are reported in the machine-readable
record rather than here, because it rises with correctness and cannot separate
the two. The paired \pmd{} effects can. At Qwen2.5-0.5B, Countdown
\texttt{pass@8} rises from $.480$ to $.672$ while \pmd{} rises $+.480$ (paired
95\% interval $[+.447,+.513]$). Python correctness rises from $.172$ to $.528$,
with a \pmd{} effect of only $+.082$ ($[-.146,+.310]$); MathIR rises from
$.512$ to $.796$ with \pmd{} $+.017$ ($[+.003,+.031]$). These intervals are
nominal, unadjusted Student-$t$ estimates over five paired training seeds.
Correctness and breadth separate cleanly here, so the same Countdown
comparison is also read on the
success-conditional axis, where Dr.GRPO ends at \pmd{} $.007$ and
Re:Dr at $.487$: the replayed policy is not merely succeeding more
often, its successes are spread across modes where the control's are not. Countdown keys identify computation routes, but
these held-out endpoints do not identify which previously discovered
routes survived training.


Fig.~\ref{fig:maxrl-factorial-scale-extensions} resolves the domain-specific
results behind the main cross-domain display. All five 3B MaxRL domain comparisons are complete:
both arms have seeds 70--74 at the registered terminal endpoint (50/50 runs).
The Dr.GRPO pair uses five seeds except Falcon Countdown, which uses four. These populations
remain separate: a factorial interaction requires the common four-arm
intersection, not differences of aggregates with different seeds.

% Not [H]: this plate is a full text-block tall, so forcing it here left the
% rest of the preceding page blank whenever it did not fit. Letting it float
% fills that page and keeps the plate near the text that cites it by number.
\begin{figure}[!htbp]
  \centering
  \includegraphics[width=.98\linewidth]
    {figures/e118_scale_extensions_appendix.pdf}
  \caption{\textbf{Per-domain results behind the cross-domain averages in
  Fig.~\ref{fig:maxrl-factorial}.} Rows show Qwen2.5-0.5B,
  Falcon3-1B, and Qwen2.5-3B; columns show correctness and \pmd{}, the same two
  axes the main results use. On \texttt{pass@8} every MaxRL--Re:Max pair is
  complete at five terminal seeds across all three models, and Dr.GRPO--Re:Dr
  is complete except for Falcon Countdown at four. Solid connectors identify
  the complete comparisons and dashed ones carry their count. The two columns
  do not share a seed population: \pmd{} is undefined for a seed whose
  policy succeeds on too few prompts, so the right column carries its own
  counts, and a row reading \emph{below the support bar} is one where no seed
  on one side of either pair defines \pmd{} --- breadth unmeasured there, not
  measured as zero. This panel applies the 30-prompt bar of
  App.~\ref{app:metric-sampling}; Table~\ref{tab:cross-scale-terminal-effects}
  pairs seeds at the lower 20-prompt bar and so
  reports its own, larger counts for the same cells. Domain tints match
  Fig.~\ref{fig:modebench-examples}. Replay arrows, trained-arm markers, and
  untrained references use the same encoding as
  Fig.~\ref{fig:maxrl-factorial}.}
  \label{fig:maxrl-factorial-scale-extensions}
\end{figure}


Write $P=\texttt{pass@8}$ (Table~\ref{tab:notation}). The endpoint tables pair
it with \pmd{}, the paper's breadth axis, and report exact paired terminal
effects on both. The registered \distk{} endpoint and the derived count
$B=\distk{}-P$ of verified modes beyond the first correct outcome are not
tabulated here: both rise with correctness, so neither distinguishes succeeding
more often from succeeding in more ways, and both are reported in full in the
adjacent machine-readable records. The one exception is
Table~\ref{tab:level2-factorial-contrasts}, where an interaction needs both of
its differences on the same seed and too few Level-2 seeds define \pmd{} on
all four arms to carry one. The two axes also carry different seed
populations, listed separately: a seed enters the \pmd{} column only where both
arms define it. A dash denotes no observed pair; partial means retain $n/5$
without an interval. Complete-block intervals are nominal 95\% estimates,
without multiple-comparison or repeated-look adjustment. All underlying arm
means and seed effects remain in the adjacent machine-readable records.

At Qwen2.5-0.5B, Re:Max raises \pmd{} by $+.439$, $+.489$, and $+.317$ on
Graph, Countdown, and PantryPlan, respectively; each five-seed interval excludes
zero. Falcon Graph instead has a negative mean breadth effect ($-.009$, interval
$[-.034,+.016]$), and its correctness effect is inconclusive ($+.014$,
$[-.016,+.044]$). Positive cross-domain effects therefore do not imply a breadth
gain in every domain. All five Qwen2.5-3B blocks are complete, with \pmd{}
effects of $+.076$, $+.077$, $+.058$, $+.006$, and $+.224$ for Graph, Countdown,
Python, MathIR, and PantryPlan; every nominal interval excludes zero, although
MathIR's gain is small enough to be of no practical size; that domain's
support leaves almost no reachable alternative to gain. Correctness intervals
exclude zero for Countdown, MathIR, and PantryPlan; Graph and Python remain
inconclusive. Read across the cohort, Re:Max ends above MaxRL in
\PMDaboveReMax{} of the \PMDblocksReMax{} comparisons that clear the paired
support bar, and Re:Dr above Dr.GRPO in \PMDaboveReDr{} of
\PMDblocksReDr{} (App.~\ref{app:registered-endpoints}).

\paragraph{Cross-domain factorial summaries.}
\label{app:factorial-numerical-detail}
Averaging five domains equally within each paired seed, Re:Max minus MaxRL
gives $+.307$ on \texttt{pass@8} at Qwen2.5-0.5B (unadjusted 95\% interval
$[+.208,+.406]$) and $+.133$ at Falcon3-1B ($[+.043,+.223]$). Both use five
paired seeds. The Dr.GRPO replay effects on correctness are $+.337$ at
Qwen2.5-0.5B, $+.442$ at Falcon3-1B, and $+.279$ at Qwen2.5-3B. Breadth is not averaged across domains here: \pmd{} admits a seed only where
both arms define it, so a mean over all five would rest on a population that
changes with the domain. Where
Table~\ref{tab:cross-scale-terminal-effects} does report a \pmd{} average it
fixes the basis first, to the domains every seed of that row defines. The
per-domain \pmd{} effects are drawn in
Fig.~\ref{fig:maxrl-factorial-scale-extensions}. The Falcon Dr.GRPO average uses four common admissible seeds
without an interval; the other two use five. At Qwen2.5-3B, MaxRL/Re:Max arm
means are $.547/.701$ on correctness and $.744/1.107$ on raw
\texttt{distinct@8}. Equal domain averages within each of seeds 70--74 give
replay effects of $+.154$ ($[+.086,+.222]$) and $+.363$ ($[+.285,+.441]$),
respectively. The common four-arm intersection contains fourteen complete
five-seed blocks and Falcon Countdown at four seeds.


\subsection{Terminal comparison on the harder level}
\label{app:level2-interim}

\begin{figure}[!htbp]
\centering
\includegraphics[width=\linewidth]{figures/replay_level2_effects.pdf}
\caption{\textbf{Replay effects within Level 2.}
Replay and control are compared on the same prompts and training seeds, with
separate Dr.GRPO and MaxRL effects, on correctness and on \pmd{} over the
successes alone. All five domains have complete five-seed factorials on
\texttt{pass@8}, including the twenty PantryPlan endpoints. The breadth panel
does not share that cohort: \pmd{} counts a seed only where both arms define it,
which at Level 2 leaves Dr.GRPO at three seeds on Graph, four on MathIR, three
on PantryPlan and none at all on Countdown, and MaxRL at one on Countdown and
four on PantryPlan. A domain with no drawn mark is one where the comparison
cannot be made rather than one where it came out at zero. MathIR's flat
breadth is the domain rather than the method: every Level-2 problem it poses
admits a single shortest derivation, so an arm that solves it the short way
has no second mode to hold. Construction and admission are described in Section~\ref{sec:levels-design}; terminal
absolute values appear in Fig.~\ref{fig:level2-admission}. Differences
between levels combine task, population, and protocol changes, so they do
not identify a causal effect of difficulty alone.}
\label{fig:replay-level2-effects}
\end{figure}

Fig.~\ref{fig:level2-admission} spans three levels, so its across-domain
\texttt{pass@8} mean is carried by the four domains --- Graph coloring,
Countdown, Python factors and PantryPlan --- that hold all five seeds at step
3,072 in all four methods at every level. Its \pmd{} mean is carried by the
three --- Graph coloring, MathIR and PantryPlan --- that additionally clear the
support bar in all four methods at every level, and MathIR enters that mean on
the seeds its arms support, which the figure's record gives per arm. Both bases are the same
for every arm and level, so no arm rests on an easier subset than the one it is
being read against. Means average domains
equally within each seed, then the seed means. All 100 Level-2 endpoints are
admitted; each within-domain contrast uses seeds 43--47, including all four
arms. The levels have separate held-out prompts, with population differences in
App.~\ref{app:data-levels}. Each level is a separate construction rather than a
point on a scale, so the figure plots what replay adds within a level rather
than four absolute positions per level. Every bar is a difference of two means
reported above, both sides sharing the level, the prompts, the seeds and the
domain basis; the differences are descriptive and carry no interval.

The untrained checkpoint each level's training starts from is recorded beside
the gains rather than drawn with them: it is a position, not a gain. The figure's machine-readable record carries it per level, averaged over every
admitted step-0 cell on the domains the bars use, with its support. Levels 1 and 3 carry one on both
axes. At Level 2 only PantryPlan clears the support bar before training, so no
untrained breadth is reported there; a value resting on one of the three
matched domains would not be comparable to the rest, and an unmeasured breadth
is a different statement from a breadth of zero.

Level-2 Graph provides the clearest replay effects on both axes under both
objectives: \pmd{} rises $+.212$ ($[+.170,+.255]$, five seeds) under Re:Max and
$+.314$ ($[+.170,+.458]$, three) under Re:Dr. Python's breadth effects are
positive in mean with intervals spanning zero ($+.088$ and $+.181$), and
MathIR's are indistinguishable from zero on both arms. Pantry's complete
five-seed Dr.GRPO replay contrast is $+.305$ on $P$; its MaxRL contrast is
$+.237$. The harder construction also costs breadth coverage rather than
correctness coverage: \pmd{} admits a seed only where both arms define it, and
at Level 2 that leaves Pantry at four seeds under Re:Max and three under Re:Dr,
Graph at three under Re:Dr, and Countdown with one seed under Re:Max and none
under Re:Dr. Those cells are gaps in what can be measured, not measurements of
no effect.

For each common seed and endpoint $Y$, the replay interaction is
\[
 I=(Y^{\rm Re:Max}-Y^{\rm MaxRL})
    -(Y^{\rm Re:Dr}-Y^{\rm Dr.GRPO}).
\]
Replay can add value under both fresh objectives without a positive
interaction. On Level-2 Graph, MaxRL alone raises $P$ relative to Dr.GRPO
by $+.233$ ($[+.140,+.325]$), while the replay interaction on $P$ is
$-.256$ ($[-.329,-.183]$). This is an additional replay benefit after
stronger fresh optimization, without a superadditive interaction.

\begin{table}[H]
  \centering
  \caption{\textbf{Fresh-objective effects and replay interactions on Level 2.}
  All three contrasts within a domain use the same four-arm seed
  intersection. The interaction is
  $(\mathrm{Re:Max}-\mathrm{MaxRL})-
  (\mathrm{Re:Dr}-\mathrm{Dr.GRPO})$.
  Each entry gives its mean and, only for a complete five-seed block, a
  paired unadjusted 95\% Student-$t$ interval. This is the one table that keeps
  the \distk{} axis: an interaction needs both differences on the same seed, and
  at Level 2 the seeds where all four arms define \pmd{} number five only on
  Python, four on MathIR, three on Graph, two on PantryPlan and none on
  Countdown --- too thin to carry the estimand.}
  \label{tab:level2-factorial-contrasts}
  \scriptsize
  \setlength{\tabcolsep}{3pt}
  \resizebox{\linewidth}{!}{%
  \begin{tabular}{@{}lclrrr@{}}
    \toprule
    \textbf{Domain} & $n$ & \textbf{Contrast} &
    $\Delta P$ & $\Delta D$ & $\Delta B$ \\
    \midrule
    \input{results/level2_factorial_contrasts_20260912_table_body.tex}
  \end{tabular}}
\end{table}


\paragraph{Level-2 Graph endpoint detail.}
\label{app:levels-numerical-detail}
The development admission comparison lowers frozen-model \texttt{pass@8} from $.438$
to $.148$. At pass 8, Dr.GRPO/MaxRL reach $.202/.434$, versus $.763/.739$
with their replay variants. The paired correctness gains are $+.561$ and
$+.305$ (unadjusted 95\% Student-$t$ intervals $[+.445,+.677]$ and
$[+.250,+.360]$, five seeds each). The weakest replay run ($.711$) exceeds
the strongest non-replay run ($.477$). Replay adds $1.019/.668$ raw
\texttt{distinct@8}, yielding $1.243/1.214$ under the two fresh objectives.
Correctness rises steeply here, so that raw gain is partly a correctness gain;
\pmd{} separates the two. On the same block Dr.GRPO ends at $.111$ and MaxRL at
$.112$, against $.346$ and $.325$ for their replay variants --- roughly a
threefold difference in the chance that two verified colorings differ, at
matched conditioning on being correct.
These observations show improved correctness and sampled support on harder
Graph; they do not establish that broader support causes the correctness
gain. Cross-level prompts are distinct, and the Level-2 aggregate
covers Graph, Countdown, Python and MathIR.

\subsection{Accuracy and verified modes throughout training}
\label{app:factorial-training-curves}

The following curves include every primary method and model scale used
in the main paper. Each objective pair uses its domain-specific terminal
seed intersection as one fixed cohort throughout training. A checkpoint
is plotted only when all members have complete, valid evaluations for both
methods. Missing or conflicting checkpoints leave gaps; seed counts never
shrink to extend a mean. Four fixed-seed draws of eight samples on the same
128 held-out prompts define each evaluation. Accuracy and mode curves use
identical seeds and checkpoints, without smoothing or outcome-based selection.

\begin{figure}[H]
  \centering
  \includegraphics[width=\linewidth]{figures/factorial_training_curves_pass8.pdf}
  \caption{\textbf{Accuracy throughout training at all three model scales.}
  Rows are Qwen2.5-0.5B, Falcon3-1B, and Qwen2.5-3B; columns are the five
  domains. Dr.GRPO, Re:Dr, MaxRL, and Re:Max use the same method
  colors as Fig.~\ref{fig:maxrl-factorial}. Lines are fixed-cohort means of
  \texttt{pass@8}; bands are seed ranges, not confidence intervals. Counts
  identify each objective's terminal paired seeds. Qwen2.5-3B MaxRL uses
  all five seeds in every domain; Falcon Countdown Dr.GRPO/replay uses four.
  Conflicted checkpoints are admitted by attempt provenance wherever exactly
  one attempt's trajectory continues past them
  (App.~\ref{app:registered-provenance}), and a checkpoint the admission leaves
  open is still drawn from the cohort seeds that report it, ringed to show that
  it rests on part of the cohort. One checkpoint, Falcon3-1B MathIR at pass
  $6.5$, is reported by no seed of either Dr.GRPO arm and is the only break in
  the figure. It is left unjoined rather than interpolated across, so that
  break is an absence of evidence rather than a measured value of zero.}
  \label{fig:factorial-training-accuracy}
\end{figure}

\begin{figure}[H]
  \centering
  \includegraphics[width=\linewidth]{figures/factorial_training_curves_pmd.pdf}
  \caption{\textbf{Verified-mode diversity throughout training at all three
  model scales.} The same model rows and methods as
  Fig.~\ref{fig:factorial-training-accuracy}, now measuring \pmd{}. The
  population is not the same one. \pmd{} is reconstructed from the saved
  responses on the evaluation grid those responses were written on, so this
  figure is drawn against training step over a finer grid than the registered
  192-step grid of Fig.~\ref{fig:factorial-training-accuracy}, and each series
  carries its own seed cohort, of one to five seeds, rather than that figure's
  paired intersection. Counts are in the machine-readable record beside the
  plate. Because
  the metric conditions on the verified draws, a falling line reports that
  successes concentrated, not that they became rarer --- the distinction the
  accuracy curves cannot draw. The two fresh objectives fall in Graph,
  Countdown, and Pantry at every scale, reaching \pmd{} $.000$ in Pantry at
  0.5B and 1B, while the replay arms rise or hold. Lines are descriptive means
  over seeds and bands show seed ranges. A count of one has no band.
  Domain-specific vertical scales are shared across model rows; absolute
  heights across different domains should be read against their labeled axes.
  A step is plotted where at least five prompts return two verified responses.
  That is below the 30-prompt bar the tables report against, and deliberately:
  suppressing every step under 30 left panels blank where the measurement
  exists, and a blank reads as an absence rather than as the thin estimate it
  is. Under five prompts nothing is plotted, because a standard error of zero
  on two or three prompts records agreement rather than precision. Segments are
  never joined across a suppressed step, so a gap marks a step whose breadth
  could not be measured rather than one measured and found to be zero. Where
  the cohort is incomplete the mean rests on part of it: a solid stretch is one
  where every seed reports, a finely dotted stretch one where only some do. The
  distinction matters because a seed qualifies by succeeding often enough to
  define \pmd{}, so a mean over whoever qualifies can move with the qualifying
  rather than with the metric. Python is the domain to read this way, at every
  one of the three model scales drawn here.}
  \label{fig:factorial-training-modes}
\end{figure}

\begin{figure}[H]
  \centering
  \includegraphics[width=\linewidth]{figures/level2_factorial_training_curves.pdf}
  \caption{\textbf{Accuracy and registered verified-mode counts during Level-2
  training.} Qwen2.5-0.5B is the model tested at this level. Rows show
  \texttt{pass@8} and the registered \distk{} for all four primary methods;
  Fig.~\ref{fig:level2-training-pmd} gives the success-conditional reading. All
  five domains contain five terminal paired seeds per objective, as the
  per-panel counts record. Both metrics follow the same evaluation rules and
  the provenance admission of App.~\ref{app:registered-provenance}. Every
  checkpoint at this level is reported by at least one cohort seed, so no curve
  breaks; the ringed points are those resting on part of the cohort, and they
  fall early, where a requeue duplicated the first evaluations.}
  \label{fig:level2-training-curves}
\end{figure}

\begin{figure}[H]
  \centering
  \includegraphics[width=\linewidth]{figures/level2_training_curves_pmd.pdf}
  \caption{\textbf{Verified-mode diversity during Level-2 training.}
  The same model and methods as Fig.~\ref{fig:level2-training-curves},
  measuring \pmd{} on the grid the saved responses were written on rather than
  on that figure's registered paired grid, so the seed cohort is per series
  here and not the paired intersection there. Graph separates
  the arms most clearly: both replay arms rise to roughly $.34$ while Dr.GRPO
  and MaxRL fall to near $.11$, so the replayed policies end with about three
  times the chance that two verified colorings differ. MathIR collapses to
  \pmd{} $.000$ under every arm at this level. A step is drawn from five
  defining prompts upward, as in Fig.~\ref{fig:factorial-training-modes} and
  for the same reason; Countdown and Python are the panels that rest on that
  band rather than on the 30-prompt bar the tables use, and are read as thin
  estimates rather than as trajectories measured to the same precision as
  Graph. Segments are not joined across a suppressed step.}
  \label{fig:level2-training-pmd}
\end{figure}


\subsection{Supporting comparators}
\label{app:ucpo-progress}

We ran four objectives that act on this same pressure as controls: UCPO's
on-policy advantage redistribution, RLEP's reuse of successful trajectories, a
fixed semantic-entropy bonus, and replay weighted by how often fresh rollouts
rediscovered each mode. Their gains are local rather than general, and they fall on
correctness rather than on breadth --- UCPO helps Countdown correctness while
its \pmd{} there is $+.002$ ($[-.018,+.023]$) and its MathIR effects are
inconclusive, and RLEP raises MathIR correctness with \pmd{} exactly flat. The weighting control is the most direct of the four, because it
changes one thing: keep the same bank and the same reuse, and weight modes by
rediscovery frequency instead of equally. Equal weighting adds $+.318$ extra
modes per prompt over it ($[+.272,+.358]$) with no aggregate correctness
difference (App.~\ref{app:frequency-replay-progress}). Reusing successes is not
the active ingredient; refusing to let popularity decide how often each one is
rehearsed is.

GRPO changes group-reward normalization; UCPO redistributes fresh on-policy
advantage; sparse RLEP-Dr replays verified successes without uniform
canonical-key weighting. Fixed Semantic MaxEnt supplies a fresh-rollout
exploration incentive (App.~\ref{app:semantic-maxent}). These are
controls for alternative explanations of the retention results, and they
are read at the coverage they have rather than at the coverage the primary
comparisons have: GRPO runs at all three scales, UCPO and sparse RLEP-Dr
only at Qwen2.5-0.5B and Falcon3-1B with 50/50 and 47/50 registered runs,
Fixed Semantic MaxEnt completes 10/15 registered blocks, and the weighting
ablation completes all nine of its.


The alternatives give local rather than uniform benefits, and the
benefits are in correctness. Qwen-0.5B Countdown UCPO raises $P$ by $+.133$
($[+.005,+.262]$) while its \pmd{} there is $+.002$ ($[-.018,+.023]$), and its
MathIR effects are inconclusive. MathIR RLEP-Dr raises $P$ by $+.126$
($[+.030,+.222]$) with \pmd{} exactly flat; its Countdown and Falcon Pantry
effects are inconclusive. Seven five-seed blocks raise $P$ by an interval
clearing zero, and exactly one moves \pmd{} that far: Falcon Countdown GRPO,
by $-.105$ ($[-.147,-.063]$), a narrowing rather than a gain. The full
five-seed 3B plain-GRPO blocks have intervals including zero on both axes.
Each comparator is read against matched Dr.GRPO on the seeds where both define
\pmd{}, which leaves Qwen-0.5B Graph at one UCPO seed and four for RLEP-Dr,
and no Python cell at any scale.
RLEP's frozen offline pool was collected at temperature $.7$, top-$p=.95$;
all reported evaluations use the common temperature-1, top-$p=1$ contract.
An offline pool versus an online bank remains an algorithmic difference.

\begin{figure}[H]
  \centering
  \includegraphics[width=\linewidth]
    {figures/direct_baseline_learning_curves_static_strip.pdf}
  \caption{\textbf{Verified-mode trajectories for the supporting comparators.}
  Rows show Qwen2.5-0.5B and Falcon3-1B; columns are the five domains.
  UCPO and sparse RLEP-Dr are compared with Dr.GRPO and Re:Dr
  at these two smaller scales.
  Labels give exact method-specific seed counts. Lines are descriptive
  registered \distk{} means and bands are seed ranges. \distk{} rises with
  correctness as well as with breadth, which is the confound \pmd{} was
  adopted to remove, so these rows are read beside the accuracy panel rather
  than as a breadth measurement standing alone;
  Fig.~\ref{fig:factorial-training-modes} gives the \pmd{} reading for the
  primary methods, and Fig.~\ref{fig:direct-baseline-accuracy-curves} gives the
  matching accuracy curves, drawn on exactly the same seed populations and the
  same registered evaluation checkpoints. A checkpoint is drawn from the cohort
  seeds that report it, and the labelled counts are each cell's cohort rather
  than its thinnest checkpoint; the per-checkpoint counts are in the record
  beside the plate. Falcon3-1B Graph and MathIR carry the one checkpoint no
  seed of their replay arm reports.}
  \label{fig:direct-baseline-curves}
\end{figure}

\begin{figure}[H]
  \centering
  \includegraphics[width=\linewidth]{figures/direct_baseline_learning_curves_pass8.pdf}
  \caption{\textbf{Accuracy trajectories for the supporting comparators.}
  The same model rows, domains, methods, terminal seeds, and registered
  checkpoints as Fig.~\ref{fig:direct-baseline-curves}, now measuring
  \texttt{pass@8}. Bands show the seed range and the vertical axes span
  $[0,1]$ throughout. MaxRL and Re:Max at all three scales are compared
  on both metrics in Figs.~\ref{fig:factorial-training-accuracy} and~\ref{fig:factorial-training-modes}.}
  \label{fig:direct-baseline-accuracy-curves}
\end{figure}


\subsection{Diversity-preserving comparators}
\label{app:diversity-comparators}

Two recent objectives target output diversity directly, and both are run here
on the Level-1 Qwen2.5-0.5B panel at five domains and five seeds.
GAPO~\citep{anschel2025groupaware} replaces the task reward with a group
frequency-aware reward, $1-(f_i-1/L)$ over verifier-positive rollouts, where
$L$ is the prompt's valid-answer count. \mb{} certifies $L$, so the method can
be run on its own assumption rather than on an estimate. SetPO~\citep{li2026setpo}
adds each rollout's leave-one-out marginal contribution to its group's
kernelized set diversity, over frozen sentence embeddings of the generated
text; its coefficient was fixed before any scientific cell ran, at the value
placing the shaping term near the ratio our adaptive controller targets.

Both give local gains. GAPO moves \pmd{} on PantryPlan and Countdown and is
inert on PythonFactors, where its uniform target is unreachable: $L$ exceeds the
sixteen-rollout group on 95\% of those prompts and reaches 3600, leaving the
$1/L$ term numerically dead and the reward a plain duplicate penalty. MathIR is
a different case and not a GAPO failure --- no method moves it, the largest
effect on that column being $+.017$, because its certified support is dominated
by derivations the policy never produces. SetPO's one clear breadth gain is
PantryPlan. Neither reaches Re:Max's breadth across the panel, and the contrast
is where each acts: GAPO flattens frequency inside the fresh group, while
replay equalizes rehearsal across a bank of verified modes over time.

These two rows are read against a re-run Dr.GRPO control rather than the one
the other rows use. Their runtime could not be inherited, they sit on different
GPUs, and they take the corrected disjoint-draw evaluator. The substitution is
measurable and small: the re-run control's terminal \pmd{} sits $+.003$ from
the published control over twenty paired cells, per-cell $-.014$ to $+.034$,
against the effects of $+.225$ and $+.320$ the table reports.

\subsection{Uniform versus fresh-frequency replay}
\label{app:frequency-replay-progress}

\begin{figure}[!htbp]
\centering
\includegraphics[width=\linewidth]{figures/replay_key_weighting.pdf}
\caption{\textbf{Equal key weights can improve breadth beyond the first success.}
Uniform minus fresh-frequency weighting of verified replay, read on extra
verified modes because neither arm of this ablation is in the resampled \pmd{}
cohort. Success and extra modes are separate endpoints; a positive extra-mode
effect does not hold correctness fixed. The displayed model, domains, seeds, and uncertainty use
the frozen weighting comparison. The domain variation and inconclusive
success contrast limit the claim to the demonstrated breadth contribution.}
\label{fig:replay-key-weighting}
\end{figure}

The weighting ablation retains the same bank construction and update code,
changing only the within-bank target: uniform weights over retained keys
versus weights proportional to validator-positive fresh-rollout counts.
Policy-dependent bank contents can subsequently differ. Positive contrasts
below favor uniform weighting. The registered co-primary contrasts are
$\Delta B$ and $\Delta P$; the five-domain Qwen2.5-0.5B aggregate first
averages domain effects within each matched seed, giving equal domain weight.

Complete-block intervals enumerate all $5^5=3{,}125$ ordered bootstrap
resamples of the five paired training seeds and use linearly interpolated
2.5th and 97.5th percentiles. A resampled seed carries its complete domain
vector when estimating the five-domain aggregate. The endpoints and
aggregate were preregistered; the exact seed-resampling and percentile
conventions are analysis choices. These nominal intervals quantify seed
variability within the fixed benchmark, without resampling prompts or
domains or adjusting across comparisons.

The registered Qwen2.5-0.5B five-domain effect is $+.318$ extra modes per
prompt ($[+.272,+.358]$), with positive within-seed domain averages in all
five seeds. Graph, Countdown, and PantryPlan drive this result; Python and
MathIR remain inconclusive. The correctness effect is $+.010$
($[-.119,+.139]$), establishing neither a gain nor equivalence. Python's
correctness effects range from $-.781$ to $+.828$ across seeds despite
small extra-mode effects in four seeds.

Falcon Graph also favors uniform weighting on extra modes ($+.233$,
$[+.182,+.320]$). Its two correctness readings point different ways, and
neither establishes a loss: its registered $\Delta P$
is $+.045$ on an interval $[-.0004,+.0910]$ that
includes zero, while mean per-response correctness, which the table does not
carry, falls by $-.0152$ ($[-.0214,-.0090]$). Falcon Pantry is
inconclusive: $\Delta B=-.146$ ($[-.564,+.414]$) and $\Delta P=-.021$
($[-.063,+.026]$). Qwen2.5-3B Graph is complete at five pairs, with an extra-mode effect of
$+.079$ ($[+.011,+.135]$); its correctness interval includes zero.
Pantry also has five pairs, with $\Delta B=+.207$
($[+.038,+.377]$) and inconclusive correctness ($+.011$, $[-.020,+.043]$).
The evidence supports a contribution from key balancing in the stated domains, not a
uniform accuracy or breadth advantage across every model--domain block.

\input{results/semantic_current_summary_20260912.tex}

\section{Collapse and Replay Mechanism Checks}
\label{app:ablations}

\input{results/conditional_concentration_20260912_appendix.tex}

\begin{figure}[!htbp]
  \centering
  % One float, two panels: the same contrast read along the scale axis and
  % along the level axis. Equal minipage widths and matched native plate
  % widths, so neither panel is magnified relative to the other.
  \begin{minipage}[t]{.49\linewidth}
    \centering
    \textbf{(A)}\\[1pt]
    \includegraphics[width=\linewidth]{figures/concentration_story_all_scales.pdf}
  \end{minipage}\hfill
  \begin{minipage}[t]{.49\linewidth}
    \centering
    \textbf{(B)}\\[1pt]
    \includegraphics[width=\linewidth]{figures/concentration_levels.pdf}
  \end{minipage}
  \caption{\textbf{Concentration along both axes.}
  \textbf{(A)} \textbf{Concentration at every scale.}
  Fig.~\ref{fig:concentration-story} for all three scales rather than
  Qwen2.5-3B alone: final minus initial \pmd{} under Dr.GRPO, GRPO and MaxRL,
  five domains by three scales. Marks are means over five seeds with nominal
  95\% seed intervals; open marks carry fewer than five seeds and no interval,
  or name an objective \pmd{} cannot score because the initial checkpoint
  returns two correct responses to no prompt. The scale ordering is the same
  one Fig.~\ref{fig:base-levels-all-scales} uses.
  \textbf{(B)} \textbf{Training narrows breadth on constructions it never trained
  on.} Final minus initial \pmd{} for Qwen2.5-3B policies trained at Level 1 and
  read on each level's held-out split, so every row above L1 is transfer rather
  than training at that level. Marks are means over training seeds against the
  one untrained checkpoint per cell, with nominal 95\% seed intervals; the
  spread is across seeds, not a paired design. Open marks fall below the
  support bar, where the estimate stands but its precision does not. Dr.GRPO
  narrows at every level while MaxRL stays near zero, and Countdown crosses to
  the broadening side by Levels 4 and 5 --- the direction is not a property of
  training alone but of training and construction together. Rows labeled
  \emph{initial below support} have no untrained \pmd{} to subtract, because the
  initial checkpoint never returns two correct responses to one prompt on that
  split; the difference is undefined there rather than small. GRPO is read at
  Level 1 only, so its marks appear on the L1 rows alone and Levels 2--5 carry
  Dr.GRPO and MaxRL.}
  \label{fig:concentration-all-scales}
  \label{fig:concentration-levels}
\end{figure}


\subsection{Breadth before and after RLVR-only training}
\label{app:pre-rl-regression}

The problem precheck evaluates unmodified Dr.GRPO and GRPO at passes 0 and 8
on all three models and five domains: 150 complete runs, each with four
fixed $K=8$ draws. The five-domain macro-average of extra verified modes
falls by $98.2\%$--$99.4\%$ at Qwen2.5-0.5B and $41.5\%$--$71.6\%$ at
larger scales. Correctness falls in 10 of the 30 domain--method--scale
comparisons, while macro raw \texttt{distinct@8} ends above its initial
value at Qwen2.5-3B. This distinguishes loss of multiple correct outcomes
from the broader question of whether any correct outcome is sampled.
These are measured changes in the registered runs: they do not establish
zero probability for the lost keys, nor a universal decrease in raw breadth
across domains.

The same decomposition on the success-conditional axis gives close to the same
numbers, which is the check that matters here: extra verified modes could in
principle have fallen only because correctness fell. Macro \pmd{} declines by
$98.5\%$ and $99.8\%$ at Qwen2.5-0.5B under GRPO and Dr.GRPO, and by
$50.8\%$--$65.1\%$ at the larger scales, against the $98.2\%$--$99.4\%$ and
$41.5\%$--$71.6\%$ measured on extra modes. Terminal macro \pmd{} reaches
$.000$ for Dr.GRPO at Qwen2.5-0.5B. The collapse is therefore a property of how
verified probability is allocated, not an artifact of how much of it there is.
Falcon3-1B Countdown is the one arm where \pmd{} rises under Dr.GRPO, from a
near-degenerate $.033$; it does not lift the macro average above zero decline.

Initial evaluations are measured within each method. Qwen2.5-3B MathIR and
Pantry replay-arm seeds 71--74 carry slightly different initial requests from
plain GRPO ($.154$ versus $.162$, and $.832$ versus $.828$ on $P$), so those cells
have separate initial references.

\begin{figure}[t]
  \centering
  \includegraphics[width=\linewidth]{figures/baseline_collapse_precheck.pdf}
  \caption{\textbf{RLVR-only training reduces average extra verified modes
  at each tested scale.} Bars split the registered \distk{} into the first
  correct mode and every verified mode beyond it, at training pass 0 and pass
  8, for
  unmodified Dr.GRPO then GRPO in each panel. All 150 evaluations are complete at five
  seeds and four fixed \(K=8\) draws; whiskers are the seed range and pass 0 is
  measured per method. Badges give the change in extra verified modes, printed only
  where pass-0 breadth exceeds \(.05\) of them; the other three domains start at
  \(.011\) or less.
  \label{fig:baseline-precheck}}
\end{figure}




Graph and PantryPlan contain substantial initial breadth at every scale;
the other three domains begin with at most $.011$ extra modes per prompt.
A retained-breadth percentage is therefore reported only when the initial
extra-mode count exceeds $.05$. Raw $\texttt{distinct@8}$ ends below its
initial value in six of the 15 Dr.GRPO domains. Re:Dr does so in
two available domain summaries, both PantryPlan; Falcon Countdown uses
four admissible seeds. On the success-conditional axis the contrast is
sharper: averaged over the domains that remain measurable at the terminal
step, Dr.GRPO ends at \pmd{} $.017$, $.103$, and $.086$ at 0.5B, 1B, and 3B
against Re:Dr's $.233$, $.201$, and $.194$. These observations concern
sampled held-out support, not the survival histories of particular stored
modes.

Uniform-action references are meaningful for the enumerable Graph and
Pantry representations, and App.~\ref{app:metric-sampling} derives them. They
are not defined for open-ended programs or derivations. Pantry's bit-uniform reference exceeds every final model, so
this domain is a support-concentration stress test rather than evidence
of complete support coverage.

\begin{table}[t]
  \centering
  \caption{\textbf{Absolute references for the registered raw \distk{}.}
  \emph{Uniform} samples the exposed action representation uniformly and is
  defined only for Graph (27 color strings) and PantryPlan (64 bit masks).
  Frozen is the matched pass-0 model; the final two columns are pass-8 means.
  Dashes avoid pretending that an open-ended program or derivation space has a
  canonical uniform distribution.}
  \label{tab:absolute-support-references}
  \scriptsize
  \setlength{\tabcolsep}{4.2pt}
  \begin{tabular}{@{}llrrrr@{}}
    \toprule
    \textbf{Model} & \textbf{Domain} & \textbf{Uniform} & \textbf{Frozen}
      & \textbf{Dr.GRPO} & \textbf{Re:Dr} \\
    \midrule
    \input{results/absolute_support_references_table_body.tex}
  \end{tabular}
\end{table}


\subsection{Fresh-gradient degeneracy and observed output collapse}
\label{app:degenerate-groups}

A Dr.GRPO group with identical rewards has exactly zero fresh task gradient.
Only mixed groups containing both correct and incorrect responses can
supply fresh reward information; optimizer momentum can still move a
policy when that task gradient is zero. The fractions of mixed groups
over all 3,072 Qwen2.5-0.5B updates, averaged over five seeds, are:

\begin{center}
\small
\begin{tabular}{@{}lrr@{}}
\toprule
\textbf{Domain} & \textbf{Dr.GRPO} & \textbf{Re:Dr} \\
\midrule
Graph coloring & .144 & .954 \\
Countdown & .140 & .456 \\
Python factors & .015 & .185 \\
MathIR & .364 & .453 \\
PantryPlan & .123 & .789 \\
\bottomrule
\end{tabular}
\end{center}

Thus 98.5\% of Python control updates carry zero fresh task gradient.
Replay supplies a verified-likelihood gradient directly and can also
change the policy so that later fresh groups become mixed.

A raw-response audit of all five Qwen2.5-0.5B Python control seeds establishes
observed output collapse directly. At pass 0, a draw has 7.99 unique raw
completions per prompt on average. By steps 288--480, every seed produces
32 byte-identical temperature-1 samples per prompt, also identical to the
logged greedy completion, on all 128 prompts. This persists through step
3,072. Terminal draws contain 1.00 unique completion per prompt in the
control and 1.97 under replay. A collapsed completion can be incorrect:
this is one observed completion per prompt, not one correct key everywhere.
Aggregate endpoint equality or zero draw spread alone would not establish
that conclusion, nor does this finite audit rule out low-probability tails.

These checks explain a vulnerability of the registered sparse binary-reward
recipe. They do not establish dominance over methods that keep fresh
on-policy signal active. Dynamic mixed-group sampling, larger groups,
reference KL, token-entropy regularization, and wider-decoding evaluations
remain relevant controls beyond the registered factorial.

\subsection{Fixed-bank exemplar score trajectories}
\label{app:fixed-bank-survival}

The fixed-bank study tracks Qwen2.5-0.5B Re:Dr on Graph for seeds
43--47. Before step 384's fresh group, bank membership and fresh counts
are frozen; admission then stops while uniform replay continues through
3,072 updates under the same training recipe. Every frozen exemplar is
identified by prompt and canonical-key fingerprints, without selection by
its score trajectory. Each scheduled replay visit records mean token log
probability $s$ and sequence log probability $\ell$. Changes compare an
identity's first and final post-freeze visits; observation steps can differ
across identities.

All 3,406 frozen identities across 1,575 seed--prompt pairs have finite,
aligned scores at every scheduled visit: 8--9 visits per identity and
29,041 observations in total. The pooled median change is $+.038$
nat/token (95\% interval $[+.025,+.048]$), but the 10th percentile is
$-.256$ and $2.61\%$ of identities lose more than $.5$ nat/token
($[1.54,3.77]\%$). The worst intermediate change is $-1.328$ nat/token.
On the sequence scale, the median is $+.154$ nat/sequence
($[+.104,+.198]$), while $21.99\%$ lose more than $.5$ nat/sequence.
Lengths range from 4 to 100 tokens, so the two threshold percentages have
different meanings.

\begin{figure}[H]
  \centering
  \includegraphics[width=.95\linewidth]{figures/e121_fixed_bank_survival.pdf}
  \caption{\textbf{Fixed-exemplar scores improve centrally but retain a declining tail.}
  Empirical cumulative distributions show first-to-final post-freeze score
  changes for every frozen identity, with one curve per seed and the pooled
  identity distribution. Mean token and sequence scores are separate panels;
  positive values indicate increased teacher-forced likelihood. Neither panel
  measures the exact probability of generating a canonical mode during policy sampling.}
  \label{fig:e121-fixed-bank-survival}
\end{figure}


Pooled summaries weight identities equally. Descriptive 95\% intervals
use 10,000 hierarchical bootstrap draws (random seed 121), resampling prompts within each
original seed and then five seeds with replacement; repeated seeds reuse
their within-draw prompt sample. Linear empirical quantiles and percentile
intervals are analysis choices. The artifact
\path{paper/results/e121_fixed_bank_survival.json} retains all seed summaries,
identity trajectories, source hashes, and audit diagnostics.

These are teacher-forced score surrogates, not exact probabilities of sampling
canonical modes, and visits do not cover every intervening update; the
declining tail excludes a claim that every exemplar improves. Without a matched
no-replay fixed-bank arm, the study does not identify a causal replay effect,
and it does not numerically validate
Theorem~\ref{thm:replay-retention}'s floor. Its contribution is longitudinal
measurement of fixed exemplars alongside the held-out evidence.

\clearpage
\phantomsection
\begin{center}\large\textbf{Part IV. What breadth represents}\end{center}
\addcontentsline{toc}{part}{Part IV. What breadth represents}

\section{Matched Prompt-Hint Ablation}
\label{app:prompt-hint-ablation}

\input{results/modebench_prompt_ablation_20260911.tex}

\paragraph{Interpreting the intervention.}
The Python intervention removes an executable nested-conditional example
along with divisor-testing and dispatch guidance; it also shortens the prompt.
Its paired effect concerns this combined removal and does not isolate a
preference cue from example or length effects. A reduction in
\texttt{distinct@8} accompanied by lower \texttt{pass@8} is not by itself
evidence of greater concentration among correct outputs: both endpoints fall
when the model simply succeeds less often. This is precisely the ambiguity
\pmd{} removes, since it is estimated from the verified draws alone, and it is
the reading to consult when the two registered endpoints move together; the two
rightmost columns of Table~\ref{tab:prompt-hints-local-strict} report it for
every cell that can carry it.
Here the initial model is the single fixed Qwen2.5-0.5B-Instruct checkpoint
before the reported RL training; it is already instruction-tuned and has no
fine-tuning-seed replication in this ablation.


\paragraph{Observed prompt dependence.}
Python shows a large loss of correctness after guidance removal. Across the
five DrGRPO seeds, mean $P_8$ falls from $.8125$ to zero at Level~2 and from
$1.0$ to zero at Level~3; $D_8$ falls identically and $B_8=0$ in both arms.
With no correct neutral DrGRPO outputs, its conditional-correct concentration
is undefined. Re:Dr retains neutral $P_8=.55625$ and $.49375$ and
$D_8=.675$ and $.65625$ at Levels~2 and~3, respectively, while losing
additional breadth: $\Delta B_8=-.2375$ and $-.25625$. The initial checkpoint
also loses Python correctness (Table~\ref{tab:prompt-hints-local-strict}).
These observations identify substantial dependence on the combined Python
wording intervention; a raw mode-count decline cannot be interpreted solely
as a redistribution among correct answers.


\paragraph{Post-hoc Python failure audit.}
A descriptive audit of all 5,120 saved DrGRPO Python draws (five seeds,
two levels, both wordings) finds that all 2,560 original outputs extract
the same lambda, whose body matches the conditional example in the prompt:
\begin{quote}\footnotesize\ttfamily
lambda n: 2 if n \% 2 == 0 else 3 if n \% 3 == 0 else 5
\end{quote}
Of the 2,560 neutral draws, 1,816 (70.94\%) are invalid Python expressions,
59 violate other restricted-language rules, 376 return non-integers,
276 fail the proper-divisor condition, and 33 lack an extracted candidate
after reaching the token limit. Thus most neutral failures involve syntax;
truncation accounts for only 1.29\%. The unchanged extractor, parser, and
executor reproduce every stored grade and successful canonical key.
This taxonomy was developed after the paired results were known and does
not identify an internal training mechanism. Counts, deterministic paired
examples, and source hashes are retained in
\path{results/modebench_prompt_ablation_python_failure_diagnostic_20260911.json}.

\paragraph{Concentration after hint removal.}
Every trained MathIR checkpoint has $B_8=0$ under both wordings at both levels.
Within each problem's eight-draw group, every observed pair of correct
outputs shares a canonical key, in every trained checkpoint--level cell:
correct-pair collision is $1.0$, versus the
certified-support uniform reference of $.2$ --- a reference this domain
cannot reach.
Table~\ref{tab:prompt-hints-local-strict}
carries the same fact on the other axis, as \pmd{} $.000$ under both wordings in
every reportable MathIR cell: the hint is not what holds those answers together. For example, at Level~2, Replay
DrGRPO retains mean $P_8=.975$ without the order hint, compared with $.99375$
originally. Across its five checkpoints, the neutral arm has 4,179 colliding
pairs out of 4,179 correct pairs, from 154 eligible checkpoint--problem
groups; the original arm has 4,382 out of 4,382 pairs from 157 such groups.
Both arms evaluate the same 32 problems at each checkpoint. This preserves
the observed concentration despite removing the specified order hint; it does
not measure all unseen modes or remove the other interface constraints.
Pantry effects are mixed. Re:Dr at Level~3 gains $.109375$ in $D_8$,
comprising $.046875$ in $P_8$ and $.0625$ in $B_8$, but this comparison has only
two training seeds. The separate completed 64-draw hosted follow-up
tests both wordings on sixteen problems per cell in
App.~\ref{app:sampling-budget-ablation}; its population and draw budget
remain distinct from this local eight-draw experiment.


\clearpage
\section{The Decoding Grid Behind the Objection}
\label{app:decoding-objection}

\paragraph{Does a different decoding setting reopen a collapsed policy?}
The cheapest remedy would be to turn the temperature up. We put every frozen
terminal checkpoint of the cohort the rest of the paper reports back through
the ordinary evaluation path at six temperatures from $0.5$ to $2.0$ --- five
domains, two arms, five seeds, 300 measured checkpoint cells
(Fig.~\ref{fig:decoding-objection}A). A RLVR-only policy stays where it
collapsed at every one of them. In Graph its \pmd{} never exceeds $.021$ across
the entire sweep: two correct colourings are almost never two different
colourings, at any setting we tried, while its paired Re:Dr arm sits at $.562$
at ordinary settings. The largest value any temperature buys any control is
Graph's $.021$, against replay arms between $.330$ and $.562$ in the three
domains where breadth is available at all. Decoding is not where the missing
modes went.


\paragraph{What one cell is.}
Each cell loads one terminal checkpoint, evaluates it once through the ordinary
training evaluation path at a chosen $(T,\text{top-}p,K)$, and exits; no
optimizer step, rollout or export occurs, so a cell is a pure function of the
checkpoint and its decoding settings. Every non-decoding argument --- prompt
template, split, response length, draw seeds --- is inherited from the
configuration the source run itself recorded rather than retyped. Pooling and
the support bar follow App.~\ref{app:metric-sampling} unchanged, and every value
is recomputed from the retained per-draw records, whose canonical keys were
written by the validator that graded them; nothing is resampled or regraded.

\paragraph{Evaluation hardware, and what it costs.}
GPU model changes floating-point reduction order and therefore the sampled
tokens, so a sweep that mixed hardware would fold it into the trade-off the
sweep exists to isolate. This grid is measured on one GPU model throughout,
which is the stronger guarantee for reading across the table: no comparison
within it carries a hardware difference. It is not the same as measuring each
checkpoint on the machine that trained it, and the difference is bounded rather
than assumed. Of 200 reproduction checks against the published terminal
endpoints, the 72 on checkpoints whose training model matches the evaluation
model reproduce \texttt{pass@8}, \texttt{mean@8}, \distk{} and the greedy trace
exactly, to the last representable unit; of the 128 measured across models, 125
fall inside the same tolerance, the three exceptions being the greedy trace
twice and \texttt{mean@8} once, which are the statistics most sensitive to
reduction order. A separate control re-measured a100-trained checkpoints on the
evaluation model at $T=1$: mean $|\Delta\pmd{}|$ of $.003$ over twelve paired
comparisons, maximum $.013$, against replay-minus-control gaps near $.5$.
Changing GPU model moves these measurements by thousandths where the effect
under study is halves.

\paragraph{A wider budget and a truncated nucleus.}
Temperature is the version of the objection that gets asked, but it is not the
only knob. An earlier Qwen2.5-0.5B cohort --- twelve passes to step 4,609, with
a replay arm predating the Re:Dr objective contract --- was swept over a wider
grid on the same benchmark, splits, verifier and evaluation path: the same six
temperatures, a quadrupled budget of $K=32\times2$ draws at
$T\in\{1.0,1.3,1.6\}$, and nucleus truncation to top-$p$ $.95$ at
$T\in\{1.0,1.6\}$, for 110 measured cells
(Fig.~\ref{fig:decoding-objection}B). Its control behaves as the current
cohort's does, and neither the larger budget nor the truncated nucleus reopens
it either. Because that cohort is not the one the rest of the paper reports,
its control range is read here as the wider sweep of the same question rather
than as a matched comparison, and its replay marks are a scale reference only.

% A two-inch strip: [t] keeps it at the head of a text page instead of
% letting it claim a float page of its own and leave half of it blank.
\begin{figure}[t]
  \centering
  \includegraphics[width=\linewidth]{figures/decoding_objection.pdf}
  \caption{\textbf{No decoding setting reopens a collapsed policy.} Both
  sweeps are Qwen2.5-0.5B terminal checkpoints, five seeds, evaluated through
  the ordinary evaluation path. Each bar is the Dr.GRPO control's \pmd{} range
  over every setting in that sweep which clears the support bar; each marker is
  its paired replay arm at ordinary settings ($T=1$, top-$p$ $1$, $K=8$), for
  scale. \textbf{A} is the cohort reported throughout, at six temperatures from
  $0.5$ to $2.0$, 300 measured checkpoint cells. \textbf{B} is an earlier
  twelve-pass cohort swept over eleven settings on the same benchmark, splits,
  verifier and evaluation path --- the same temperatures, a quadrupled budget of
  $K=32\times2$ draws at $T\in\{1.0,1.3,1.6\}$, and nucleus truncation to
  top-$p$ $.95$ at $T\in\{1.0,1.6\}$, for 110 measured cells. That cohort's
  replay arm predates the Re:Dr objective contract, so its markers are a scale
  reference rather than a matched contrast, and the control range is what the
  wider grid contributes. Python never clears the bar under either control,
  which is itself the collapse reported in Section~\ref{sec:results} rather
  than a gap in the sweeps. A bar drawn narrower than $.004$ is padded to stay
  visible.}
  \label{fig:decoding-objection}
\end{figure}


\paragraph{What the grids do and do not establish.}
Both measure decoding of frozen terminal policies. Neither shows that a control
\emph{trained} differently, or trained at another temperature, would not close
the gap; temperature at evaluation is not temperature during rollout. They
cover one model scale. One reading to avoid: temperature $2.0$ degrades
correctness in Countdown and Python, so a \pmd{} that fails to rise there says
the policy has stopped solving the task, not that its breadth was unreachable.
The records \path{paper/results/decoding_objection_e125.json} and
\path{paper/results/decoding_objection_e72.json} retain every cell, its seeds,
its defined-prompt count and its source path;
\path{ops/build_paper_decoding_objection_e125.py} and
\path{ops/build_paper_decoding_objection.py} regenerate them and the displayed
tables.

\clearpage
\section{Sampling-Budget Ablation}
\label{app:sampling-budget-ablation}

\input{results/modebench_discovery_curves_20260911.tex}

\input{results/discovery_hosted_interpretation_20260912.tex}

\input{results/gpt56_all_levels32_sampling_20260913.tex}

\subsection{Observed sampling-budget effects}
\label{app:sampling-budget-effects}

The full experiment is complete: 36,864 hosted and 110,592 local responses,
147,456 in total. Each of the six hosted model--wording cohorts has all
6,144 planned draws; strict and normalized analyses retain every response.

The completed local panel contains 110,592 fresh responses from all 25
registered checkpoints. Each trained checkpoint is evaluated in its own domain
at both levels; the initial checkpoint covers all six cells. Serial regrading
corrects one original Python grade from failure to success, confirmed once
in an independent warmed worker using the unchanged verifier and timeouts.
The saved receipt does not identify the cause of the original discrepancy.
No response is rescued by formatting normalization, so strict and normalized
analyses coincide throughout this panel. The original responses and grades
remain available alongside the correction audit and source identities.

\paragraph{Two weightings of the same correct pairs.}
Tables~\ref{tab:discovery-breadth-frontier} and~\ref{tab:discovery-breadth-local}
carry the full-64 pooled collision beside the promptwise weighting this paper
uses everywhere else (App.~\ref{app:metric-sampling}): every prompt with at
least two verified draws contributes one $\widehat{\pmd}$, and eligible
prompts count equally rather than in proportion to $\binom{c}{2}$. Across the
70 defined model--cell--wording entries the two readings agree closely, with a
median absolute difference of $.003$ between promptwise \pmd{} and pooled
$1-C$, so no reading in this appendix turns on the choice. They do separate
where one prompt supplies most of a cell's correct pairs, and in both
directions: under original wording the initial checkpoint has pooled
$1-C=.151$ against promptwise $.587$ on Level-2 PantryPlan, and pooled $.625$
against promptwise $.307$ on Level~3. Nineteen entries differ by at least
$.05$. The promptwise column is the one to read as breadth at matched
correctness; the pooled column remains a descriptive rate over a cell's
retained pairs, in which the prompts a model solves most often weigh most.
What no weighting removes is the selection: a prompt enters either column
only when it clears two verified draws, so the eligible population still
tracks correctness. It is thinnest where the two readings diverge most ---
both initial-checkpoint PantryPlan entries above rest on 7 of 16 problems
--- and it can differ sharply between the wordings, as at initial-checkpoint
Level-3 Python, eligible on 16 problems under the original wording and on
one under the neutral. Promptwise \pmd{} equalizes what each eligible prompt
contributes, not which prompts are eligible, which is the joint-selection
caveat of App.~\ref{app:metric-sampling}.

\paragraph{MathIR: concentration largely persists through 64 draws.}
At Level~2, both trained methods have $B_k=0$ at every measured budget under
both wordings. Re:Dr has original $P_8=.9774$ and $P_{64}=.9875$;
neutral wording gives $.9767$ and $.9875$. Its $D_k$ equals $P_k$ throughout,
so the small late gain consists entirely of first successes. At exactly eight
correct draws, both wordings have $R_8=1$ on the same 76 eligible
checkpoint--problem groups, versus the certified-support uniform lower
reference $5[1-(4/5)^8]=4.1611$. These groups reuse the same 16 problems
across five checkpoints. At Level~3, all neutral trained groups and all
original Re:Dr groups also contain at most one observed correct key.
There is one exception in original DrGRPO: a single problem at seed~46
contains two keys with counts 28 and 9, giving mean $B_{64}=.0125$ and
$B_{64}-B_8=.0035$. Its equal-seed correct-pair collision is $.9803$;
the other trained MathIR cells have collision $1.0$, against the uniform
upper reference $.2$. Thus the larger budget reveals a rare additional mode
without supporting an exact single-mode claim for every trained condition.

\paragraph{Python: prompt dependence and late discovery coexist.}
Original DrGRPO has unchanged $P_k=D_k=.8125$ at Level~2 and $1.0$ at
Level~3 from eight through 64 draws. Neutral DrGRPO has zero correct responses
out of 5,120 at Level~2 and three out of 5,120 at Level~3. Each of those
three successes occurs alone in its checkpoint--problem group, so Level~3
mean $P_{64}=D_{64}=.0375$ while correct-pair collision remains undefined.
Re:Dr behaves differently. Under original wording its pass probability
is already unchanged from eight to 64, but $D_k$ gains $.0582$ at Level~2
and $.1529$ at Level~3. Under neutral wording, $P_k$ rises from $.5186$ to
$.7500$ and from $.5109$ to $.7500$, respectively; $D_k$ rises from $.6618$
to $1.1250$ and from $.6561$ to $1.1875$. These final mode counts remain
below the original $1.2750$ and $1.5750$. The corresponding additional-mode
gains are $.2318$ [$.0574,.4565$] and $.2923$ [$.0659,.5445$], with the
registered seed-and-prompt intervals. Larger budgets therefore recover both
first successes and additional verified alternatives in the neutral replay
condition. These unconditional gains do not identify a change in the
conditional distribution of correct solutions by themselves.

\paragraph{Pantry: eight draws understate observed discovery.}
Across the trained Pantry cells and both wordings, $D_{64}-D_8$ ranges
from $.4151$ to $1.0247$, with additional-mode gains ranging from $.2739$
to $.7664$. The initial checkpoint also gains $.8565$--$1.0778$ modes.
In the two observed training seeds, Level~3 point estimates even change
method ordering: original DrGRPO/Re:Dr have $D_8=.2075/.3381$
but $D_{64}=1.1563/.9063$; neutral wording gives $.2253/.3515$ at eight
and $1.2500/.8125$ at 64. This is a descriptive ranking of these checkpoints,
not evidence of a population-level method advantage. Pantry has only two
training seeds, and fixed-correctness comparisons at larger $m$ often have
few or no jointly eligible problems. The eligibility table and undefined
means remain essential to interpreting the conditional curves.

These observations support persistent concentration in several trained
conditions and substantial late discovery in others, which is why no single
saturation claim covers all domains: a flat finite-sample curve still cannot
exclude rarer unseen modes. The $k=8$ points rarefy the fresh 64-draw pools on
16 problems per cell rather than restating the previous experiment's
32-problem scores. The completed hosted panel below
adds a separate comparison of fixed deployments on these historical tasks;
it is not evidence from newly trained Level-3 checkpoints.


% Concise theory for both papers; full derivations are preserved in
% paper/audits/streamlining_20260911/full_theory.tex.

\input{results/inference_followups_20260912.tex}
\input{results/portfolio_withdrawals_20260917.tex}

\clearpage
\phantomsection
\begin{center}\large\textbf{Part V. Fixed deployments}\end{center}
\addcontentsline{toc}{part}{Part V. Fixed deployments}

\input{results/frontier_hosted_20260911_appendix.tex}

\subsection{Hosted concentration on the success-conditional axis}
\label{app:hosted-mode-diversity}

Fig.~\ref{fig:hosted-verified-breadth} reports \pmd{} for the same audited
responses, per domain and level and for every deployment in the cohort rather
than this one alone. The hosted tables above report \texttt{correct\_pair\_collision},
which is one minus \pmd{} pooled over pairs; the values here instead average
over prompts without weighting by how many verified pairs each contributed,
matching the aggregation used everywhere else in this paper and stated in
App.~\ref{app:metric-sampling}. The two readings agree closely,
and the pooled reading recomputed from the per-sample records reproduces the
published collision estimate exactly, which checks the two pipelines against
each other.

This deployment is highly accurate and, on three of five domains, close to
deterministic in which verified answer it gives. MathIR reaches \pmd{} $.000$
at Levels 2 and 3: across every prompt with two or more verified draws, no two
verified answers ever differed. Level-3 Graph sits at $.049$ against six
available colorings. PantryPlan is the consistent exception. Python at Levels 2
and 3 is left blank because the model clears two verified draws on only 14 and
15 of 128 prompts there, which is too few to estimate breadth.

\paragraph{The Level-4 probe.}
Fig.~\ref{fig:base-levels-all-scales} carries one Level-4 point per domain for
GPT-5.6 Sol, at \pmd{} $\MDfrontierLpmdFour$ over $\MDfrontierLdomainsFour$
domains. Two things separate it from every other hosted cell and it is read
on its own account for both of them. It was
served direct through OpenAI rather than the Azure deployment that served the
other levels, which may be a different build; the served model string is
recorded per response rather than assumed. And its Python cell rests entirely on
the frozen formatting normalizer: strict grading rejects the LaTeX-escaped
Python this deployment emits, verifying two of its 1,024 draws where the
normalizer recovers 675. That is a presentation artifact rather than a capability
measurement, and it affects the Levels 1--3 Python cells in the same direction.

Table~\ref{tab:hosted-cohort-mode-diversity} extends the measurement to the
whole completed cohort. Every one of the seven deployments is concentrated on
this axis, and none reaches $1.5$ effective verified modes: macro \pmd{} runs
from $.142$ for Claude Opus 5 to $.320$ for DeepSeek V4 Pro, or $\neff$ between
$1.17$ and $\MDneffceiling{}$. Those macro averages are taken over the thirteen
model--domain--level cells every deployment reports. That ceiling is partly
MathIR's doing, and the last two columns say by how much. MathIR keys a mode by
its derivation rather than its answer, and most of the derivations it certifies
are longer routes no measured policy produces, so it sits near zero for every
deployment for reasons that are partly the construction's. Averaged over the four domains whose modes are alternative answers rather than
alternative derivations of one answer, every deployment rises and the ceiling
reaches $\MDneffceilingnomathir{}$ --- still under one and three fifths
effective verified modes, against domains admitting between four and several
hundred.
The concentration is not an artifact of that domain; the exact height of the
bound is partly one.

 Two of them, GPT-5.6 Sol
and Claude Opus 5, return no measurable Python breadth at Levels 2 and 3
because the provider refused most of those requests, and averaging each
deployment over whatever it can report would rank them against an easier set of
cells than the other five: it puts DeepSeek V4 Pro at $.280$ rather than $.320$
and moves it below Kimi K3. The ordering does not follow correctness, and it is not stable
across domains --- GPT-5.4 and Kimi K3 lead on Python while DeepSeek V4 Pro
leads on Graph --- so this is a property of individual deployments rather than
a single frontier-wide constant. What is common to all of them is the level:
the most varied deployment measured still repeats its own verified mode on
roughly seven of every ten pairs of correct answers. Every model also declines
from Level 1 to Level 3. Fig.~\ref{fig:frontier-level-grid} draws this
cohort one deployment per row, beside the open-scale grid it mirrors.

\begin{table}[!htbp]
  \centering
  \caption{\textbf{Verified-mode diversity across seven hosted deployments.}
  Macro \pmd{} averages the cells that clear the support bar, weighting the
  five domains and three levels equally, with $\neff=1/(1-\pmd{})$ beside it
  and the per-level macro to its right. Each deployment contributes 128
  held-out prompts per cell and eight stateless requests per prompt under its
  own default sampling. GPT-5.6 Sol and Claude Opus 5 clear the bar on 13 of
  15 cells; the other five clear all 15.}
  \label{tab:hosted-cohort-mode-diversity}
  \small
  \begin{tabular}{lcccccc}
    \toprule
    & & & \multicolumn{3}{c}{Macro \pmd{} by level} \\
    \cmidrule(lr){4-6}
    Deployment & Macro \pmd{} & $\neff$ & L1 & L2 & L3 \\
    \midrule
    \input{results/mode_diversity_hosted_cohort_table_body.tex}
  \end{tabular}
\end{table}

\paragraph{What the level axis carries besides difficulty.}
The Level-2 and Level-3 system messages add explicit strategy guidance to
MathIR, Python and PantryPlan --- an algebraic isolation order, a nested
small-divisor chain, a preference over ingredient families --- which the
Level-1 messages do not carry;
App.~\ref{app:prompt-hint-ablation} gives the exact text and removes it. A
deployment that follows an instruction concentrates on the mode that
instruction names, so some of the decline above belongs to the prompt rather
than to the task. How much is measured here rather than assumed. Averaged over
the seven deployments, \pmd{} falls by \MDhintguidedfall{} from Level 1 to
Level 3 across the \MDhintguidedpairs{} guided deployment--domain pairs and by
\MDhintplainfall{} across the \MDhintplainpairs{} pairs in Graph and Countdown,
which carry no strategy text at any level. Averaged over those two unguided
domains alone, \MDhintplainfalling{} of \MDhintplaindeployments{} deployments
still decline. The guidance inflates the cross-level fall without accounting
for it.

Two further readings bear on the same question, and they do not agree. Removing
the MathIR guidance from local checkpoints leaves \pmd{} at $.000$ under both
wordings, so at 0.5B the hint is not what holds those answers together
(App.~\ref{app:prompt-hint-ablation}). At 64 draws on hosted deployments,
neutral MathIR wording instead raises breadth at matched correctness for both
GPT deployments at Level 2, by intervals excluding zero, while Grok's include
it (App.~\ref{app:hosted-discovery-complete}). The second is the one measured
on deployments of this kind. Read the cross-level narrowing as concentration
under the prompts as specified --- which is what these deployments are asked,
and what a user of them receives --- rather than as an effect of difficulty
alone.


% Render this retained display inside its generated numerical appendix.
\AddToHookNext{file/after/gpt56_temperature_curve_20260911_appendix.tex}{%
\begin{figure}[H]
  \centering
  \includegraphics[width=.98\linewidth]{figures/gpt56_temperature_curve.pdf}
  \caption{\textbf{Temperature traces pass@8 and verified breadth.}
  GPT-5.6 Sol, with 480 fixed prompts and eight draws each at $T\in\{0,0.5,1,1.5,2\}$.
  Left: all domains and levels; right: separate levels, averaging five domains
  equally. Empirical \texttt{pass@8} is the fraction with at least one correct
  draw; the registered \distk{} averages unique correct modes, including
  zeroes, and so tracks \texttt{pass@8} across temperatures by construction.
  Connected points use reasoning \texttt{none} and the frozen formatting
  normalizer.
  App.~\ref{app:gpt56-temperature} gives pointwise uncertainty intervals for
  every benchmark level at each of the five measured temperatures, together with
  the per-domain values behind each plotted point.}
  \label{fig:gpt56-temperature-curve}
\end{figure}
}

\input{results/frontier_comparison_20260911_appendix.tex}

\input{results/hosted_level_averages_20260911.tex}

\begin{figure}[H]
  \centering
  \includegraphics[width=\linewidth]{figures/hosted_verified_breadth.pdf}
  \caption{\textbf{Accurate hosted outputs can still cover few verified modes.}
  Columns show domains; markers distinguish Levels 1--3. Each point uses all
  eight responses to each of 128 prompts. Top: per-response accuracy.
  Bottom: \pmd{}, which conditions on success and so separates breadth among
  correct answers from accuracy itself. A fixed formatting normalizer applies
  to every
  deployment; Opus 5 Python uses revised task wording without a system message.
  Hollow marks in that one column repeat the original-prompt cells the revised
  ones replace, on the wording the other six deployments answer. This is the
  per-domain and per-level reading for every audited deployment.
  Table~\ref{tab:hosted-level-averages-medium} reports the corresponding
  per-level averages of \texttt{pass@8} and \texttt{distinct@8} across all five
  domains, beside the \pmd{} of the same draws.
  Table~\ref{tab:hosted-level-averages-parity} repeats the first two over the
  four domains that share one prompt.}
  \label{fig:hosted-verified-breadth}
\end{figure}

\input{results/hosted_level_averages_20260911_parity.tex}

\input{results/gpt56sol_python_parity_20260918.tex}

\input{results/hosted_reasoning_off_20260912_main.tex}

\input{results/hosted_reasoning_off_20260912_appendix.tex}

\clearpage
\phantomsection
\begin{center}\large\textbf{Part VI. Idealized analysis}\end{center}
\addcontentsline{toc}{part}{Part VI. Idealized analysis}

\section{Mode Collapse and Verified Replay}
\label{app:theory}

Under the explicit categorical assumptions below, centered binary-reward
estimators amplify a pre-existing imbalance among correct modes,
whereas persistent verified replay supplies a loss barrier for banked modes.
The complete-response argument then states what would be needed to transfer
that protection to a language model. These results concern idealized update
rules; they do not certify the implemented AdamW/PPO trajectories.
Table~\ref{tab:notation} lists the symbols used here.

\subsection{Categorical abstraction and mean updates}
\label{app:theory-grpo-mean}

Every result in this appendix is stated under one model, fixed here once.

\begin{assumption}[Categorical mean-flow model]
\label{ass:categorical-mean-flow}
Fix one isolated prompt and a finite category set
$\mathcal A=\mathcal C\mathbin{\dot\cup}\mathcal N$ with $m\ge2$ correct
execution modes and $n\ge1$ incorrect categories, and set
$p_a=\operatorname{softmax}(z)_a$, $P=\sum_{c\in\mathcal C}p_c$, and
$q_c=p_c/P$, so that $P$ measures correctness while $q$ describes how that
correct probability is allocated among the individual correct modes.
\begin{enumerate}
\item[(A1)] \emph{Independent logits.} The policy carries one independently
  optimized logit for every category in $\mathcal A$, including one separate
  logit for each canonical correct mode.
\item[(A2)] \emph{Common length normalization.} Every response shares a single
  length scale, which is absorbed into the positive coefficient of the group
  estimator, rather than a scale that varies from one response to the next
  within the same rollout group.
\item[(A3)] \emph{Euclidean mean flow.} Updates are the infinitesimal expected
  updates in these logits, taken without optimizer preconditioning and without
  momentum carried between steps.
\item[(A4)] \emph{No competing regularizer.} Explicit reference KL is zero and
  no entropy term is present.
\item[(A5)] \emph{Nondegenerate start.} Logits are finite initially,
  $0<P(0)<1$, and the group size is $G\ge2$.
\end{enumerate}
\end{assumption}

We write ``under Assumption~\ref{ass:categorical-mean-flow}'' for (A1)--(A5)
together. Each limitation of the analysis can be traced to a specific item:
AdamW and PPO violate (A3), variable-length GRPO violates (A2), and entropy
regularization violates (A4).

A language model induces probabilities over canonical outcomes, but this
identification of probabilities does not identify its update geometry.
Aggregating several response strings into a mode does not generally commute
with gradient flow in response logits. One independent logit per canonical
outcome is therefore an explicit abstraction, not a reduction proved for the
neural policy or its validator equivalence classes.

Draw a group $Z_1,\ldots,Z_G\stackrel{\mathrm{iid}}{\sim}p$ with $G\ge2$,
write $R_i=\1[Z_i\in\mathcal C]$, and let $R=\sum_iR_i$. At the on-policy
point where the PPO ratio equals one, consider the common-length estimator
\begin{equation}
  \widehat g_G
  =\frac1G\sum_{i=1}^G
   w(R)\left(R_i-\frac{R}{G}\right)\nabla_z\log p_{Z_i}.
  \label{eq:group-estimator}
\end{equation}
For the Dr.GRPO abstraction, $w(R)=1$ after absorbing the shared positive
length scale. Reward-standardized GRPO with equal or common response-length
normalization takes $w(R)$ to be the reciprocal group reward standard
deviation on mixed groups (with a fixed numerical stabilizer if used).
Standard variable-length GRPO also weights each response by its own inverse
length \citep{shao2024deepseekmath,liu2025understanding}; that factor cannot
generally be absorbed into $w(R)$ and is outside the equivalence proved here.
In either included case, $0<w(\ell)<\infty$ for
$\ell\in\{1,\ldots,G-1\}$. We define the entire update as zero on all-equal
groups without evaluating an undefined reciprocal standard deviation;
likewise, $w(R)R(G-R)$ below is defined as zero when $R=0$ or $R=G$. All
advantages and normalization weights are detached in the score estimator.

\begin{lemma}[A group baseline does not remove frequency amplification]
\label{lem:grpo-mean}
For $0<P<1$, the expected update in Equation~\eqref{eq:group-estimator} is
\begin{equation}
  \E[\widehat g_G]
  =c_G(P)\nabla_z P,
  \qquad
  c_G(P)=
  \frac{\E\!\big[w(R)R(G-R)\big]}{G^2P(1-P)}>0.
  \label{eq:expected-group-gradient}
\end{equation}
For the Dr.GRPO abstraction, $c_G(P)=(G-1)/G$. Group centering and reward
standardization rescale the expected gradient without changing its direction
under these assumptions.
\end{lemma}

\begin{proof}
Condition on the complete binary reward vector, whose sum is $R$. The score
identities give
\[
  \E[\nabla_z\log p_{Z_i}\mid R_i=1]=\nabla_z\log P,
  \qquad
  \E[\nabla_z\log p_{Z_i}\mid R_i=0]
  =\nabla_z\log(1-P).
\]
There are $R$ correct rows with centered reward $(G-R)/G$ and $G-R$
incorrect rows with centered reward $-R/G$. Therefore the conditional
expectation of the unaveraged sum in~\eqref{eq:group-estimator} is
\[
  w(R)\frac{R(G-R)}{G}
  \left(\nabla_z\log P-\nabla_z\log(1-P)\right)
  =w(R)\frac{R(G-R)}{GP(1-P)}\nabla_zP.
\]
The outer factor $1/G$ and expectation over
$R\sim\operatorname{Binomial}(G,P)$ yield~\eqref{eq:expected-group-gradient}.
For $w\equiv1$,
$\E[R(G-R)]=G(G-1)P(1-P)$, proving the Dr.GRPO value.
\end{proof}

The binomial identity also gives the denominator-free expression
\[
 c_G(P)=\frac{G-1}{G}\sum_{j=0}^{G-2}\binom{G-2}{j}
 w(j+1)P^j(1-P)^{G-2-j}.
\]
The binomial weights sum to one, so $c_G$ lies between
$(G-1)\min_j w(j+1)/G$ and $(G-1)\max_j w(j+1)/G$.
This establishes the smooth, positive extension to $P=0,1$ used in the
subsequent flow and convergence arguments.

The next lemma specializes the practical-estimator analysis of
\citet[Section~4.3 and App.~D, Theorem~5]{tajwar2026maxrl}
to our categorical notation; we include the derivation for completeness.

\begin{lemma}[Binary MaxRL has the same correctness-gradient direction]
\label{lem:maxrl-mean}
Under Assumption~\ref{ass:categorical-mean-flow}, let MaxRL assign
$A_i^{\mathrm{MaxRL}}=GR_i/R-1$ when $R>0$ and zero when $R=0$, as in
Section~\ref{sec:maxrl-method}. For
$\widehat g_G^{\mathrm{MaxRL}}=G^{-1}\sum_i
 A_i^{\mathrm{MaxRL}}\nabla_z\log p_{Z_i}$ at the on-policy point, the
expected update is
\begin{equation}
  \E[\widehat g_G^{\mathrm{MaxRL}}]
  =c_G^{\mathrm{MaxRL}}(P)\nabla_zP,
  \qquad
  c_G^{\mathrm{MaxRL}}(P)
  =\frac{1-(1-P)^{G-1}}{P}>0,
  \label{eq:maxrl-expected-gradient}
\end{equation}
with continuous endpoint values
$c_G^{\mathrm{MaxRL}}(0)=G-1$ and
$c_G^{\mathrm{MaxRL}}(1)=1$.
\end{lemma}

\begin{proof}
Condition on $R=r>0$. Each successful row has advantage $(G-r)/r$ and each
failed row has advantage $-1$. Using the same conditional score identities as
above, the expected unaveraged sum is
\[
  (G-r)\!\left(\nabla_z\log P-\nabla_z\log(1-P)\right)
  =\frac{G-r}{P(1-P)}\nabla_zP.
\]
After the outer factor $1/G$ and averaging over
$R\sim\operatorname{Binomial}(G,P)$,
\[
 \E[(G-R)\1\{R>0\}]
 =G(1-P)-G\Pr(R=0)
 =G\big[(1-P)-(1-P)^G\big].
\]
Substitution gives Equation~\eqref{eq:maxrl-expected-gradient}. Its endpoint
values follow by continuity (or the finite geometric-series identity), which
also gives the stated continuous extension at both boundaries.
\end{proof}

The same finite-series identity gives the potential used below:
\begin{equation}
 \Psi_{\mathrm{MaxRL}}(P)
 :=\int_0^P c_G^{\mathrm{MaxRL}}(s)\,ds
 =\sum_{k=1}^{G-1}\frac{1-(1-P)^k}{k}.
 \label{eq:maxrl-potential-finite-sum}
\end{equation}
In particular, $\Psi_{\mathrm{MaxRL}}(1)=\sum_{k=1}^{G-1}1/k$.
For binary rewards, best@\(k\) equals pass@\(k\), so this is a harmonic
mixture of sampling objectives for $k=1,\ldots,G-1$. The upper limit is
$G-1$, not $G$: the implemented centered estimator drops the entire update
on all-failure groups. Subtracting an unconditional score baseline even on
those groups instead gives the order-$G$ estimator. This distinction concerns
the exact finite-group expectation at the on-policy point; it does not assert
unbiasedness of subsequent clipped or adaptive-optimizer updates.


\subsection{Winner-take-all collapse in the categorical mean flow}
\label{app:theory-collapse}

Consider $\dot z=c_G(P)\nabla_zP$, with any of the smooth positive
coefficients established above: Dr.GRPO, common-length reward-standardized
GRPO, or the specified centered binary MaxRL estimator. This infinitesimal
expected-update model removes
sampling noise and retains independently trained Euclidean logits. The
conditional dynamics reduce to the standard frequency-dependent replicator
system with fitness $f_c(q)=q_c$ \citep{hofbauer1998evolutionary}. The result is
a specialization to the group estimators, not a new general theorem about
replicator dynamics. Related correctness-indifferent collapse analyses include
\citet{sinha2026expected,lochab2026ucpo}.

\begin{theorem}[Standard replicator specialization for binary-reward RL]
\label{thm:grpo-collapse}
Under Assumption~\ref{ass:categorical-mean-flow}, if the conditional correct-mode
distribution $q(0)$ has a unique largest coordinate $c^\star$, then
\[
  P(t)\longrightarrow1,
  \qquad
  q_{c^\star}(t)\longrightarrow1,
  \qquad
  q_c(t)\longrightarrow0\quad(c\ne c^\star).
\]
Thus these categorical mean flows converge to a single correct execution
mode for generic initial conditions, even though every mode in $\mathcal C$
has exactly the same reward.
\end{theorem}

\begin{remark}
Theorem~\ref{thm:grpo-collapse} is asymptotic and internal to (A1)--(A3): it
claims neither finite-time support loss nor convergence of the stochastic
neural optimizers used in our experiments.
\end{remark}

\begin{proof}
The coefficient $c_G$ is smooth and bounded on $[0,1]$, and the softmax
reward gradient is bounded. Thus the logit flow exists for all finite times
and cannot develop infinite logits at finite time.
For a correct mode $c$ and incorrect category $a$,
\[
  \dot z_c=c_G(P)p_c(1-P),
  \qquad
  \dot z_a=-c_G(P)p_aP.
\]
Also $\dot P=c_G(P)\lVert\nabla_zP\rVert_2^2>0$. If $P$ converged to an
interior value, continuity and positivity of $c_G$, together with
\[
 \lVert\nabla_zP\rVert_2^2
 \ge P^2(1-P)^2\left(\frac1m+\frac1n\right),
\]
would keep $\dot P$ bounded away from zero, a contradiction. Hence $P\to1$.

Define the effective optimization time
\[
  \tau(t)=\int_0^t c_G(P(s))P(s)(1-P(s))\,ds.
\]
If $\tau(\infty)<\infty$, every correct and incorrect logit would have finite
total variation, because the absolute velocity of each is at most $\dot\tau$.
All logits would then converge to finite values, contradicting $P\to1$.
Therefore $\tau(t)\to\infty$, leaving infinite effective time for the
correct-simplex dynamics below.

On the correct simplex, changing from $t$ to $\tau$ gives the replicator flow
\begin{equation}
  \frac{dq_c}{d\tau}
  =q_c\left(q_c-\sum_{d\in\mathcal C}q_d^2\right),
  \qquad
  \frac{d}{d\tau}\log\frac{q_c}{q_d}=q_c-q_d.
  \label{eq:grpo-replicator}
\end{equation}
The initially unique maximum remains the unique maximum. For $c\ne c^\star$,
set $\xi_c=q_c/q_{c^\star}$, so $\xi_c(0)<1$ and $\xi_c$ is nonincreasing.
Since $q_{c^\star}\ge1/m$,
\[
 \frac{d\log\xi_c}{d\tau}
 =-q_{c^\star}(1-\xi_c)
 \le-\frac{1-\xi_c(0)}{m}.
\]
Hence $\xi_c(\tau)\le\xi_c(0)\exp[-(1-\xi_c(0))\tau/m]\to0$.
Every competing mode vanishes relative to $c^\star$; normalization then
yields $q_{c^\star}\to1$, establishing the claimed single-mode limit.
\end{proof}

If several modes tie for the largest initial probability, symmetry preserves
the tie; the same ratio argument eliminates strictly smaller coordinates and
leaves the tied maxima uniform. An exactly uniform correct distribution is
therefore a special initial condition. The theorem concerns limiting
concentration: finite softmax logits already give positive probabilities at
every finite time, so finite-time positivity alone is not a retention result.

\subsection{No outcome-binary objective can change which mode wins}
\label{app:theory-objective-inertness}

Theorem~\ref{thm:grpo-collapse} is stated for three estimators, which invites
the reading that a fourth might escape it. It does not. Under
Assumption~\ref{ass:categorical-mean-flow} the entire class of group-centered
objectives built from binary verifier outcomes induces the same conditional
trajectory, and differs only in how fast that trajectory is traversed. The
reason is one line of conditioning: a binary reward cannot tell the estimator
which correct mode it just saw.

\begin{theorem}[Conditional inertness of outcome-binary objectives]
\label{thm:objective-inertness}
Under Assumption~\ref{ass:categorical-mean-flow}, let $a(\cdot,\cdot)$ be any
detached advantage measurable with respect to the group's binary reward
vector, applied as
$\widehat g_a=G^{-1}\sum_{i}a(R_i,R)\nabla_z\log p_{Z_i}$ at the on-policy
point. Then
\begin{equation}
  \E[\widehat g_a]=c_a(P)\,\nabla_zP,
  \qquad
  c_a(P)=\frac1G\left(
  \frac{\E[R\,a(1,R)]}{P}-\frac{\E[(G-R)\,a(0,R)]}{1-P}\right).
  \label{eq:objective-inertness}
\end{equation}
If $c_a>0$ on $(0,1)$, then under the time change
$d\tau=c_a(P)P(1-P)\,dt$ the conditional distribution $q$ obeys
Equation~\eqref{eq:grpo-replicator} exactly, for every such $a$. Dr.GRPO,
reward-standardized GRPO, a leave-one-out baseline, the centered binary MaxRL
estimator, and any reweighting shaped by a pass@$k$ target are points in this
class. The choice among them rescales optimization time along the correctness
axis and leaves both the trajectory in $q$ and the identity of the surviving
mode exactly where they were.
\end{theorem}

\begin{proof}
Condition on the binary reward vector. The score identities used in
Lemma~\ref{lem:grpo-mean} give
$\E[\nabla_z\log p_{Z_i}\mid R_i=1]=\nabla_z\log P$ and
$\E[\nabla_z\log p_{Z_i}\mid R_i=0]=\nabla_z\log(1-P)$; neither depends on
which correct mode or which incorrect category was drawn, which is the only
place the binary reward is used. Since $a$ is measurable with respect to that
vector, summing the $R$ successful and $G-R$ failed rows and taking
expectations gives
\[
  \E[\widehat g_a]
  =\frac1G\Big(\E[R\,a(1,R)]\nabla_z\log P
   +\E[(G-R)\,a(0,R)]\nabla_z\log(1-P)\Big),
\]
and substituting $\nabla_z\log P=\nabla_zP/P$ and
$\nabla_z\log(1-P)=-\nabla_zP/(1-P)$ yields
Equation~\eqref{eq:objective-inertness}. The flow
$\dot z=c_a(P)\nabla_zP$ has correct-coordinate velocities
$c_a(P)p_c(1-P)$, so $\frac{d}{dt}\log(q_c/q_d)=c_a(P)P(1-P)(q_c-q_d)$ and the
stated time change removes $c_a$ entirely, leaving
Equation~\eqref{eq:grpo-replicator}.
\end{proof}


The same binomial regrouping that gave Lemma~\ref{lem:grpo-mean} its
denominator-free form applies to $c_a$. Using $R\binom GR=G\binom{G-1}{R-1}$
and $(G-R)\binom GR=G\binom{G-1}{R}$,
\begin{equation}
 c_a(P)=\sum_{j=0}^{G-1}\binom{G-1}{j}P^j(1-P)^{G-1-j}
 \big[a(1,j+1)-a(0,j)\big],
 \label{eq:inertness-polynomial}
\end{equation}
a polynomial in $P$ of degree at most $G-1$. So whenever the applied advantages
are bounded, $c_a$ is smooth and bounded on the closed interval $[0,1]$, which
is every regularity property that the proof of the collapse theorem drew on.

Equation~\eqref{eq:inertness-polynomial} is a Bernstein expansion, and its
control points have a reading that makes the sign of $c_a$ checkable by
inspection rather than assumed. Write
\begin{equation}
 d_j:=a(1,j+1)-a(0,j),\qquad j=0,\ldots,G-1 .
 \label{eq:leave-one-out-gap}
\end{equation}
Holding the other $G-1$ rows fixed at $j$ successes, $d_j$ is the change in the
advantage a single row receives when that row flips from failure to success:
the leave-one-out gap the objective puts between being right and being wrong.
Since the Bernstein basis is positive and sums to one, $c_a(P)$ lies between
$\min_jd_j$ and $\max_jd_j$ at every $P$, so the sign of the coefficient is
settled by the signs of these $G$ numbers and by nothing else.
For Dr.GRPO every $d_j$ equals $(G-1)/G$, recovering the constant coefficient
of Lemma~\ref{lem:grpo-mean}; for the centered MaxRL estimator $d_0=G-1$ and
$d_j=G/(j+1)$ for $j\ge1$, and summing the basis against them returns
$[1-(1-P)^{G-1}]/P$ as in Lemma~\ref{lem:maxrl-mean}. The consequence is that
the collapse theorem was never about the three estimators it was stated for.

\begin{corollary}[Every outcome-binary objective collapses]
\label{cor:class-collapse}
Under Assumption~\ref{ass:categorical-mean-flow}, let $a$ be bounded and
reward measurable, and suppose it rewards correctness in the weak sense that
its leave-one-out gaps satisfy $d_j\ge0$ for every $j$ with $d_j>0$ for at
least one, which by Equation~\eqref{eq:leave-one-out-gap} is a condition on
$2G$ numbers that can be read off the estimator. Let $q(0)$ have a unique
largest coordinate $c^\star$. Then the conclusions of Theorem~\ref{thm:grpo-collapse}
hold verbatim: $P\to1$, $q_{c^\star}\to1$, and $q_c\to0$ for every other
correct mode. The surviving mode is whichever one led at initialization, and
no admissible choice of the advantage $a$ changes its identity --- only the
rate at which it takes the conditional distribution over.
\end{corollary}

\begin{proof}
By Theorem~\ref{thm:objective-inertness} the flow is
$\dot z=c_a(P)\nabla_zP$, which is the flow of
Theorem~\ref{thm:grpo-collapse} with $c_G$ replaced by $c_a$. That proof uses
the coefficient only through its smoothness and boundedness on $[0,1]$, which
Equation~\eqref{eq:inertness-polynomial} supplies, and its positivity on
$(0,1)$, which the hypothesis now supplies: the Bernstein basis functions are
strictly positive on $(0,1)$, so a nonnegative set of control points with one
strictly positive member gives $c_a(P)\ge\binom{G-1}{j_0}P^{j_0}(1-P)^{G-1-j_0}
d_{j_0}>0$ there. The identity of $c^\star$ is fixed by $q(0)$ alone, because
the replicator equation~\eqref{eq:grpo-replicator} reached after the time
change carries no trace of the advantage at all.
\end{proof}

This closes the escape route rather than leaving it open. An objective outside
the corollary is one with some $d_j<0$: for some group composition it
\emph{reduces} the advantage a response receives when that response becomes
correct. That is not a diversity mechanism, and it does not buy breadth ---
Theorem~\ref{thm:objective-inertness} is proved without any sign condition, so
such an objective still produces $c_a(P)\nabla_zP$, still cannot tell two
correct modes apart, and still moves along the very same orbit in $q$, merely
traversing part of it backwards while it also undoes correctness. The
hypothesis therefore excludes only objectives that escape collapse by failing
at the task, which is why no reweighting of binary outcomes is a route to
breadth.

\begin{remark}[Group size is not a remedy]
\label{rem:group-size-inert}
The group size enters the class only through $c_a$, which
Equation~\eqref{eq:inertness-polynomial} shows is a coefficient on
$\nabla_zP$. The time change of Theorem~\ref{thm:objective-inertness} removes
it, so for every $G\ge2$ the conditional dynamics are the same replicator
equation and Corollary~\ref{cor:class-collapse} applies unchanged: enlarging
the rollout group rescales how fast the flow moves along the correctness axis
and does not alter which correct mode survives it. What a larger group does
change is the starvation bound of
Lemma~\ref{lem:replay-gradient-availability} and the MaxRL coefficient of
Corollary~\ref{cor:maxrl-starvation}, both of which concern discovery near
$P=0$ rather than retention near $p_b=0$. Drawing more samples per prompt is
therefore an answer to the first of those problems and not to the second one.
\end{remark}

\begin{corollary}[What a stronger fresh objective does buy]
\label{cor:maxrl-starvation}
For the coefficients of Lemmas~\ref{lem:grpo-mean} and~\ref{lem:maxrl-mean},
\[
 \frac{c_G^{\mathrm{MaxRL}}(P)}{c_G^{\mathrm{Dr.GRPO}}}
 =\frac{G\big[1-(1-P)^{G-1}\big]}{(G-1)P},
\]
which decreases from $G$ at $P\to0$ to $G/(G-1)$ at $P=1$. MaxRL therefore
multiplies progress along the correctness axis by a factor approaching the
group size in exactly the regime where
Lemma~\ref{lem:replay-gradient-availability} bounds fresh signal to be scarce,
and by at most $G/(G-1)$ once correctness is high.
\end{corollary}

\begin{proof}
Divide $c_G^{\mathrm{MaxRL}}(P)=[1-(1-P)^{G-1}]/P$ by $(G-1)/G$. The quotient
is $c_G^{\mathrm{MaxRL}}$ up to a positive constant and
$c_G^{\mathrm{MaxRL}}$ is decreasing, with endpoint values $G-1$ and $1$ from
Lemma~\ref{lem:maxrl-mean}.
\end{proof}

\subsection{A replay signal when fresh groups supply none}
\label{app:theory-gradient-availability}

\begin{lemma}[Fresh-group starvation]
\label{lem:replay-gradient-availability}
Under Assumption~\ref{ass:categorical-mean-flow} with correctness probability
$P$, draw $G$ fresh categories independently. The probability of a
mixed-reward group is
\[
 h_G(P)=1-P^G-(1-P)^G\le G\min\{P,1-P\}.
\]
Every all-equal group has zero fresh task component for
Equation~\eqref{eq:group-estimator} and for the specified centered MaxRL
estimator, so a nonzero fresh task component requires a mixed group, an event
whose probability is at most $h_G(P)$ at every correctness level.
\end{lemma}

\begin{proof}
The all-correct and all-incorrect events are disjoint and have probabilities
$P^G$ and $(1-P)^G$. A mixed group requires both a correct and an incorrect
sample, so the union bound supplies each of the two stated bounds, $GP$ and
$G(1-P)$, and the minimum of the pair is therefore also a bound.
\end{proof}

\begin{lemma}[Replay ascent persists as fresh signal vanishes]
\label{lem:replay-ascent}
Under Assumption~\ref{ass:categorical-mean-flow}, let
$\mathcal B\subseteq\mathcal C$ be a nonempty fixed verified bank with
normalized positive weights $w_b$, let $R_w=-\sum_b w_b\log p_b$ with
$\rho_w>0$, and extend $w$ by zero outside the bank to obtain $\bar w$. The
deterministic replay ascent vector is
\[
 v^{\rm replay}=-\rho_w\nabla R_w=\rho_w(\bar w-p),
\]
which vanishes only at $p=\bar w$. If $p_b\le w_b/2$ for a banked mode, then
$v^{\rm replay}_b\ge\rho_w w_b/2$, even when $b$ is absent from the fresh
group. For a uniform bank of $k\ge2$ modes and fixed $\rho_w>0$, as
$p\to\mathbf e_b$ at one banked correct mode, $h_G(P)\to0$ whereas
$\|v^{\rm replay}\|_2\to\rho_w\sqrt{1-1/k}>0$.
\end{lemma}

\begin{proof}
For $R_w=-\sum_b w_b\log p_b$, the softmax score identity gives
$\partial R_w/\partial z_a=p_a-\bar w_a$, proving the replay formula and
coordinate bound. Finally,
\[
 \|\bar w-\mathbf e_b\|_2^2=(1-1/k)^2+(k-1)/k^2=1-1/k,
\]
which gives the strictly positive limiting norm for $k\ge2$.
\end{proof}

\begin{remark}
Together the two lemmas say only that a banked target can supply a verified
signal when the fresh task gradient is unavailable. Lemma~\ref{lem:replay-gradient-availability}
bounds an opportunity, not a gradient magnitude, and says nothing about other
loss terms, clipping, or momentum carried across updates. The vector in
Lemma~\ref{lem:replay-ascent} averages the fixed bank, so a stochastic
exemplar draw realizes it only in expectation, and a nonzero logit coordinate
can still be annihilated by a degenerate neural Jacobian. Replay also need not
increase the next mixed-group probability, since $h_G(P)$ is largest at
$P=1/2$.
\end{remark}

\subsection{Fixed-bank replay and qualitative retention}
\label{app:theory-replay}

Let $\mathcal B\subseteq\mathcal C$ be a nonempty fixed bank of $k$ modes that
the policy has produced and the validator has accepted. Uniform verified replay
uses one exemplar per banked mode and, in the categorical reduction, minimizes
\begin{equation}
  R_{\mathcal B}(z)
  =-\frac1k\sum_{b\in\mathcal B}\log p_b.
  \label{eq:verified-replay-loss}
\end{equation}
This is ordinary cross-entropy against the distribution that assigns mass
$1/k$ to each banked mode and zero elsewhere. It is not conditioned on the
bank: decreasing~\eqref{eq:verified-replay-loss} can raise total bank mass as
well as redistribute mass within the bank.

In an averaged continuous-time model, absorb a fixed round-robin replay
frequency and loss scale into an effective dose $\rho>0$. This assumes a
fixed eligible-bank set and constant relative weighting; it does not identify
finite alternating updates with simultaneous gradient flow. The prompt-isolated categorical mean flow for any of the three fresh
binary-reward objectives above plus verified replay is
\begin{equation}
  \dot z
  =c_G(P)\nabla_zP-\rho\nabla_zR_{\mathcal B}(z).
  \label{eq:replay-mean-flow}
\end{equation}
Let $\Psi(P)=\int_0^P c_G(s)\,ds$. For the finite-group estimators in Lemmas~\ref{lem:grpo-mean}
and~\ref{lem:maxrl-mean}, $c_G$ has a finite positive extension to $[0,1]$;
in particular, Dr.GRPO has $c_G=(G-1)/G$ and MaxRL has
$c_G=[1-(1-P)^{G-1}]/P$. Nothing below uses more than that. By
Theorem~\ref{thm:objective-inertness} and
Equation~\eqref{eq:inertness-polynomial}, every bounded reward-measurable
objective contributes $c_a(P)\nabla_zP$ with $c_a$ a polynomial, so it admits
the same $\Psi_a$ and the retention argument runs unchanged: what follows is a
statement about verified replay added to \emph{any} outcome-binary fresh
objective, of which Re:Dr and Re:Max are the two this paper trains. Define
\begin{equation}
  F_{\mathcal B}(z)=\rho R_{\mathcal B}(z)-\Psi(P(z)).
  \label{eq:replay-potential}
\end{equation}
Then~\eqref{eq:replay-mean-flow} is exactly gradient flow
$\dot z=-\nabla_zF_{\mathcal B}$.

\begin{theorem}[Verified replay prevents banked-mode extinction]
\label{thm:replay-retention}
Suppose logits are finite at time $T$ and thereafter the nonempty verified
bank $\mathcal B$ is fixed, fits within replay capacity, and receives a
fixed positive effective replay dose $\rho$, the fresh objective being any
bounded reward-measurable one in the sense of
Theorem~\ref{thm:objective-inertness}. Under~\eqref{eq:replay-mean-flow}, define
\[
  C_T
  =R_{\mathcal B}(z(T))
   +\frac{\Psi(1)-\Psi(P(T))}{\rho}.
\]
Then for every $b\in\mathcal B$ and every $t\ge T$,
\begin{equation}
  p_b(t)\ge \exp(-kC_T)>0.
  \label{eq:replay-retention-bound}
\end{equation}
Thus each banked probability is uniformly bounded away from zero for all
$t\ge T$, including asymptotically, even when the bank does not cover every
correct mode. This is stronger than finite-time positivity, which already
follows from finite softmax logits without replay.
\end{theorem}

\begin{proof}
The smooth vector field has bounded velocity because $c_G$,
$\nabla_zP$, and $\nabla_zR_{\mathcal B}$ are bounded, so finite initial
logits give a flow for every finite $t\ge T$. Along this flow,
\[
  \frac{d}{dt}F_{\mathcal B}(z(t))
  =-\lVert\nabla_zF_{\mathcal B}(z(t))\rVert_2^2\le0.
\]
Since $P(t)\le1$ and $\Psi$ is increasing,
\[
  R_{\mathcal B}(z(t))
  \le R_{\mathcal B}(z(T))
     +\frac{\Psi(1)-\Psi(P(T))}{\rho}
  =C_T.
\]
Every $-\log p_b$ is nonnegative and their sum is
$kR_{\mathcal B}$. Therefore
$-\log p_b(t)\le kC_T$ for every banked mode, which is exactly
Equation~\eqref{eq:replay-retention-bound} for every time after $T$.
\end{proof}


The floor is qualitative and can be numerically vacuous. For example, at
$k=16$ even a hypothetical $C_T=160$ gives $e^{-2560}$, far below any useful
evaluation probability. The argument establishes a uniform-in-time lower
bound in the stated ideal flow, not practical visibility or a prediction for
the experimental replay loss. The effective dose also depends on relative
fresh/replay frequencies, not the replay coefficient alone.

A common nonnegative, locally integrable rate $\eta(t)$ multiplying both
terms of~\eqref{eq:replay-mean-flow} preserves the energy bound, since
$\dot F_{\mathcal B}=-\eta(t)\|\nabla F_{\mathcal B}\|_2^2\le0$.
The convergence statements below additionally require
$\int_T^\infty\eta(t)\,dt=\infty$ and then follow by reparameterizing
optimization time. Infinitely many replay visits alone establish neither
this condition nor a fixed replay-to-fresh weighting. Different time-varying
relative weights are outside the fixed-potential argument.

\paragraph{Full coverage and the categorical target.}
\begin{corollary}[Conditional full-coverage categorical limit]
\label{cor:replay-no-collapse}
Assume the conditions of Theorem~\ref{thm:replay-retention} and that by time
$T$ the bank covers the full correct support, $\mathcal B=\mathcal C$. Let
$u_c=1/m$ on $\mathcal C$. Then the categorical mean flow converges to
\[
  P(t)\longrightarrow1,
  \qquad q(t)\longrightarrow u.
\]
Consequently
\[
  \min_{c\in\mathcal C}q_c(t)\longrightarrow\frac1m,
  \qquad
  \mathcal H(q(t))\longrightarrow\log m,
\]
and conditional correct-mode collapse is impossible.
\end{corollary}

\begin{proof}
For full coverage,
\begin{equation}
 R_{\mathcal C}=\log m-\log P+\mathrm{KL}(u\|q).
 \label{eq:replay-mass-uniformity-decomposition}
\end{equation}
The potential is bounded below by $-\Psi(1)$, so its descent gives
$\int_T^\infty\|\nabla F_{\mathcal C}\|_2^2dt<\infty$. The finite-group
$c_G$ is smooth on $[0,1]$, and softmax derivatives and the cross-entropy
gradient are bounded. Thus the Hessian of $F_{\mathcal C}$ and the velocity
are bounded; the gradient along the trajectory is uniformly continuous.
Square integrability then implies $\nabla F_{\mathcal C}\to0$: recurring
gradients of fixed positive size would otherwise contribute energy over
intervals of a uniform minimum length.

Every limiting probability vector therefore satisfies the correct-coordinate
equations
\[
 0=\rho(p_c-1/m)-c_G(P)p_c(1-P).
\]
Summing them gives $0=-(1-P)(\rho+c_G(P)P)$, hence $P=1$. Each correct
coordinate is then $1/m$ and every incorrect coordinate is zero. Compactness
of the probability simplex and uniqueness of this limit establish convergence
of the probability vector.
\end{proof}

\paragraph{Positive weights and the implemented length normalization.}
Uniformity of the target is separate from protection against extinction.
For fixed weights $w_b>0$ with $\sum_{b\in\mathcal B}w_b=1$, define
$R_w=-\sum_b w_b\log p_b$ and use replay coefficient $\rho_w>0$.
The same potential argument gives, for all $t\ge T$,
\[
 p_b(t)\ge\exp(-C_{T,w}/w_b),\qquad
 C_{T,w}=R_w(z(T))+\frac{\Psi(1)-\Psi(P(T))}{\rho_w}.
\]
If $\mathcal B=\mathcal C$, then
$R_w=H(w)-\log P+\mathrm{KL}(w\|q)$, and the independent-logit limit is
$P\to1$, $q\to w$. The preceding stationarity proof applies with $1/m$
replaced by $w_c$: summing the correct-coordinate equations forces $P=1$,
then each equation forces $p_c=w_c$.

A partial bank has a stronger limitation in this same ideal flow. For any
fixed nonempty $\mathcal B\subseteq\mathcal C$ and fixed positive weights
summing to one, extend $w$ by zero to $\bar w$ on all categories. Then
$p(t)\to\bar w$, so even correct modes outside the frozen bank vanish
asymptotically. For $F=\rho_wR_w-\Psi(P)$, the preceding energy and
bounded-Hessian argument still gives $\nabla F\to0$; neither argument
requires full bank coverage.
At every subsequential probability limit, the correct-coordinate equations are
\[
 0=\rho_w(p_c-\bar w_c)-c_G(P)p_c(1-P).
\]
Summing over correct categories gives
$0=-(1-P)(\rho_w+c_G(P)P)$, hence $P=1$; each equation then gives
$p_c=\bar w_c$, and all incorrect mass is zero. Compactness and uniqueness
of this limiting probability vector prove convergence. Thus protecting a
fixed set of banked modes need not preserve unbanked correct alternatives.
This conclusion concerns the fixed-bank categorical flow, not changing-bank
or stochastic neural training.

This distinction matters for the implemented mean-token exemplar loss.
Even if each key were represented by exactly one categorical exemplar of
fixed length $L_b$, its loss would be $A R_w$, where
\[
 A=\frac1k\sum_{b\in\mathcal B}\frac1{L_b},\qquad
 w_b=\frac{L_b^{-1}}{\sum_{j\in\mathcal B}L_j^{-1}}.
\]
Its effective categorical replay coefficient would be $\rho_w=\rho A$.
The full-coverage limit would therefore weight shorter exemplars more, and
would be uniform only for equal lengths. For an actual language model,
$\pi_\theta(e_b\mid x)$ is only one contribution to the total key probability.
Equal averaging of per-token exemplar scores does not imply uniform learned
key probabilities; the categorical uniformity corollary is not that claim.

Full support is an additional premise, not a consequence of replay. In
Python factors, evaluation prompts admit \MDvmPythonRange{} correct execution
modes (mean \MDvmPythonMean{}), while the bank capacity is 16. The
full-coverage corollary therefore cannot supply a domain-wide guarantee there. The retention theorem protects only
admitted banked modes; it supplies no probability floor for an unbanked key.

\subsection{Complete responses and finite evaluation visibility}
\label{app:theory-exemplar-bridge}

The learner scores a particular response string, while evaluation groups
strings by execution key. Event containment provides a conditional bridge:
one complete verified response is one way of producing its key. This bridge
also makes the loss's length factor explicit. Here
$p_\theta(b\mid x)=\Pr_{Y\sim\pi_\theta(\cdot\mid x)}[V(Y,s_x)=b]$
is the unconditional probability of producing key $b$, rather than the
probability conditional on correctness.

\begin{lemma}[Retention from a joint exemplar potential]
\label{lem:shared-exemplar-retention}
Fix finitely many verified complete-response exemplars $(x_j,e_j)$, each
realizing key $b_j$, with lengths $L_j>0$ and fixed weights $\alpha_j>0$.
Let
\[
 F(\theta)=\sum_j\alpha_j
 \frac{-\log\pi_\theta(e_j\mid x_j)}{L_j}-J(\theta),
 \qquad J(\theta)\le J_{\max}<\infty.
\]
Suppose $F(\theta(T))$ is finite and
$F(\theta(t))\le F(\theta(T))$ for every $t\ge T$. With
$D=F(\theta(T))+J_{\max}$, every protected exemplar satisfies
\[
 p_\theta(b_j\mid x_j)\ge\pi_\theta(e_j\mid x_j)
 \ge\exp\!\left(-\frac{L_jD}{\alpha_j}\right)>0
 \qquad(t\ge T).
\]
\end{lemma}
\begin{proof}
The weighted exemplar loss equals $F(\theta(t))+J(\theta(t))\le D$.
Each summand is nonnegative, so
$-\log\pi_\theta(e_j\mid x_j)\le L_jD/\alpha_j$.
The complete-response event is a subset of the event producing key $b_j$,
which proves the stated lower bound on every protected key's probability.
\end{proof}

Exact descent of a smooth joint potential is sufficient for the energy
inequality, even with shared parameters and multiple prompts. For example,
the ideal categorical model permits
$J(\theta)=\sum_x\nu_x\Psi_x(P_x(\theta))$ with fixed nonnegative prompt
weights and $J_{\max}=\sum_x\nu_x\Psi_x(1)$. The lemma does not establish
this descent condition for stochastic, clipped, alternating AdamW updates,
and it neither implies uniform neural key probabilities nor requires them.
Its bound applies only to the listed exemplar--prompt pairs. Shared
parameters do not extend it to held-out prompts; transfer to those prompts
remains an empirical question.

The scored probability must describe the complete response, including its
termination event or a fixed deterministic horizon, under the same prompt,
temperature, masks, truncation rules, and verifier as the sampling law being
bounded. A prefix score or a score under another decoding law does not
supply this bound without a separate event or likelihood comparison.
For a complete response of length $L$ and mean-token log score $s$, its
probability is $e^{Ls}$, not $e^s$. A change in this exemplar's likelihood
need not equal the change in its whole key's probability.

\paragraph{Retention and finite evaluation visibility.}
At a fixed checkpoint, suppose $k$ distinct banked keys for one prompt have
probability floors $p_b\ge\delta_b>0$. For $K$ independent draws from that
same policy, the banked distinct count $D_{\mathcal B,K}$ satisfies
\[
 \begin{aligned}
 \mathbb E[D_{\mathcal B,K}]
 &\ge\sum_{b\in\mathcal B}\left[1-(1-\delta_b)^K\right],\\
 \Pr(\text{some banked key is unseen})
 &\le\sum_{b\in\mathcal B}(1-\delta_b)^K
 \le k e^{-K\delta_{\min}},
 \end{aligned}
\]
where $\delta_{\min}=\min_b\delta_b$. The first statement follows from
indicator expectations and the second from the union bound. Thus
$K\ge\log(k/\varepsilon)/\delta_{\min}$ suffices to observe all banked keys
with probability at least $1-\varepsilon$, for $0<\varepsilon<1$.
The energy floors can make this budget impractically large; asymptotic
retention does not imply visibility in an eight-draw evaluation.

Conversely, not observing a mode does not establish extinction. At a fixed
checkpoint, zero sightings of a prespecified key in $N$ independent draws
give a one-sided $(1-\alpha)$ binomial upper bound
$1-\alpha^{1/N}$ and lower bound zero \citep{clopper1934binomial}. For
$N=32$ and $\alpha=.05$, the upper bound is approximately $8.94\%$.
These are single-key bounds; simultaneous statements need error allocation.
Counts from changing checkpoints cannot be pooled to certify a current
policy, and teacher-forced scores are not independent Bernoulli observations.

\subsection{Reference KL retains the reference, not the discovery}
\label{app:theory-kl}

Assumption~\ref{ass:categorical-mean-flow}(A4) sets the reference-KL
coefficient to zero, which is what our runs train at
(Table~\ref{tab:run-contract}, $\beta_{\mathrm{KL}}=0$). The collapse theorem
therefore says nothing about the regularizer most RLVR recipes carry, and a
reader is entitled to ask whether that regularizer already does the work the
bank is introduced for. This subsection answers in the idealized flow. The
answer is not that reference KL fails to prevent winner-take-all --- for every
$\beta>0$ it prevents it --- but that the allocation it holds the policy at is
the reference model's own, that it holds it there with a restoring force which
vanishes as a mode vanishes, and that the two properties have the same cause:
a bounded penalty cannot be a barrier.

Replace (A4) by
\begin{itemize}
\item[(A4$'$)] \emph{Reference KL.} A fixed reference $\mu$ with full support
  on the category set $\mathcal A$ enters with a positive coefficient $\beta$,
  and no entropy term is present alongside it.
\end{itemize}
The low-variance estimator standard in GRPO implementations is unbiased for
$\mathrm{KL}(p\,\|\,\mu)$ under samples drawn from $p$
\citep{shao2024deepseekmath}, so the reverse direction is the one to analyze;
the forward direction is treated at the end of this subsection and behaves
differently in exactly the way that matters. Write
$D(z)=\mathrm{KL}(p(z)\,\|\,\mu)$. The prompt-isolated categorical mean flow is
\begin{equation}
  \dot z=c_G(P)\nabla_zP-\beta\nabla_zD(z),
  \qquad
  \nabla_{z_b}D=p_b\big(\log(p_b/\mu_b)-D\big),
  \label{eq:kl-mean-flow}
\end{equation}
which is exact gradient flow of the potential
$F_\beta(z)=-\Psi(P(z))+\beta D(z)$, built exactly as
Equation~\eqref{eq:replay-potential} was built for verified replay.

\begin{lemma}[A bounded penalty supplies no probability floor]
\label{lem:kl-no-barrier}
For a full-support reference, $D\le\log(1/\min_a\mu_a)<\infty$ everywhere on
the simplex, whereas $R_{\mathcal B}\to\infty$ whenever a banked coordinate
$p_b\to0$. Consequently the energy argument of
Theorem~\ref{thm:replay-retention}, applied to $F_\beta$, returns only
$p_b\ge0$: for every $\beta>0$ it yields no uniform-in-time positive lower
bound on any coordinate.
\end{lemma}

\begin{proof}
$D=\sum_ap_a\log p_a-\sum_ap_a\log\mu_a\le0+\log(1/\min_a\mu_a)$. Descent of
$F_\beta$ bounds $\beta D(z(t))$ by $F_\beta(z(T))+\Psi(1)$, but $D$ stays
finite as $p_b\to0$, so this bounds no individual $-\log p_b$. The
corresponding step in Theorem~\ref{thm:replay-retention} works only because
$kR_{\mathcal B}$ dominates each $-\log p_b$ and is itself unbounded on the
simplex, which is exactly the property $D$ lacks.
\end{proof}

\begin{proposition}[The retained conditional is the reference's]
\label{prop:kl-stationary}
Under Assumption~\ref{ass:categorical-mean-flow} with (A4$'$), the flow
\eqref{eq:kl-mean-flow} has exactly one stationary point of full support,
\[
  p^\star_a\propto\mu_ae^{\lambda R_a},
  \qquad
  \lambda=c_G(P^\star)/\beta,
  \qquad
  R_a=\1[a\in\mathcal C],
\]
with $\lambda=(G-1)/(G\beta)$ in closed form for the Dr.GRPO coefficient, and
a unique $\lambda$ for the MaxRL coefficient. Consequently
\[
  q^\star=\mu|_{\mathcal C},
  \qquad
  \mathcal H(q^\star)=\mathcal H(\mu|_{\mathcal C}),
  \qquad
  \operatorname{logit}P^\star=\operatorname{logit}\mu(\mathcal C)+\lambda .
\]
This stationary conditional does not depend on $\beta$: the coefficient sets
the correctness level there and leaves the allocation among correct modes at
the reference's.
\end{proposition}

\begin{proof}
Stationarity of~\eqref{eq:kl-mean-flow} in a coordinate with $p_b>0$ reads
$c_G(P)(R_b-P)=\beta(\log(p_b/\mu_b)-D)$, so
$\log(p_b/\mu_b)=D+(c_G(P)/\beta)(R_b-P)$ and the $b$-independent terms are
absorbed by normalization, giving the stated form with $\lambda=c_G(P)/\beta$.
Conversely each $\lambda\ge0$ determines
$P(\lambda)=\mu(\mathcal C)e^\lambda/(\mu(\mathcal C)e^\lambda+1-\mu(\mathcal C))$,
strictly increasing from $\mu(\mathcal C)$ to $1$. The consistency requirement
$\lambda=c_G(P(\lambda))/\beta$ has a solution by the intermediate value
theorem, since the left side increases from $0$ without bound while the right
side is positive and bounded; it is unique whenever $c_G$ is non-increasing,
which covers the constant Dr.GRPO coefficient and
$c_G^{\mathrm{MaxRL}}(P)=[1-(1-P)^{G-1}]/P$, decreasing from $G-1$ to $1$.
Two correct modes give $\log(p_c/\mu_c)=\log(p_d/\mu_d)$, hence
$q^\star\propto\mu$ on $\mathcal C$, which is $\beta$-free. The logit identity
is $P^\star/(1-P^\star)=[\mu(\mathcal C)/(1-\mu(\mathcal C))]e^\lambda$.
\end{proof}

Read against Corollary~\ref{cor:replay-no-collapse}, this is the substantive
difference. Full-coverage replay sends $q\to u$ and
$\mathcal H(q)\to\log m$ for \emph{every} $\rho>0$; the best a reference KL
offers, at every $\beta>0$, is to hold $q$ at $\mu|_{\mathcal C}$. A reference that is already
concentrated among correct modes is returned unchanged, and
Section~\ref{sec:results-collapse} measures exactly that concentration at pass~0.
Reference KL is an anchor and not a target: it can hold breadth where it was
and cannot raise it above what the reference had. The proposition is stated for
an arbitrary full-support $\mu$, so it also answers the natural repair. One
could anchor to a broader reference and inherit its allocation instead --- but
$q^\star=\mu|_{\mathcal C}$ says the result is only ever as broad as the model
supplying it, so the repair presupposes a policy that already allocates its
correct mass the way the training was supposed to make it allocate. A bank is
not circular in this way: it is assembled from the policy's own verified
successes as they are found, and Corollary~\ref{cor:replay-no-collapse} sends
$q$ to a distribution no available model had to exhibit first.

\begin{remark}
The correctness statement in Proposition~\ref{prop:kl-stationary} is structural
rather than numerical, in the sense of the remark following
Theorem~\ref{thm:replay-retention}. With the Dr.GRPO coefficient at $G=16$ and
any $\beta\le.1$, $\lambda\ge9.4$ and $1-P^\star$ is at most of order
$10^{-4}$: the flow's correctness ceiling sits far below a measurable deficit.
What survives at practical $\beta$ is the conditional statement, which is
$\beta$-free, together with the rate statement established in the corollary
that follows it below.
\end{remark}

\begin{lemma}[Reference KL has no probability-independent boundary force]
\label{lem:kl-boundary-force}
Under (A4$'$) the logit velocity contributed by the reference-KL term is
\[
  \dot z^{\mathrm{KL}}_b=\beta p_b\big(\log(\mu_b/p_b)+D\big)
  \xrightarrow[\,p_b\to0\,]{}0,
  \qquad
  \big|\dot z^{\mathrm{KL}}_b\big|
  \le\beta\Big(\frac{\mu_b}{e}+\log\frac1{\min_a\mu_a}\Big),
\]
whereas the verified-replay term contributes $\rho(w_b-p_b)\to\rho w_b>0$ at
the same limit. Reference KL therefore belongs to the class described by
Lemma~\ref{lem:sampled-score-bound}, with the difference that this is a
property of the exactly evaluated penalty rather than of a sampled estimator.
\end{lemma}

\begin{proof}
The displayed velocity is the softmax score identity used above. For the
bound, $\sup_{0<s\le1}s\log(\mu_b/s)=\mu_b/e$ at $s=\mu_b/e$, the same
expression is bounded below by $\log\mu_b$ at $s=1$, and $0\le D\le
\log(1/\min_a\mu_a)$ by Lemma~\ref{lem:kl-no-barrier}. The replay expression is
Lemma~\ref{lem:replay-ascent}.
\end{proof}

\begin{corollary}[Recovery from depth $s$: $\Theta(1/s)$ against $\Theta(\log(1/s))$]
\label{cor:kl-recovery-time}
Fix correct modes $b$ and $c^\star$ in a state with $p_b=s$ small and
$p_{c^\star}=1-O(s)$. The normalization drift cancels in a ratio,
$\frac{d}{dt}\log(p_b/p_{c^\star})=\dot z_b-\dot z_{c^\star}$, so under
(A4$'$)
\[
  \frac{d}{dt}\log\frac{p_b}{p_{c^\star}}
  =\beta s\log\frac{\mu_b}{s\mu_{c^\star}}-c_G(P)(1-P)+O\!\big(\beta s\log(1/s)\big),
\]
whereas under the fixed-bank replay term with $b,c^\star\in\mathcal B$ and
uniform weights it is exactly $\rho(p_{c^\star}-p_b)=\rho(1-O(s))$. Hence, for
$s$ small enough that the leading term dominates the correctness drag, the
time to raise $p_b$ from $s$ to a fixed $\delta$ is
\[
  T_{\mathrm{KL}}(s)=\frac{1+o(1)}{\beta s\log(1/s)},
  \qquad
  T_{\mathrm{replay}}(s)=\frac{\log(\delta/s)}{\rho}\,(1+o(1)),
\]
and $T_{\mathrm{KL}}/T_{\mathrm{replay}}\to\infty$ as $s\to0$.
\end{corollary}

\begin{proof}
For the KL term, $D=\log(1/\mu_{c^\star})+O(s\log(1/s))$ at such a state, so
$p_{c^\star}(\log(\mu_{c^\star}/p_{c^\star})+D)=O(s\log(1/s))$ while
$p_b(\log(\mu_b/p_b)+D)=s\log(\mu_b/(s\mu_{c^\star}))+O(s\log(1/s))$; the
correctness term contributes $c_G(P)(1-P)(p_b-p_{c^\star})$. For the replay
term, $\dot z_b-\dot z_{c^\star}=\rho[(w_b-w_{c^\star})-(p_b-p_{c^\star})]$,
which is $\rho(p_{c^\star}-p_b)$ for uniform weights. Since
$p_b=s(1+O(s))$ in both cases, separating variables gives
$dt=dp_b/(\beta p_b^2\log(1/p_b))$ and $dt=dp_b/(\rho p_b)$ respectively;
integrating each from $s$ to $\delta$ under $s\ll\delta$ then yields the two
displayed recovery times.
\end{proof}

Table~\ref{tab:kl-recovery} integrates the two flows numerically and reaches
both rates at depths a run could plausibly visit. The reference there is
uniform, which is the case most favorable to reference KL: by
Proposition~\ref{prop:kl-stationary} its stationary conditional then sits
exactly on the replay target, so the two regularizers aim at the same
distribution and only the rate separates them. Recovery time grows by a factor
of \MDklDecadeKL{} per decade of depth under reference KL, at both
coefficients, and by a constant additive amount under replay that reproduces
the predicted $\log(10)/\rho$ to four significant figures at both doses.

\begin{table}[t]
  \centering
  \caption{\textbf{Time for a depleted correct mode to return to
  $p_b=\MDklTarget{}$, by the depth it starts from.} Numerical integration of
  the categorical mean flow with $m=\MDklModes{}$ correct modes, $G=16$, and a
  uniform reference, in the flow's own time units. Reference KL multiplies its
  recovery time by \MDklDecadeKL{} per decade of depth, matching the
  $\Theta(1/s)$ of Corollary~\ref{cor:kl-recovery-time}; verified replay adds
  a constant per decade, matching $\Theta(\log(1/s))$. The coefficients scale
  the constants and not the exponents.}
  \label{tab:kl-recovery}
  \small
  \setlength{\tabcolsep}{6pt}
  \begin{tabular}{@{}lrrrr@{}}
    \toprule
    \textbf{Regularizer} & $s=10^{-3}$ & $10^{-4}$ & $10^{-5}$ & $10^{-6}$ \\
    \midrule
    \input{results/kl_replay_recovery_table_body.tex}
  \end{tabular}
\end{table}

\paragraph{The direction of the divergence decides whether it is a barrier.}
Forward KL to the same reference is a different object:
$\mathrm{KL}(\mu\,\|\,p)=-\sum_a\mu_a\log p_a-\mathcal H(\mu)$ is exactly the
weighted replay loss $R_w$ of
App.~\ref{app:theory-replay} with $w=\mu$, so it is unbounded at the
boundary and does supply the floor $p_a\ge\exp(-C_{T,\mu}/\mu_a)$. Three
differences remain, and they are the design of the bank. Its weights are the
reference's own probabilities, so a mode the base model rarely emits receives
a correspondingly weak floor, where uniform replay weights every admitted key
alike. Its support is all of $\mathcal A$, so it protects incorrect categories
on the same footing as correct ones, where the bank is verifier-filtered.
And it cannot be estimated from on-policy samples without importance weights or
draws from $\mu$, where the bank is a stored sample of the policy's own
verified history. Verified replay is the barrier construction with the target
replaced by a policy-generated, validator-accepted, uniformly weighted one.

\paragraph{A regularizer whose target carries incorrect mass cannot reach
$P=1$.} The structural reason replay costs no correctness is that
$\mathcal B\subseteq\mathcal C$. For any fixed target $\nu$ used in a
cross-entropy or KL penalty, the stationary equations summed over
$\mathcal C$ give $P^\star=1$ when $\nu(\mathcal N)=0$ and $P^\star<1$
strictly when $\nu(\mathcal N)>0$, as in
Proposition~\ref{prop:kl-stationary}. A full-support reference necessarily has
$\mu(\mathcal N)>0$. This is the same accounting as the remark above and
inherits its numerical weakness; it is recorded because it identifies which
property of the target, rather than which coefficient, produces the
difference.

\paragraph{What this does not establish.} (A1)--(A3) still bind, so these are
statements about independently optimized logits under a mean flow, not about
AdamW on a language model; the caveats of
App.~\ref{app:theory-collapse} apply unchanged. A $\beta$ schedule, an
annealed or adaptively controlled coefficient, and a reference that is itself
updated are all outside the fixed-potential argument. The sampled $k_3$
estimator is weaker still than the exactly evaluated penalty analyzed here,
since a mode absent from the rollout group contributes nothing to it at all,
which places it under Lemma~\ref{lem:sampled-score-bound} as well; we have not
run it, and $\beta_{\mathrm{KL}}=0$ in every arm we report, so nothing here is
an empirical comparison. Reference KL has its own guarantees, for staying near
a reference and for the regularized control problem
\citep{geist2019regmdp,vieillard2020kl}, and those are not what
Corollary~\ref{cor:kl-recovery-time} contradicts. The claim is the narrow one
the two lemmas support jointly: among correct modes that have already become
rare, a bounded penalty restores at a rate that vanishes with the mode, and an
unbounded one does not.

\begin{corollary}[The three limits, in the reported breadth metric]
\label{cor:pmd-limits}
The breadth axis this paper reports is
$\pmd{}=1-C(q)=1-\sum_{c\in\mathcal C}q_c^2$. In the categorical mean flow, a
bounded outcome-binary objective alone sends $\pmd{}\to0$
(Corollary~\ref{cor:class-collapse}), and adding full-coverage uniform verified
replay sends $\pmd{}\to1-1/m$ for every dose $\rho>0$
(Corollary~\ref{cor:replay-no-collapse}). Adding a reference KL instead leaves
the flow with a single interior stationary point, at which
\[
 \pmd{}=1-\sum_{c\in\mathcal C}
   \left(\frac{\mu_c}{\mu(\mathcal C)}\right)^{\!2}
 \qquad\text{for every }\beta>0,
\]
the frozen reference's own success-conditional breadth
(Proposition~\ref{prop:kl-stationary}).
\end{corollary}

\begin{proof}
Substitute $\mathbf e_{c^\star}$, the uniform $u$ on $\mathcal C$, and
$\mu|_{\mathcal C}$ into $1-\sum_cq_c^2$, using continuity of $C$ on the
simplex.
\end{proof}

The difference in what the three statements assert is itself a result and not
an artifact of how hard we tried. The first two are limits, reached because
each flow descends a potential whose stationary set the trajectory provably
attains. The third is only the location of a stationary point, because
Lemma~\ref{lem:kl-no-barrier} leaves the reference-KL flow with nothing to
keep it off the boundary: a vertex is a degenerate stationary point of that
flow too, and a bounded penalty supplies no argument excluding it. Reference KL
therefore cannot be shown, by this route, to reach even its own fixed point,
whereas verified replay reaches its target for every positive dose. We report
that asymmetry rather than papering over it, and it is the sharper form of the
claim that a bounded penalty is the wrong instrument for a boundary problem.

On the axis the tables report, then, these are not three degrees of the same
effect. One drives breadth to its floor, one to the ceiling the bank's coverage
allows, and one returns the value the base model began with --- a quantity
Section~\ref{sec:results-collapse} measures at pass~0 rather than a free
parameter, which is what makes the middle and last rows distinguishable in
practice and not only in the limit.

\paragraph{Two boundaries, and which instrument reaches each.}
The results in this appendix separate along the two boundaries of the simplex
that training actually visits, and they are reached by different means.

At $P\to0$ no verified success is available anywhere, and
Lemma~\ref{lem:replay-gradient-availability} bounds the probability of a
mixed group --- the only kind that carries fresh signal --- by $GP$. This is a
\emph{discovery} boundary, and a fresh objective is the instrument for it:
Corollary~\ref{cor:maxrl-starvation} gives MaxRL a factor approaching $G$
there, and Equation~\eqref{eq:maxrl-potential-finite-sum} identifies what it
is maximizing as a harmonic mixture of pass@$k$ for $k=1,\ldots,G-1$.

At $p_b\to0$ for an already-discovered correct mode, the policy is losing
something it found. This is a \emph{retention} boundary, and
Theorem~\ref{thm:objective-inertness} says no member of the fresh class
reaches it, while Lemma~\ref{lem:kl-boundary-force} says a bounded penalty
does not either. Only an unbounded penalty on that coordinate supplies a force
that does not vanish with it, which is Theorem~\ref{thm:replay-retention}.

The two instruments are therefore not substitutes and not alternatives to each
other. They also compose without interfering: adding verified replay to the
MaxRL coefficient gives the single potential
$F=\rho R_{\mathcal B}-\Psi_{\mathrm{MaxRL}}(P)$, whose descent carries the
$\Theta(G)$ starvation factor of Corollary~\ref{cor:maxrl-starvation} and the
probability floor of Theorem~\ref{thm:replay-retention} at once, because the
extension of that theorem to the whole outcome-binary class asked nothing of
the fresh objective beyond boundedness. That combination is the method this
paper reports (Section~\ref{sec:results-maxrl}). Their prices differ in the same direction as
their roles. A fresh-objective change is arithmetic on advantages that were
computed anyway and costs nothing. Verified replay costs a bank with no
parameters and, by Equation~\eqref{eq:replay-flop-ratio}, a per-update
overhead of $3\E|\mathcal B|/G$ against a reference-KL forward, measured below
one on four of our five domains (App.~\ref{app:replay-cost}). Reference KL
costs a resident copy of the policy, measured at \MDcostStateMin{} to
\MDcostStateMax{} times the bank it would replace, in exchange for a
conditional limit it inherits from the reference rather than chooses
(Proposition~\ref{prop:kl-stationary}). That ordering is an accounting
statement about these three additions, not a ranking of published methods on
their own benchmarks.

\subsection{Scope relative to entropy, KL, admission, and neural optimization}
\label{app:theory-entropy}
\label{app:theory-entropy-comparison}

\paragraph{Dependence on optimization geometry.}
For a smooth parameterization $z=z(\theta)$ with Jacobian
$J_\theta=\partial z/\partial\theta$, Euclidean parameter flow instead
induces $\dot z=c_G(P)J_\theta J_\theta^\top\nabla_zP$. The independently
trained logit theorem uses $J_\theta J_\theta^\top=I$ and does not automatically
extend to arbitrary shared parameters or another optimization metric. Equal
binary reward alone does not force this winner selection: the categorical
natural-gradient reward flow can be represented by
$\dot z_a=c_G(P)(R_a-P)$, which gives identical velocities to correct logits
and preserves their ratios. Here $R_a=\1[a\in\mathcal C]$.

This caveat is specific to the collapse conclusion, and
Theorem~\ref{thm:objective-inertness} does not inherit it. For any fixed
preconditioner $M\succeq0$ the flow is $\dot z=c_a(P)M\nabla_zP$, so two
members of the class differ by a positive state-dependent scalar and trace the
same orbit at different speeds, whatever $M$ is. Whichever endpoint the
geometry selects --- a single surviving mode under the independent-logit
metric, preserved ratios under the natural-gradient one --- it selects the same
one for every outcome-binary objective. So the conclusion that a fresh
objective is not the instrument for this problem does not rest on the choice of
geometry; only the identification of the endpoint as winner-take-all does.
Parameter count alone does not specify the optimization geometry. The empirical
scale
comparisons also change pretrained models and learning-rate schedules; they
neither measure the geometry nor identify the cause of cross-scale differences.
The theorem supplies a mechanism to test in controlled training; it cannot
identify the training history or cause of concentration in a hosted model.

A bounded advantage based on fresh samples has a different boundary behavior.
The next bound includes group-dependent detached scores and distinguishes
sampled updates from gradients of an exactly evaluated entropy functional.

\begin{lemma}[Bounded sampled scores have no probability-independent boundary force]
\label{lem:sampled-score-bound}
Under Assumption~\ref{ass:categorical-mean-flow}, condition on the training history, draw
$Z_1,\ldots,Z_G$ independently from $p$, and let $a_i$ be any detached,
possibly group-dependent advantages with $|a_i|\le B$. To interpret the
bound uniformly as $p_b\to0$, assume a finite $B$ independent of that limit.
At the on-policy point define
\[
 \widehat g=\frac1G\sum_{i=1}^G a_i(\mathbf e_{Z_i}-p).
\]
For every category $b$,
\[
 |\E[\widehat g_b]|\le\E|\widehat g_b|
 \le 2B p_b(1-p_b).
\]
On a group with no draw of $b$, exactly
$\widehat g_b=-p_bG^{-1}\sum_i a_i$; it is zero if the applied advantages
sum to zero. By comparison, the verified-replay component for a fixed target
$w$ has logit velocity $\rho(w_b-p_b)$, which approaches $\rho w_b>0$
for a banked category as $p_b\to0$.
\end{lemma}

\begin{proof}
Use the triangle inequality and $|a_i|\le B$. Each marginal draw satisfies
$\E|\1\{Z_i=b\}-p_b|=2p_b(1-p_b)$, irrespective of the dependence of
$a_i$ on the group. On the absence event all indicator terms are zero.
The replay formula follows by differentiating $-\sum_b w_b\log p_b$.
\end{proof}

The bound compares the isolated logit-gradient components. It does not prove
inevitable collapse, rule out retention by a bounded estimator, or guarantee
that replay dominates the full update. An unsampled mode can still change
through softmax normalization or shared parameters. Exact entropy is outside
the uniformly bounded-score premise because its surprisal is unbounded near
zero; the lemma is not an impossibility result for entropy regularization.
Nor does the absence of a mode from a finite sample establish zero probability.

\paragraph{Exact entropy is a different comparator.}
For a known full correct support and uniform $u$, the elementary identities
\[
 H(q)=\log m-\mathrm{KL}(q\|u),\qquad
 R_u(p)=\log m-\log P+\mathrm{KL}(u\|q)
\]
show that exact conditional entropy and categorical replay both favor a
uniform conditional target. With strictly increasing correctness potential
$\Psi$, both $\Psi(P)+\beta H(q)$ and $\Psi(P)-\rho R_u(p)$ have the unique
global optimum $P=1,q=u$ in probability space on the simplex with $P>0$,
for positive coefficients. This is a boundary limit for finite-logit softmax
policies. The optimum follows from nonnegativity of KL and increasing
$\Psi$; it does not assert convergence of either neural update. The KL direction distinguishes their
loss level sets \citep{pereyra2017confidencepenalty,agarwal2021policygradient}:
replay loss diverges when a protected coordinate vanishes, whereas
conditional entropy can remain $\log(m-1)$ with one missing mode.

Exact full-output entropy is a third objective and also regularizes incorrect
mass. Exact tabular entropy-regularized policy gradient has its own positive
probability and convergence guarantees under specified assumptions
\citep{mei2020softmaxpg}; our collapse theorem contains no entropy term.
The sampled semantic scores in the experiments are not exact gradients of
full-output or conditional verified-key entropy. Neither their empirical
comparison nor the bounded-score lemma is an impossibility result for entropy.

\paragraph{Discovery and optimization remain separate requirements.}
Retention starts after a verified key enters active memory. Actual insertion
depends on generation, validation, capacity, and eligibility; a positive
sampling probability alone does not establish eventual admission. Once a
non-evicting capacity-16 bank fills, an unbanked key can remain ineligible.
Changing membership or exemplar weights changes the loss being bounded, so
fixed-bank descent cannot simply be carried across those changes. In
particular, the fixed-weight argument does not cover a target whose weights
continue to change with fresh observation counts.

For a fixed joint potential, a discrete bound
$F(\theta_{n+1})\le F(\theta_n)+\epsilon_n$ with
$\epsilon_n\ge0$ and $\sum_{n\ge n_0}\epsilon_n=E<\infty$ would extend
Lemma~\ref{lem:shared-exemplar-retention}, using
$D=F(\theta_{n_0})+J_{\max}+E$ for steps $n\ge n_0$.
This follows by telescoping and bounds retention, not convergence. The noisy
optimizer would need a separate argument establishing such control of the
joint loss's upward drift. A positive replay coefficient or small nominal
learning rate does not establish it. Adam moments, clipping, momentum, prompt
scheduling, and bank changes are not covered by the exact-flow proofs.
The experiments have not established the required joint-energy or
update-error assumptions. The mathematical contribution is the stated
categorical mechanism and conditional exemplar argument; empirical support
and longitudinal exemplar scores remain separate measurements.

% Shared scientific provenance; operational inventories remain in the repository.

\section{What Verified Replay Costs}
\label{app:replay-cost}

App.~\ref{app:theory-kl} argues that a memory over verified successes and a
reference-KL penalty are different tools. They also carry different prices, and
the price is the easier of the two claims to get wrong by asserting it. Two
costs are separated here because they scale differently, and only one of them
favors the bank asymptotically.

Everything below is read from the \texttt{train\_metrics.jsonl} of the replay
arms of the three paired cohorts, \MDcostRuns{} runs across five domains and
three model families. No arm is re-executed and no new training is run: these
are the same logs the retention results are read from. Runs are keyed by the
ledger's arm field rather than by run name, because a replay experiment's name
appears on both of its arms.

\paragraph{Resident state, which diverges.}
A reference KL needs a frozen copy of the policy resident for the whole run,
$\Theta(|\theta|)$ bytes, growing with the model. A bank needs no parameters at
all: it holds token identifiers for verified exemplars, one prompt and one
response per admitted key, $\Theta(\text{prompts}\times\text{keys}\times
\text{length})$ bytes, a quantity that does not move when the model grows.
Table~\ref{tab:replay-cost} reports both. The bank stays at
\MDcostBankKB{}~kB or below at every scale while the reference copy grows from
$0.99$ to $6.17$~GB, so the ratio rises from \MDcostStateMin{} to
\MDcostStateMax{} across a model-size range of about six --- the linear
divergence in $|\theta|$ that the accounting predicts.

\paragraph{Per-update compute, which does not diverge.}
This cost does not favor replay asymptotically and is not claimed to. A
reference KL adds one reference forward pass over the fresh group, about
$2|\theta|T_{\mathrm{fresh}}$ FLOPs; replay adds a forward and a backward over
the selected bank, about $6|\theta|T_{\mathrm{replay}}$. Replay is therefore
cheaper exactly when $3T_{\mathrm{replay}}<T_{\mathrm{fresh}}$, a
constant-factor question. Because a bank stores the policy's own rollouts on
the same prompts, exemplar and fresh response lengths coincide and the
condition collapses to a statement about how full the banks get,
\begin{equation}
  \frac{\text{replay FLOPs}}{\text{reference-KL FLOPs}}
  =\frac{3\,\E|\mathcal B|}{G},
  \label{eq:replay-flop-ratio}
\end{equation}
so replay is the cheaper addition when mean occupancy stays below $G/3$. The
measured ratios in Table~\ref{tab:replay-cost} reproduce
Equation~\eqref{eq:replay-flop-ratio} to three decimals in all
\MDcostCellsTotal{} cells.

Mean occupancy runs from \MDcostOccMin{} to \MDcostOccMax{} keys against a
capacity of 16, and the ratio from \MDcostRatioMin{} to \MDcostRatioMax{}.
It is below one in \MDcostCellsBelow{} of \MDcostCellsTotal{} cells and above
one in the remaining three, all of them PantryPlan, whose banks fill furthest.
The honest summary is that replay is the cheaper of the two additions on four
of five domains and the more expensive one on PantryPlan, and that neither
ordering is asymptotic: both terms are linear in $|\theta|$, so the compute
comparison is a property of bank occupancy in a domain and not of the method.

\paragraph{Scope of the accounting.}
These are analytic FLOP and byte counts driven by measured token counts, not
measured runtimes. Every arm we report trains at $\beta_{\mathrm{KL}}=0$
(Table~\ref{tab:run-contract}), so no reference-KL wall clock exists anywhere
in this work to compare against, and none is claimed. The counts assume full
fine-tuning, which is what we run; a policy trained as a low-rank adapter
recovers its reference by disabling the adapter and pays no resident-state cost
at all, which removes the first argument and leaves the second. Bank bytes are
token identifiers at four bytes each and exclude the checkpoint bookkeeping
that the atomic state carries alongside them. The measured occupancies are a
property of these domains at this capacity, and
Equation~\eqref{eq:replay-flop-ratio} says exactly how the comparison would
move at another one.

\begin{table}[t]
  \centering
  \caption{\textbf{What each regularizer would add, measured on the replay
  arms.} Per family, the first row is mean bank occupancy $\E|\mathcal B|$ by
  domain against a capacity of 16, then the bytes a resident reference copy
  would need, the mean bank footprint, and their ratio. The second row is the
  per-update FLOP ratio of Equation~\eqref{eq:replay-flop-ratio}: below one
  means replay adds less compute than a reference-KL forward would. Read from
  \MDcostRuns{} replay runs; five seeds per cell.}
  \label{tab:replay-cost}
  \scriptsize
  \setlength{\tabcolsep}{4.2pt}
  \begin{tabular}{@{}lrrrrrrrr@{}}
    \toprule
    \textbf{Family} & \textbf{Graph} & \textbf{Countd.} & \textbf{Python}
      & \textbf{MathIR} & \textbf{Pantry} & \textbf{Ref (GB)}
      & \textbf{Bank (kB)} & \textbf{Ratio} \\
    \midrule
    \input{results/replay_cost_table_body.tex}
  \end{tabular}
\end{table}

\clearpage
\phantomsection
\begin{center}\large\textbf{Part VII. Extended comparison}\end{center}
\addcontentsline{toc}{part}{Part VII. Extended comparison}

\section{Diversity-Preserving RLVR: Extended Comparison}
\label{app:related-extended}

That RLVR concentrates its outputs is not our discovery, and this appendix
states what three contemporary lines establish and exactly which of our claims
each one already covers.
\citet{cui2025entropy} identify entropy collapse in reasoning RL, attribute the
fall to a covariance between token probability and advantage, and fit a ceiling
on downstream reward as a function of remaining entropy; Clip-Cov and KL-Cov
restrict the high-covariance tokens and postpone the fall.
\citet{lochab2026ucpo} show that RLVR's objective is indifferent to how correct
probability is divided among correct answers, characterize the uniform-correct
policy as the optimum of a robustness and entropy-regularized criterion, and
add a conditional uniformity penalty to GRPO.
\citet{li2026distribution} attribute the same concentration to reverse-KL
mode seeking and replace it with an approximate forward-KL objective that
aligns the group to a target proportional to sampled rewards.
We take the phenomenon as established by this literature and do not claim it.
Our three separable claims are a success-conditional measurement, an
intervention that is not a fresh-update objective, and a stated reason the two
classes differ in reach.

\paragraph{The axis each work reads.}
\citet{cui2025entropy} track policy entropy over all sampled tokens. That
quantity mixes correct and incorrect mass, so it falls both when a policy stops
offering alternative correct answers and when it stops offering wrong ones, and
a training run that improves only by suppressing errors registers as a loss of
diversity. \citet{lochab2026ucpo} report \texttt{pass@}$K$ and an
equation-level count within the correct set; \citet{li2026distribution} read
diversity through downstream score on problems with many feasible solutions.
\texttt{pass@}$K$ and downstream score rise with the success rate on their
own: \eqref{eq:main-occupancy} makes the dependence explicit, and across the
\MDcells{} untrained measurements of App.~\ref{app:base-levels-all-scales}
\distk{} correlates $\MDdistinctcorr$ with \texttt{pass@8}.
\pmd{} removes that dependence by construction --- it is a function of $q$
alone --- so correctness and breadth are two axes a method can move in opposite
directions, and the concentration Sec.~\ref{sec:results-collapse} reports is
not a restatement of a change in the success rate. \mb{} also fixes sameness with the executable check that assigns
reward, so the equivalence used to define a mode is the one used to grant
credit, with no judge or embedding threshold between them.

\paragraph{Where the intervention can act.}
All three interventions modify the fresh on-policy update. Clip-Cov and KL-Cov
reweight tokens of the current rollouts; UCPO reallocates a prompt group's
positive advantage mass across that group's verifier-positive rows; DMPO forms
its target over the trajectories the group sampled. Each therefore requires a
mode to be present in the current group before it can act on that mode, and
that requirement is the pressure App.~\ref{app:theory} formalizes.
Lemma~\ref{lem:sampled-score-bound} applies to this whole class whenever the
per-row scores $a_i$ are detached and uniformly bounded by some $B$: the
expected logit velocity on category $b$ is at most $2Bp_b(1-p_b)$, which
vanishes with $p_b$, and on a group containing no draw of $b$ the component is
exactly $-p_bG^{-1}\sum_i a_i$. UCPO's estimator satisfies the premise as
implemented here. Write $G^{+}$ for the number of verifier-positive rows in the
group, $A^{+}$ for the per-row positive advantage a uniform split would assign,
$\hat q_i$ for row $i$'s estimated within-group mode frequency, and
$\omega\in[0,1]$ for the weight UCPO puts on inverse-frequency rather than
uniform reallocation. Such a row receives
$A_i=G^{+}A^{+}\bigl((1-\omega)/G^{+}+\omega v_i/\sum_j v_j\bigr)$ with
$v_i=1/\hat q_i$ self-normalized within the group, so
$|A_i|\le G^{+}|A^{+}|$ and the total positive mass is conserved --- it is a
redistribution among sampled correct rows, not an addition to an absent one. Verified replay is outside the premise because its
target is not a function of the current group: for a banked key the logit
velocity is $\rho(w_b-p_b)$, which approaches $\rho w_b>0$ as $p_b\to0$. This
is a statement about which mechanism retains a signal at the boundary, and
App.~\ref{app:theory-entropy} records what it is not: it does not prove that
collapse is inevitable, does not rule out retention by a bounded estimator, and
is not an impossibility result for entropy regularization, whose surprisal is
unbounded near zero and therefore fails the bounded-score premise.

\paragraph{The target is shared; the weighting is not.}
UCPO's uniform-correct target is the target our replay loss descends, and
App.~\ref{app:theory-entropy} shows that exact conditional entropy shares that
optimum too, the three differing in the direction of the KL and hence in what
happens to the loss when a mode's probability vanishes. So the disagreement
with \citet{lochab2026ucpo} is not about where to aim. DMPO is the one of the
three that aims elsewhere: a target proportional to rewards over sampled
trajectories is, under a binary verifier, uniform over the correct sampled
trajectories, which places mass on each mode in proportion to how many times
the group produced it. That is frequency weighting, and it is the control we
run directly. Holding the bank and the reuse fixed and weighting modes by
rediscovery frequency instead of equally costs $+.318$ extra verified modes per
prompt ($[+.272,+.358]$) with no aggregate correctness difference
(App.~\ref{app:frequency-replay-progress}). Reuse is not the active
ingredient; declining to let popularity set the rehearsal rate is.

\paragraph{What we measured rather than cited.}
Two of the three lines above are implemented directly and run as matched
comparators rather than only discussed; the third is covered by the
frequency-weighting control already described, and sparse RLEP-Dr is added
beside them for the memory side. UCPO is implemented from its Equations 16, 47
and 51, and runs
against seed-matched Dr.GRPO at Qwen2.5-0.5B and Falcon3-1B
(App.~\ref{app:ucpo-progress}); its gains are local, raising Countdown
correctness and \pmd{} there while its MathIR effects are inconclusive. An
entropy-shaped objective on the success-conditional axis is implemented as
Fixed Semantic MaxEnt (App.~\ref{app:semantic-maxent}), which supplies a
bounded surprisal bonus to underrepresented verified outcomes that appear in a
fresh rollout; its effects are likewise heterogeneous, and App.~\ref{app:semantic-maxent}
states why it carries neither an exact entropy-gradient interpretation nor an
individual-mode preservation guarantee. Sparse RLEP-Dr covers the memory side
\citep{zhang2025rlep}: the same reuse of verified successes, without uniform
canonical-key weighting.

\paragraph{Scope of the comparison.}
These comparators are read at the coverage they have. UCPO and sparse RLEP-Dr
run at two of our three scales on \mb{}, whose executable domains are not the
mathematical reasoning suites of \citet{lochab2026ucpo} or the NP-Bench
problems of \citet{li2026distribution}, and the token-level interventions of
\citet{cui2025entropy} act at a different granularity from the outcome-level
comparators we run. Nothing here ranks these
methods on their own benchmarks. The claim is narrower and is the one the
lemma and the weighting ablation support jointly: an objective confined to the
current group and a memory over verified keys reach different modes, and
uniform rehearsal over those keys is what separates our memory from replay that
revisits successes as often as it found them.

\clearpage
\phantomsection
\begin{center}\large\textbf{Part VIII. Records}\end{center}
\addcontentsline{toc}{part}{Part VIII. Records}

\section{Source Integrity and Reproducibility}
\label{app:reproducibility}

\subsection{Admissible sources and missing checkpoints}
\label{app:source-integrity}

Falcon3-1B Countdown Re:Dr seed 59 contains two complete four-draw
terminal evaluation sequences with identical request identities but
conflicting metrics. That source is excluded at every
checkpoint and from every efficacy endpoint, paired contrast, and aggregate
that reuses it; neither sequence is selected or averaged. The independently
valid Dr.GRPO and GRPO seed-59 sources remain admissible. The core replay
comparison therefore has 74 pairs: 14 complete five-seed blocks and Falcon
Countdown at $n=4$ (seeds 55--58). Primary partial intersections are
descriptive and receive no confidence interval.

A second kind of conflict is more common and has a diagnosable cause. The job
watchdog can requeue a run; the restart finds no resumable checkpoint, begins
from initialization, and evaluates the initial model a second time, appending a
second step-0 block to the same record. Separate attempts can land on different
GPU models, and GPU model changes floating-point reduction order and therefore
the sampled tokens, so the two blocks disagree and the checkpoint is invalidated
under the rule below. This is why several MaxRL initial checkpoints are
unavailable, leaving the longitudinal contrast with three defined seeds in
Falcon3-1B Graph and one in Falcon3-1B PantryPlan. The same hardware
sensitivity means that initial evaluations of the identical frozen base model
agree within a GPU pool and can differ across pools, which is one reason
initial evaluations are measured within each method rather than shared across
arms; a before--after contrast compares a method with its own initial
checkpoint and is unaffected. Of the step-0 blocks in a run's records, exactly one is the block its admitted
trajectory continues from, and App.~\ref{app:initial-provenance} admits that
block for the 27 cells where it is unique. That amendment was adopted after the
analysis was specified, so it reports every affected block under both
admissions.

Trajectory records use registered sources and approved continuations or
replacements, with superseded attempts excluded. Identical duplicates are
deduplicated; conflicting records invalidate the affected checkpoint rather
than being selected by metric value. Missing or ambiguous checkpoints remain gaps;
curves do not interpolate or carry values forward across them. An ambiguous
nonterminal checkpoint does not invalidate an independently admissible
terminal endpoint. Each curve declares its fixed source-admissible seed cohort; paired core
curves use the admitted intersection. Accuracy and mode-count views use
the same populations and gap rules. Supporting methods retain their
explicit method-specific cohorts; available checkpoint counts and gaps
are recorded alongside the curves. A curve point is drawn from the cohort seeds
that report at that checkpoint and is ringed where that is fewer than all of
them, so a checkpoint whose evidence exists for part of the cohort is read as
the partial measurement it is rather than as an absence; a checkpoint no seed
reports remains a gap. Machine-readable records identify the
sources, exclusions, exact seed sets, checkpoint rule, and source hashes.

\subsection{Conflicted initial checkpoints, admitted by provenance}
\label{app:initial-provenance}

A conflicted step-0 endpoint is a requeue duplicate, not a missing evaluation,
and the two attempts disagree only because they executed on different GPU
models. A fourth amendment to the concentration protocol therefore admits, for
each affected cell, the step-0 block that the cell's admitted trajectory
continues from --- chosen by file order, step number and the presence of
training rows, never by a metric value. The rule resolves 27 of the 36
conflicted cells and leaves nine, where more than one such block exists,
unavailable as before.

This amendment was adopted after the analysis was specified, so it reports
every affected block under both admissions (Table~\ref{tab:initial-provenance}) and is stated as a
post-specification change rather than a confirmation. Two checks bound what it
can have moved. The twelve Dr.GRPO and GRPO before--after comparisons that
carry Section~\ref{sec:results-collapse} are unaffected and reproduce to
$10^{-9}$, and no terminal replay-effect block changes at all, since those
compare two trained checkpoints and use no initial endpoint.

What the re-admission buys is precision, not a new direction: every affected
block moves from one to four defined seeds up to the registered five and
receives a nominal interval, while its sign is unchanged. It does not license a
broader collapse claim. Falcon3-1B Graph under MaxRL moves the other way, and
now does so with an interval excluding zero, as does its replay variant; MaxRL
concentrates on Countdown, MathIR and PantryPlan at that scale but not on
Graph. We therefore continue to state the training-induced concentration result
for the RLVR-only objectives it was measured on, and report the MaxRL
longitudinal cohort as the mixed result it is.

\begin{table}[H]
  \centering
  \caption{\textbf{Every block the provenance amendment touches, under both
  admissions.} Entries are the change in conditional collision from the initial
  to the terminal checkpoint on jointly eligible prompts, so positive is more
  concentrated. $n$ is the defined seed count out of five registered; intervals
  are nominal paired Student-$t$ and appear only for complete blocks. The
  preserved columns reproduce the amendment-3 analysis, which this amendment
  leaves on disk unchanged beside its own result.}
  \label{tab:initial-provenance}
  \small
  \setlength{\tabcolsep}{4pt}
  \begin{tabular}{@{}lll rl rl@{}}
    \toprule
    & & & \multicolumn{2}{c}{\textbf{Preserved}} & \multicolumn{2}{c}{\textbf{Provenance-admitted}} \\
    \cmidrule(lr){4-5}\cmidrule(l){6-7}
    \textbf{Model} & \textbf{Domain} & \textbf{Method} &
      $n$ & $\Delta C$ & $n$ & $\Delta C$ \\
    \midrule
    \input{results/initial_provenance_comparison_table_body.tex}
  \end{tabular}
\end{table}

\subsection{Conflicted registered checkpoints, admitted by provenance}
\label{app:registered-provenance}

The requeue that duplicates a step-0 block duplicates blocks at registered
training steps in the same way, and the frozen snapshot invalidates those
checkpoints under the same rule, which the training-curve figures draw as gaps.
A fifth amendment extends the step-0 rule of App.~\ref{app:initial-provenance}
to any registered step: within a source file whose attempt trained past that
step, a block of four draws counts as continued when the next evaluation row
carries a later step, and the checkpoint is admitted only where exactly one
continued block exists across the cell's files. At step 0 this is the fourth
amendment's rule verbatim, so that amendment's 27 initial admissions are
carried by the same mechanism; its own artifacts are left unchanged. The rule
reads file order, step numbers, draw indices and the presence of later training
rows, never a metric value.

Of the 70 conflicted registered checkpoints away from step 0, the rule resolves
27. The remaining 43 carry more than one continued block, because the run was
requeued and more than one attempt trained past that step, so nothing in the
source structure separates them. They stay unavailable, and the figures draw
them as gaps rather than interpolating across them. No re-evaluation can settle
them: only the terminal export of each run was retained, so the mid-training
weights those checkpoints evaluated no longer exist. The amendment is a partial
repair by construction. It was adopted after the figures it changes were first
drawn, and extends a rule itself adopted after the original analysis was
specified; both are recorded in
\texttt{paper/audits/training\_curve\_provenance\_20260919}.

\subsection{Data, models, and scientific records}
\label{app:repro-artifacts}
\label{app:repro-code}
\label{app:archived-models}

The \href{https://huggingface.co/datasets/od2961/ModeBench}{ModeBench dataset}
provides separate level/domain configurations, fixed split roles, full row
identities, and construction and admission records. Its
\href{https://huggingface.co/datasets/od2961/ModeBench\#loading}{loading guide}
and \href{https://huggingface.co/datasets/od2961/ModeBench\#evaluation}{evaluation guide}
identify the reported splits and preserve auxiliary data separately.
The \href{https://huggingface.co/od2961/maxent-grpo-models}{model archive}
links each available model--domain--method--seed export to an immutable
upload revision containing original weights, configuration, tokenizer,
file hashes, completion provenance, and the applicable base-model license.
The \href{https://huggingface.co/od2961/maxent-grpo-models/blob/main/EXPERIMENTS.md}{experiment index}
maps the core scale comparisons, MaxRL factorial, Level-2 factorial,
replay-weighting ablation, and supporting controls to their model families;
the \href{https://huggingface.co/od2961/maxent-grpo-models/blob/main/DATA.md}{scientific-record index}
and \href{https://huggingface.co/od2961/maxent-grpo-models/blob/main/RESTORE.md}{restore guide}
identify preserved analysis inputs and SHA-256 verification procedures.

Each empirical figure has a machine-readable record identified in the
figure manifest, containing its displayed values and frozen numerical inputs.
\path{paper/FIGURE_MANIFEST.md} maps both paper versions to their builders,
records, and reproduction commands. Detailed source and recovery inventories
remain in those records and the repository's archived audit documents.
Model availability may extend beyond the frozen analysis population; it
does not change a figure's source-admission or seed-intersection rule.
Terminal exports restore inference models, while intermediate trajectories
require the recorded evaluation checkpoints.

\subsection{Exact continuation and validation}
\label{app:repro-recovery}

Exact training continuation requires the atomic optimizer, bank, token,
cursor, and request-stream state in App.~\ref{app:algorithm-control},
in addition to model weights. State mismatch, non-finite values, missing
completion evidence, changed dataset identities, and source or aggregate
hash drift are hard failures. Placement-only recovery may change hardware
or wall time but must preserve the scientific update. Figure reproduction
uses frozen records; missing results are never replaced by shallower
checkpoints or silently incorporated from a newer archive.

\end{document}
